% Generated by scripts/render_blueprint.py; edit blueprint/manifest.json instead.
% Planned names are prose, never links to nonexistent Lean declarations.
\section{Scope and reading guide}
\paragraph{Main result: $\omega_F\le9/4$ over every field.}
\hyperref[thm:omega-le-nine-fourths]{Jump to the theorem and its Lean proof.}
The main theorem is kernel-checked.
Here $\omega_F$ is defined from the asymptotic rank of the square matrix multiplication
tensor, exactly as in the manuscript. The final theorem has no assumptions
about characteristic, cardinality, value functionals, or unproved tensor inequalities.
All mathematical entries in the audited blueprint are formalized.
The exponent-convention remark and computational interpretation are expository.
The library defines $\omega_F$ using asymptotic tensor rank and also identifies
it with the infimum of ordinary-rank logarithmic bounds. The introduction's
identification with an arithmetic-operation exponent is a conventional
interpretation, not a formalized computational-model equivalence.

Every active numbered definition, lemma, corollary, theorem, and remark in
\texttt{material/main.tex} has a corresponding node below. Original labels are
preserved; previously unlabelled environments have acquired stable labels.
Paper results retain their manuscript numbers, and labelled equations retain
the manuscript's separate equation numbering. Added foundations use auxiliary
labels A1, A2, and so on. The same labels appear in the dependency graph and its
popups; source locations identify the corresponding manuscript passages.

The canonical inventory is \texttt{blueprint/manifest.json}. Status
\texttt{planned} means that a proof route is specified, not that a proof has been
checked. Status \texttt{notready} means that more design work is needed.
Only checked declarations receive Lean links and formalization marks. Generic
lemmas and their tensor-theoretic applications have separate nodes.
In the dependency graph, boxes represent theorem-like results and equations; ellipses
represent definitions and remarks. Bold borders mark results whose proof or derivation is
supplied in the manuscript, including prior results reproved there. Thin borders
mark the other nodes, including cited inputs and added formalization helpers.
Border thickness is independent of the colors recording Lean progress.
Light green backgrounds mark completed formalized definitions, which have no
separate theorem proof to supply. Dark green backgrounds mark results whose
proofs and dependencies are formalized, including remarks with checked proofs.
The main result has a larger label and a double border, retaining its status color.

The main dependency chain passes from tensor and asymptotic-rank foundations
through separation and fractional value to the polynomial multiplication
degenerations, the diagonal recurrence, and elementary four-thirds growth.
The Gamma-function sharpening and value outside the admissible interval remain
part of full-manuscript coverage, but are not required for the main theorem.

\paragraph{Finite values without circularity.}
Bound individual witnesses first, using asymptotic rank at $\tau=\omega/3$ and
flattenings at $\tau=2/3$. Monotonicity of the exponent in each witness ratio
bounds its intermediate scores by an endpoint before taking the supremum.
Thus finiteness does not depend on value normalization or value
multiplicativity, and the value calculus has no circular justification.

\paragraph{Partial-sum values without entropy limits.}
The checked proof first sums the finite type-separation bounds using the
multinomial theorem. There are only polynomially many count vectors, so taking
roots removes the overhead and proves superadditivity of $V_\tau^{1/\tau}$.
Weighted AM--GM then gives the entropy bound for every real probability vector,
including zero coordinates. This reverses the order of the two manuscript
inequalities and removes multinomial entropy limits and rational approximation
from their dependencies. The correspondence notes record this alternative proof.

\paragraph{Polynomial multiplication and the completed exponent bound.}
The concavity degeneration is checked using four explicit coefficient maps
and three block identities, avoiding the general exact-sequence theorem.
The three-block degeneration keeps the shared first factor fixed and reverses
the middle block's other two coordinates. Their value inequalities supply
all numerical growth hypotheses. At $\tau=\omega_F/3$, elementary growth and
the uniform diagonal degeneration give
$a^{4\omega_F/9}\le V_{\omega_F/3}(\mathrm{PM}_{a,a})\le2a-1$ for all $a\ge1$.
The numerical comparison proves the main theorem without the exact polynomial
rank formula or the Gamma-function estimate. The revised manuscript now also
uses the uniform constant-one bound; the correspondence notes distinguish
its cited exact rank identity from the upper bound sufficient for Lean.

\section{Additional foundations}

\renewcommand{\thetheorem}{A1}
\begin{definition}[Finite tensors, restrictions, and equivalences]
\label{aux:tensor-calculus}
\lean{MatrixMultiplication.Tensor,MatrixMultiplication.Tensor.map,MatrixMultiplication.Tensor.Restricts,MatrixMultiplication.Tensor.Equivalent,MatrixMultiplication.Tensor.kronecker,MatrixMultiplication.Tensor.directSum,MatrixMultiplication.Tensor.permuteS3,MatrixMultiplication.Tensor.kroneckerPower,MatrixMultiplication.Tensor.reindex}
\leanok
\uses{aux:coefficient-tensors,aux:factor-permutations,aux:asymptotic-rank-definition,aux:ordinary-rank-definition}
A tensor over a field $F$ has three finite coordinate sets. A restriction is the image under three linear maps, and equivalence means mutual restriction. Define factor permutations, Kronecker products, powers, direct sums, and coordinate reindexings.
\end{definition}

\noindent\textit{Auxiliary foundation. Source:} \texttt{material/main.tex:299}.
\textit{Status:} \texttt{formalized}. Additional manuscript coverage.
\textit{Checked Lean declarations:} \leandecl{MatrixMultiplication.\allowbreak{}Tensor}, \leandecl{MatrixMultiplication.\allowbreak{}Tensor.\allowbreak{}map}, \leandecl{MatrixMultiplication.\allowbreak{}Tensor.\allowbreak{}Restricts}, \leandecl{MatrixMultiplication.\allowbreak{}Tensor.\allowbreak{}Equivalent}, \leandecl{MatrixMultiplication.\allowbreak{}Tensor.\allowbreak{}kronecker}, \leandecl{MatrixMultiplication.\allowbreak{}Tensor.\allowbreak{}directSum}, \leandecl{MatrixMultiplication.\allowbreak{}Tensor.\allowbreak{}permuteS3}, \leandecl{MatrixMultiplication.\allowbreak{}Tensor.\allowbreak{}kroneckerPower}, \leandecl{MatrixMultiplication.\allowbreak{}Tensor.\allowbreak{}reindex}.
\textit{Formalization note:} This definition node records exactly the coordinate definitions in its statement, all of which are implemented. Its former open-ended proof plan bundled additional structural identities; individual proved operation laws are tracked in the dedicated restriction, power, permutation, and sum nodes. No unlisted associativity or distributivity theorem is claimed here.

\begin{proof}
\proves{aux:tensor-calculus}
\leanok
\textit{Checked proof.} Use finite coefficient arrays, finite-sum factor-map actions, and mutual restriction for equivalence. Define permutations by rearranging coordinate types, powers by finite tuples, direct sums by disjoint coordinate blocks, and reindexing by coordinate bijections. Operation laws are recorded in their separate support nodes.
\end{proof}

\renewcommand{\thetheorem}{A2}
\begin{lemma}[Mutual restriction and invertible factor maps]
\label{aux:invertible-tensor-equivalence}
\lean{MatrixMultiplication.exists_linearEquiv_comp_eq_of_finrank_range_eq,MatrixMultiplication.matrix_exists_unit_mul_eq_of_rank_eq,MatrixMultiplication.Tensor.equivalent_of_unit_maps,MatrixMultiplication.Tensor.equivalent_iff_invertible_factor_maps}
\leanok
For tensors in the same ambient coordinate spaces, $T\cong S$ if and only if $S=(A\otimes B\otimes C)T$ for invertible matrices $A,B,C$.
\end{lemma}

\noindent\textit{Auxiliary foundation. Source:} \texttt{material/main.tex:299}.
\textit{Status:} \texttt{formalized}. Additional manuscript coverage.
\textit{Checked Lean declarations:} \leandecl{MatrixMultiplication.\allowbreak{}exists\_\allowbreak{}linearEquiv\_\allowbreak{}comp\_\allowbreak{}eq\_\allowbreak{}of\_\allowbreak{}finrank\_\allowbreak{}range\_\allowbreak{}eq}, \leandecl{MatrixMultiplication.\allowbreak{}matrix\_\allowbreak{}exists\_\allowbreak{}unit\_\allowbreak{}mul\_\allowbreak{}eq\_\allowbreak{}of\_\allowbreak{}rank\_\allowbreak{}eq}, \leandecl{MatrixMultiplication.\allowbreak{}Tensor.\allowbreak{}equivalent\_\allowbreak{}of\_\allowbreak{}unit\_\allowbreak{}maps}, \leandecl{MatrixMultiplication.\allowbreak{}Tensor.\allowbreak{}equivalent\_\allowbreak{}iff\_\allowbreak{}invertible\_\allowbreak{}factor\_\allowbreak{}maps}.
\textit{Formalization note:} The final correspondence audit identified this prose assertion separately from the definition of equivalence. The checked proof works over every field, including zero tensors and empty coordinate spaces.

\begin{proof}
\proves{aux:invertible-tensor-equivalence}
\leanok
\uses{aux:tensor-calculus,aux:flattening-calculus}
\textit{Checked proof.} Mutual restriction preserves all flattening ranks. Each factor map is therefore injective on the corresponding tensor support. Extend its induced isomorphism of support subspaces to an ambient automorphism and replace the three factor maps successively. Inverse factor maps prove the converse.
\end{proof}

\renewcommand{\thetheorem}{A3}
\begin{definition}[Algebraic degeneration over every field]
\label{aux:degeneration-calculus}
\lean{MatrixMultiplication.Tensor.Degenerates,MatrixMultiplication.Tensor.RationalDegenerates,MatrixMultiplication.Tensor.rationalDegenerates_iff_degenerates,MatrixMultiplication.Tensor.Degenerates.refl,MatrixMultiplication.Tensor.Restricts.degenerates,MatrixMultiplication.Tensor.Degenerates.trans,MatrixMultiplication.Tensor.Degenerates.kronecker,MatrixMultiplication.Tensor.Degenerates.directSum}
\leanok
Write $T\trianglerighteq S$ when polynomial linear maps send $T$ to $\varepsilon^dS+\varepsilon^{d+1}E(\varepsilon)$ for some nonnegative integer $d$. This agrees with rational-function specialization in the paper. Degeneration is reflexive and transitive, contains restriction, and respects tensor products and direct sums.
\end{definition}

\noindent\textit{Auxiliary foundation. Source:} \texttt{material/main.tex:312}.
\textit{Status:} \texttt{formalized}. Main proof path.
\textit{Checked Lean declarations:} \leandecl{MatrixMultiplication.\allowbreak{}Tensor.\allowbreak{}Degenerates}, \leandecl{MatrixMultiplication.\allowbreak{}Tensor.\allowbreak{}RationalDegenerates}, \leandecl{MatrixMultiplication.\allowbreak{}Tensor.\allowbreak{}rationalDegenerates\_\allowbreak{}iff\_\allowbreak{}degenerates}, \leandecl{MatrixMultiplication.\allowbreak{}Tensor.\allowbreak{}Degenerates.\allowbreak{}refl}, \leandecl{MatrixMultiplication.\allowbreak{}Tensor.\allowbreak{}Restricts.\allowbreak{}degenerates}, \leandecl{MatrixMultiplication.\allowbreak{}Tensor.\allowbreak{}Degenerates.\allowbreak{}trans}, \leandecl{MatrixMultiplication.\allowbreak{}Tensor.\allowbreak{}Degenerates.\allowbreak{}kronecker}, \leandecl{MatrixMultiplication.\allowbreak{}Tensor.\allowbreak{}Degenerates.\allowbreak{}directSum}.
\textit{Formalization note:} The rational-function relation, its equivalence with polynomial witnesses, and all listed operation laws are checked over every field.

\begin{proof}
\proves{aux:degeneration-calculus}
\leanok
\uses{aux:polynomial-degeneration-definition,aux:polynomial-degeneration-operations,aux:rational-degeneration-equivalence}
\textit{Checked proof.} Clear denominators and normalize the leading coefficient to compare the two definitions, as proved in the rational-degeneration support node. For composition substitute a sufficiently high power of the parameter before composing the polynomial witnesses. The parameter is formal, so no infinitude or characteristic assumption is introduced.
\end{proof}

\renewcommand{\thetheorem}{A4}
\begin{lemma}[Asymptotic rank calculus]
\label{aux:asymptotic-rank-calculus}
\lean{MatrixMultiplication.Tensor.asymptoticRank_eq_limit,MatrixMultiplication.Tensor.asymptoticRank_pow,MatrixMultiplication.Tensor.asymptoticRank_kronecker_le,MatrixMultiplication.Tensor.asymptoticRank_directSum_le,MatrixMultiplication.Tensor.asymptoticRank_mono_degeneration,MatrixMultiplication.Tensor.rank_mono,MatrixMultiplication.Tensor.rank_kronecker_le,MatrixMultiplication.Tensor.asymptoticRank_le_borderRank,MatrixMultiplication.Tensor.borderRank_le_rank}
\leanok
Rank is monotone under restriction and submultiplicative. Its normalized powers converge to asymptotic rank. Asymptotic rank is power multiplicative, submultiplicative, subadditive, and monotone under degeneration; in particular $\widetilde R(T)\le\underline R(T)\le R(T)$.
\end{lemma}

\noindent\textit{Auxiliary foundation. Source:} \texttt{material/main.tex:354}.
\textit{Status:} \texttt{formalized}. Additional manuscript coverage.
\textit{Checked Lean declarations:} \leandecl{MatrixMultiplication.\allowbreak{}Tensor.\allowbreak{}asymptoticRank\_\allowbreak{}eq\_\allowbreak{}limit}, \leandecl{MatrixMultiplication.\allowbreak{}Tensor.\allowbreak{}asymptoticRank\_\allowbreak{}pow}, \leandecl{MatrixMultiplication.\allowbreak{}Tensor.\allowbreak{}asymptoticRank\_\allowbreak{}kronecker\_\allowbreak{}le}, \leandecl{MatrixMultiplication.\allowbreak{}Tensor.\allowbreak{}asymptoticRank\_\allowbreak{}directSum\_\allowbreak{}le}, \leandecl{MatrixMultiplication.\allowbreak{}Tensor.\allowbreak{}asymptoticRank\_\allowbreak{}mono\_\allowbreak{}degeneration}, \leandecl{MatrixMultiplication.\allowbreak{}Tensor.\allowbreak{}rank\_\allowbreak{}mono}, \leandecl{MatrixMultiplication.\allowbreak{}Tensor.\allowbreak{}rank\_\allowbreak{}kronecker\_\allowbreak{}le}, \leandecl{MatrixMultiplication.\allowbreak{}Tensor.\allowbreak{}asymptoticRank\_\allowbreak{}le\_\allowbreak{}borderRank}, \leandecl{MatrixMultiplication.\allowbreak{}Tensor.\allowbreak{}borderRank\_\allowbreak{}le\_\allowbreak{}rank}.
\textit{Formalization note:} All stated laws and the comparison with the manuscript rational-function degeneration relation are checked over every field.

\begin{proof}
\proves{aux:asymptotic-rank-calculus}
\leanok
\uses{aux:ordinary-rank-definition,aux:asymptotic-rank-definition,aux:rank-operation-bounds,aux:degeneration-calculus,aux:asymptotic-rank-limit,aux:asymptotic-rank-powers-products,aux:asymptotic-rank-sums,aux:asymptotic-rank-degeneration,aux:border-rank-bound}
\textit{Checked proof.} Use submultiplicativity and the subadditive-sequence limit theorem. To handle a polynomial degeneration witness, extract coefficient $dN$ after taking $N$ powers. There are only polynomially many triples of coefficient maps, giving a polynomial overhead times $R(T^{\otimes N})$. Its $N$-th root tends to one over every field. Derive subadditivity by the binomial expansion and uniform exponential rank bounds.
\end{proof}

\renewcommand{\thetheorem}{A5}
\begin{lemma}[Flattenings and their ranks]
\label{aux:flattening-calculus}
\lean{MatrixMultiplication.Tensor.flatteningRank,MatrixMultiplication.Tensor.flatteningRank_kronecker,MatrixMultiplication.Tensor.flatteningRank_mono_degeneration,MatrixMultiplication.Tensor.flatteningRank_mono,MatrixMultiplication.Tensor.flatteningRank_le_rank}
\leanok
Each tensor determines three flattening matrices. Their ranks are monotone under restriction and degeneration, multiplicative under Kronecker product, and bounded above by tensor rank.
\end{lemma}

\noindent\textit{Auxiliary foundation. Source:} \texttt{material/main.tex:393}.
\textit{Status:} \texttt{formalized}. Additional manuscript coverage.
\textit{Checked Lean declarations:} \leandecl{MatrixMultiplication.\allowbreak{}Tensor.\allowbreak{}flatteningRank}, \leandecl{MatrixMultiplication.\allowbreak{}Tensor.\allowbreak{}flatteningRank\_\allowbreak{}kronecker}, \leandecl{MatrixMultiplication.\allowbreak{}Tensor.\allowbreak{}flatteningRank\_\allowbreak{}mono\_\allowbreak{}degeneration}, \leandecl{MatrixMultiplication.\allowbreak{}Tensor.\allowbreak{}flatteningRank\_\allowbreak{}mono}, \leandecl{MatrixMultiplication.\allowbreak{}Tensor.\allowbreak{}flatteningRank\_\allowbreak{}le\_\allowbreak{}rank}.
\textit{Formalization note:} All stated laws are checked. The rational/polynomial degeneration equivalence transfers monotonicity to the manuscript relation.

\begin{proof}
\proves{aux:flattening-calculus}
\leanok
\uses{aux:flattening-operations,aux:flattening-powers,aux:flattening-degeneration,aux:degeneration-calculus}
\textit{Checked proof.} Use the checked flattening product and power identities. Group each flattening into a matrix tensor and transport the polynomial degeneration witness by Kronecker products of the complementary factor maps. Its asymptotic rank equals matrix rank, so asymptotic-rank monotonicity gives degeneration monotonicity.
\end{proof}

\renewcommand{\thetheorem}{A6}
\begin{lemma}[Matrix multiplication as a bilinear structure tensor]
\label{aux:matrix-bilinear-interpretation}
\lean{MatrixMultiplication.Tensor.matrixMul_bilinear,MatrixMultiplication.Tensor.structureTensor_matrixMul}
\leanok
\uses{lem:degeneration-sequences}
Contracting the first two factors of $\langle\ell,m,n\rangle$ with matrices $X,Y$ gives the entries of $XY$. In the standard matrix bases, this coefficient tensor equals the structure tensor of the bilinear multiplication map.
\end{lemma}

\noindent\textit{Auxiliary foundation. Source:} \texttt{material/main.tex:404}.
\textit{Status:} \texttt{formalized}. Additional manuscript coverage.
\textit{Checked Lean declarations:} \leandecl{MatrixMultiplication.\allowbreak{}Tensor.\allowbreak{}matrixMul\_\allowbreak{}bilinear}, \leandecl{MatrixMultiplication.\allowbreak{}Tensor.\allowbreak{}structureTensor\_\allowbreak{}matrixMul}.
\textit{Formalization note:} The final audit identified this interpretation as a separate background assertion. The contraction identity holds over semirings, and the literal structure-tensor equality is checked over every field. Zero dimensions are included.

\begin{proof}
\proves{aux:matrix-bilinear-interpretation}
\leanok
\uses{aux:explicit-tensor-definitions,aux:coefficient-tensors}
\textit{Checked proof.} Expand matrix multiplication as a finite sum and evaluate the equality indicators in the coefficient tensor. Compute the product of standard basis matrices to identify all structure-tensor coefficients.
\end{proof}

\renewcommand{\thetheorem}{A7}
\begin{lemma}[Basic bounds for the exponent]
\label{aux:omega-bounds}
\lean{MatrixMultiplication.two_le_omega,MatrixMultiplication.omega_le_three}
\leanok
\uses{aux:matrix-multiplication-exponent}
For every field $F$, $2\le\omega_F\le3$.
\end{lemma}

\noindent\textit{Auxiliary foundation. Source:} \texttt{material/main.tex:404}.
\textit{Status:} \texttt{formalized}. Main proof path.
\textit{Checked Lean declarations:} \leandecl{MatrixMultiplication.\allowbreak{}two\_\allowbreak{}le\_\allowbreak{}omega}, \leandecl{MatrixMultiplication.\allowbreak{}omega\_\allowbreak{}le\_\allowbreak{}three}.
\textit{Formalization note:} Both bounds are checked over every field. The lower bound comes from the first flattening; the upper bound comes from the elementary eight-term decomposition at matrix size two.

\begin{proof}
\proves{aux:omega-bounds}
\leanok
\uses{aux:omega-lower-bound,aux:matrix-multiplication-calculus}
\textit{Checked proof.} The square matrix multiplication tensor has flattening rank $n^2$ and the elementary decomposition has $n^3$ terms. Apply these bounds at $n=2$ and use monotonicity of logarithms.
\end{proof}

\renewcommand{\thetheorem}{A8}
\begin{definition}[Finite probability distributions and entropy]
\label{aux:entropy}
\lean{MatrixMultiplication.entropy}
\leanok
For a probability distribution $q$ on a nonempty finite set, define $H(q)=-\sum_iq_i\log_2q_i$, with $0\log0=0$.
\end{definition}

\noindent\textit{Auxiliary foundation. Source:} \texttt{material/main.tex:522}.
\textit{Status:} \texttt{formalized}. Additional manuscript coverage.
\textit{Checked Lean declarations:} \leandecl{MatrixMultiplication.\allowbreak{}entropy}.
\textit{Formalization note:} The definition extends to finite real weight vectors; probability nonnegativity and unit-sum conditions are hypotheses of the inequalities. Lean uses Real.log 0 = 0, giving the required zero-weight convention.

\renewcommand{\thetheorem}{A9}
\begin{lemma}[Multinomial entropy limit and approximation with fixed support]
\label{aux:multinomial-asymptotics}
\lean{MatrixMultiplication.multinomial_log_div_tendsto,MatrixMultiplication.log_two_mul_entropy_natCounts,MatrixMultiplication.multinomial_root_tendsto_entropy,MatrixMultiplication.probability_count_approximation,MatrixMultiplication.rational_probability_approximation,MatrixMultiplication.entropy_tendsto,MatrixMultiplication.entropy_weightedProduct_tendsto,MatrixMultiplication.entropy_weightedProduct_le_of_rational,MatrixMultiplication.entropy_weightedProduct_le_of_counts}
\leanok
If nonnegative integers $K_i$ sum to $K>0$ and $q_i=K_i/K$, then $\binom{nK}{nK_1,\ldots,nK_m}^{1/(nK)}\longrightarrow2^{H(q)}$. Every probability distribution is a limit of rational probability distributions with exactly the same support. Entropy-weighted product inequalities for rational probabilities therefore extend to all probabilities, including zero block values.
\end{lemma}

\noindent\textit{Auxiliary foundation. Source:} \texttt{material/main.tex:837}.
\textit{Status:} \texttt{formalized}. Additional manuscript coverage.
\textit{Checked Lean declarations:} \leandecl{MatrixMultiplication.\allowbreak{}multinomial\_\allowbreak{}log\_\allowbreak{}div\_\allowbreak{}tendsto}, \leandecl{MatrixMultiplication.\allowbreak{}log\_\allowbreak{}two\_\allowbreak{}mul\_\allowbreak{}entropy\_\allowbreak{}natCounts}, \leandecl{MatrixMultiplication.\allowbreak{}multinomial\_\allowbreak{}root\_\allowbreak{}tendsto\_\allowbreak{}entropy}, \leandecl{MatrixMultiplication.\allowbreak{}probability\_\allowbreak{}count\_\allowbreak{}approximation}, \leandecl{MatrixMultiplication.\allowbreak{}rational\_\allowbreak{}probability\_\allowbreak{}approximation}, \leandecl{MatrixMultiplication.\allowbreak{}entropy\_\allowbreak{}tendsto}, \leandecl{MatrixMultiplication.\allowbreak{}entropy\_\allowbreak{}weightedProduct\_\allowbreak{}tendsto}, \leandecl{MatrixMultiplication.\allowbreak{}entropy\_\allowbreak{}weightedProduct\_\allowbreak{}le\_\allowbreak{}of\_\allowbreak{}rational}, \leandecl{MatrixMultiplication.\allowbreak{}entropy\_\allowbreak{}weightedProduct\_\allowbreak{}le\_\allowbreak{}of\_\allowbreak{}counts}.
\textit{Formalization note:} The original entropy proof now has checked multinomial asymptotics and rational approximation with exactly unchanged support. Zero counts are included in the limit. Rounded-up natural counts give positive denominators at every approximation stage. The entropy-weighted product passes to the limit even when fixed block values vanish. These results are independent of the alternative multinomial-sum proof used on the main theorem path.

\begin{proof}
\proves{aux:multinomial-asymptotics}
\leanok
\uses{aux:entropy}
\textit{Checked proof.} Stirling's formula gives $\log((nc)!)/n-c\log n\to c\log c-c$, treating $c=0$ separately. Sum over the counts and cancel the $K\log n$ terms, then divide by $K$ and exponentiate. For approximation, take $c_r(i)=\lceil(r+1)q_i\rceil$ and normalize by $\sum_i c_r(i)$. Zero coordinates stay zero and the normalized counts converge to $q$. Continuity of $x\log x$ gives entropy convergence; unchanged zero coordinates handle the otherwise discontinuous $0^0$ corner in the weighted products.
\end{proof}

\renewcommand{\thetheorem}{A10}
\begin{lemma}[Sharp matrix-multiplication border-subrank lower bound]
\label{aux:matrix-mul-border-subrank}
\lean{MatrixMultiplication.Tensor.matrixMul_degenerates_diagonal_of_sumBand,MatrixMultiplication.centralBand,MatrixMultiplication.card_centralBand_eq_sub,MatrixMultiplication.card_centralBand,MatrixMultiplication.sharpSubrankSize_eq_natCeil,MatrixMultiplication.Tensor.matrixMul_degenerates_diagonal_sharp_nat,MatrixMultiplication.Tensor.matrixMul_degenerates_diagonal_sharp,MatrixMultiplication.Tensor.matrixMul_borderSubrank_lower_bound}
\leanok
For every positive integer $n$, $\underline Q(\langle n,n,n\rangle)\ge\lceil3n^2/4\rceil$.
\end{lemma}

\noindent\textit{Auxiliary foundation. Source:} \texttt{material/main.tex:697}.
\textit{Status:} \texttt{formalized}. Additional manuscript coverage.
\textit{Checked Lean declarations:} \leandecl{MatrixMultiplication.\allowbreak{}Tensor.\allowbreak{}matrixMul\_\allowbreak{}degenerates\_\allowbreak{}diagonal\_\allowbreak{}of\_\allowbreak{}sumBand}, \leandecl{MatrixMultiplication.\allowbreak{}centralBand}, \leandecl{MatrixMultiplication.\allowbreak{}card\_\allowbreak{}centralBand\_\allowbreak{}eq\_\allowbreak{}sub}, \leandecl{MatrixMultiplication.\allowbreak{}card\_\allowbreak{}centralBand}, \leandecl{MatrixMultiplication.\allowbreak{}sharpSubrankSize\_\allowbreak{}eq\_\allowbreak{}natCeil}, \leandecl{MatrixMultiplication.\allowbreak{}Tensor.\allowbreak{}matrixMul\_\allowbreak{}degenerates\_\allowbreak{}diagonal\_\allowbreak{}sharp\_\allowbreak{}nat}, \leandecl{MatrixMultiplication.\allowbreak{}Tensor.\allowbreak{}matrixMul\_\allowbreak{}degenerates\_\allowbreak{}diagonal\_\allowbreak{}sharp}, \leandecl{MatrixMultiplication.\allowbreak{}Tensor.\allowbreak{}matrixMul\_\allowbreak{}borderSubrank\_\allowbreak{}lower\_\allowbreak{}bound}.
\textit{Formalization note:} The sharp cited lower bound is checked over every field, including dimension zero. The explicit monomial degeneration works over every commutative semiring. The central-band cardinality and its equality to the ceiling formula are proved exactly. Lemma 4.5 continues to use the simpler quadratic construction. The sharp construction is additional checked support for the manuscript’s cited bound.

\begin{proof}
\proves{aux:matrix-mul-border-subrank}
\leanok
\uses{def:matrix-multiplication,def:rank-invariants,aux:monomial-degeneration}
\textit{Checked proof.} Set $m=\lfloor n/2\rfloor$ and retain triples $(i,j,i+j-m)$ for $0\le i,j<n$ and $m\le i+j<m+n$. The two removed triangular corners have total size $\lfloor n/2\rfloor\lceil n/2\rceil$, leaving $\lceil3n^2/4\rceil$ labels. Use weights $(i+j)^2$, $n^2+u^2-j^2$, and $4n^2-2i(i+z)$. On the tensor support $u+j=i+z$, their total is $5n^2+(i-u)^2$, whose minimum occurs exactly at diagonal labels. Reindex those labels by their cardinality and apply the border-subrank definition.
\end{proof}

\renewcommand{\thetheorem}{A11}
\begin{lemma}[Witness bounds and finiteness without circularity]
\label{aux:value-finiteness}
\lean{MatrixMultiplication.Tensor.ValueWitness.bound_omega,MatrixMultiplication.Tensor.ValueWitness.bound_two,MatrixMultiplication.Tensor.valueWitness_le_asymptoticRank,MatrixMultiplication.Tensor.fractionalValue_le_asymptoticRank,MatrixMultiplication.Tensor.fractionalValue_ne_top,MatrixMultiplication.Tensor.one_le_fractionalValue,MatrixMultiplication.Tensor.fractionalValue_unit}
\leanok
\uses{aux:fractional-value-definition}
For $2/3\le\tau\le\omega/3$, every defining witness score is bounded by $\widetilde R(T)$. Consequently $V_\tau(T)\le\widetilde R(T)<+\infty$. For nonzero $T$ the witness set is nonempty and $V_\tau(T)\ge1$; the unit diagonal has value exactly one in the admissible interval.
\end{lemma}

\noindent\textit{Auxiliary foundation. Source:} \texttt{material/main.tex:791}.
\textit{Status:} \texttt{formalized}. Main proof path.
\textit{Checked Lean declarations:} \leandecl{MatrixMultiplication.\allowbreak{}Tensor.\allowbreak{}ValueWitness.\allowbreak{}bound\_\allowbreak{}omega}, \leandecl{MatrixMultiplication.\allowbreak{}Tensor.\allowbreak{}ValueWitness.\allowbreak{}bound\_\allowbreak{}two}, \leandecl{MatrixMultiplication.\allowbreak{}Tensor.\allowbreak{}valueWitness\_\allowbreak{}le\_\allowbreak{}asymptoticRank}, \leandecl{MatrixMultiplication.\allowbreak{}Tensor.\allowbreak{}fractionalValue\_\allowbreak{}le\_\allowbreak{}asymptoticRank}, \leandecl{MatrixMultiplication.\allowbreak{}Tensor.\allowbreak{}fractionalValue\_\allowbreak{}ne\_\allowbreak{}top}, \leandecl{MatrixMultiplication.\allowbreak{}Tensor.\allowbreak{}one\_\allowbreak{}le\_\allowbreak{}fractionalValue}, \leandecl{MatrixMultiplication.\allowbreak{}Tensor.\allowbreak{}fractionalValue\_\allowbreak{}unit}.
\textit{Formalization note:} Checked over every field, including T=0 and omega=2. No value identity or prior finiteness assumption is used. Equality with the manuscript rational-function value is now proved separately.

\begin{proof}
\proves{aux:value-finiteness}
\leanok
\uses{cor:asymptotic-rank-lower-bound,aux:symmetrization-rank-bound,aux:flattening-degeneration,aux:flattening-repeated-blocks,lem:flattening-lb,aux:witness-endpoint-interpolation,aux:unit-value-witness}
\textit{Checked proof.} At the upper endpoint apply the equal-block asymptotic sum identity, degeneration monotonicity, and $\widetilde R(T^{\otimes\mathfrak S_3})\le\widetilde R(T)^6$ to obtain $ab^\omega\le cd^\omega\widetilde R(T)^{6N}$. At the lower endpoint use flattening-rank monotonicity, the repeated-block formula, and $R_1(T^{\otimes\mathfrak S_3})\le\widetilde R(T)^6$ to obtain the quadratic analogue. Bound intermediate witness scores by endpoint monotonicity of $(b/d)^{3\tau}$, then take the extended-real supremum. The all-one witness gives the lower bound; diagonal normalization gives unit value.
\end{proof}

\renewcommand{\thetheorem}{A12}
\begin{lemma}[Polynomial value satisfies the numerical hypotheses]
\label{aux:polynomial-value-hypotheses}
\lean{MatrixMultiplication.Tensor.polynomialMulValue,MatrixMultiplication.Tensor.polynomialMulValue_nonneg,MatrixMultiplication.Tensor.polynomialMulValue_symmetry,MatrixMultiplication.Tensor.polynomialMulValue_boundary,MatrixMultiplication.Tensor.polynomialMul_value_growthHypotheses}
\leanok
For admissible $\tau$, the function $P(a,b)=V_\tau(\mathrm{PM}_{a,b})^{1/\tau}$ is symmetric, satisfies $P(1,a)=a$, and satisfies the concavity and tripling inequalities. These properties are packaged as the numerical growth hypotheses.
\end{lemma}

\noindent\textit{Auxiliary foundation. Source:} \texttt{material/main.tex:1126}.
\textit{Status:} \texttt{formalized}. Main proof path.
\textit{Checked Lean declarations:} \leandecl{MatrixMultiplication.\allowbreak{}Tensor.\allowbreak{}polynomialMulValue}, \leandecl{MatrixMultiplication.\allowbreak{}Tensor.\allowbreak{}polynomialMulValue\_\allowbreak{}nonneg}, \leandecl{MatrixMultiplication.\allowbreak{}Tensor.\allowbreak{}polynomialMulValue\_\allowbreak{}symmetry}, \leandecl{MatrixMultiplication.\allowbreak{}Tensor.\allowbreak{}polynomialMulValue\_\allowbreak{}boundary}, \leandecl{MatrixMultiplication.\allowbreak{}Tensor.\allowbreak{}polynomialMul\_\allowbreak{}value\_\allowbreak{}growthHypotheses}.
\textit{Formalization note:} All numerical hypotheses are now supplied by the defined polynomial value over every field. The only parameter assumptions are admissibility; no tensor inequalities are assumed.

\begin{proof}
\proves{aux:polynomial-value-hypotheses}
\leanok
\uses{aux:numerical-growth-hypotheses,cor:value-diag-matrix,aux:polynomial-coordinate-models,aux:binary-partial-value-support,cor:pm-concavity,cor:pm-tripling}
\textit{Checked proof.} Identify $\mathrm{PM}_{1,a}$ with $\langle1,1,a\rangle$ and interchange the first two factors for symmetry. Combine normalization with the two polynomial-value inequalities.
\end{proof}

\renewcommand{\thetheorem}{A13}
\begin{lemma}[Elementary four-thirds growth]
\label{aux:polynomial-value-elementary-growth}
\manuscriptproof
\lean{MatrixMultiplication.Tensor.polynomialMul_value_elementary_lower_bound}
\leanok
For the manuscript function $P(a,b)=V_\tau(\mathrm{PM}_{a,b})^{1/\tau}$ with admissible $\tau$ and every $a\ge1$, $P(a,a)\ge a^{4/3}$.
\end{lemma}

\noindent\textit{Auxiliary foundation. Source:} \texttt{material/main.tex:1168}.
\textit{Status:} \texttt{formalized}. Main proof path.
\textit{Proof supplied in manuscript.}
\textit{Checked Lean declarations:} \leandecl{MatrixMultiplication.\allowbreak{}Tensor.\allowbreak{}polynomialMul\_\allowbreak{}value\_\allowbreak{}elementary\_\allowbreak{}lower\_\allowbreak{}bound}.
\textit{Formalization note:} The concrete polynomial-value bound is proved for every positive integer a, without Gamma-function estimates.

\begin{proof}
\proves{aux:polynomial-value-elementary-growth}
\leanok
\uses{aux:polynomial-value-hypotheses,aux:numerical-elementary-growth}
\textit{Checked proof.} Start from $P(1,1)=1$. The recurrence multiplier is at least $(1+1/(3a))(1+1/a)\ge(1+1/a)^{4/3}$. Alternatively cube the desired multiplier comparison and prove a polynomial inequality; induction gives $a^4\le P(a,a)^3$, and monotonicity converts it to the real-power statement.
\end{proof}

\renewcommand{\thetheorem}{A14}
\begin{lemma}[Comparison of real powers]
\label{aux:power-growth-comparison}
\lean{MatrixMultiplication.exponent_le_one_of_nat_rpow_le_two_mul}
\leanok
If $r\in\mathbb R$ and $n^r\le2n$ for every positive integer $n$, then $r\le1$.
\end{lemma}

\noindent\textit{Auxiliary foundation. Source:} \texttt{material/main.tex:1182}.
\textit{Status:} \texttt{formalized}. Main proof path.
\textit{Checked Lean declarations:} \leandecl{MatrixMultiplication.\allowbreak{}exponent\_\allowbreak{}le\_\allowbreak{}one\_\allowbreak{}of\_\allowbreak{}nat\_\allowbreak{}rpow\_\allowbreak{}le\_\allowbreak{}two\_\allowbreak{}mul}.
\textit{Formalization note:} The all-positive-integer formulation supplies the numerical comparison used by the elementary main proof.

\begin{proof}
\proves{aux:power-growth-comparison}
\leanok
\textit{Checked proof.} If $r>1$, choose an integer $n>2^{1/(r-1)}$. Monotonicity of real powers implies $n^{r-1}>2$, contradicting $n^r\le2n$. This checked argument avoids limits.
\end{proof}

\renewcommand{\thetheorem}{A15}
\begin{lemma}[Implemented coefficient tensor API]
\label{aux:coefficient-tensors}
\lean{MatrixMultiplication.Tensor,MatrixMultiplication.Tensor.map,MatrixMultiplication.Tensor.Restricts,MatrixMultiplication.Tensor.Restricts.refl,MatrixMultiplication.Tensor.Restricts.trans,MatrixMultiplication.Tensor.Equivalent,MatrixMultiplication.Tensor.Equivalent.refl,MatrixMultiplication.Tensor.Equivalent.symm,MatrixMultiplication.Tensor.Equivalent.trans,MatrixMultiplication.Tensor.kronecker,MatrixMultiplication.Tensor.directSum,MatrixMultiplication.Tensor.partialSum,MatrixMultiplication.Tensor.map_one,MatrixMultiplication.Tensor.map_map}
\leanok
Coefficient tensors, factor-map actions, restrictions, mutual-restriction equivalence, Kronecker products, binary direct sums, and binary partial sums are defined. Factor maps compose; restriction and equivalence have their expected reflexivity and transitivity laws.
\end{lemma}

\noindent\textit{Auxiliary foundation. Source:} \texttt{material/main.tex:299}.
\textit{Status:} \texttt{formalized}. Main proof path.
\textit{Checked Lean declarations:} \leandecl{MatrixMultiplication.\allowbreak{}Tensor}, \leandecl{MatrixMultiplication.\allowbreak{}Tensor.\allowbreak{}map}, \leandecl{MatrixMultiplication.\allowbreak{}Tensor.\allowbreak{}Restricts}, \leandecl{MatrixMultiplication.\allowbreak{}Tensor.\allowbreak{}Restricts.\allowbreak{}refl}, \leandecl{MatrixMultiplication.\allowbreak{}Tensor.\allowbreak{}Restricts.\allowbreak{}trans}, \leandecl{MatrixMultiplication.\allowbreak{}Tensor.\allowbreak{}Equivalent}, \leandecl{MatrixMultiplication.\allowbreak{}Tensor.\allowbreak{}Equivalent.\allowbreak{}refl}, \leandecl{MatrixMultiplication.\allowbreak{}Tensor.\allowbreak{}Equivalent.\allowbreak{}symm}, \leandecl{MatrixMultiplication.\allowbreak{}Tensor.\allowbreak{}Equivalent.\allowbreak{}trans}, \leandecl{MatrixMultiplication.\allowbreak{}Tensor.\allowbreak{}kronecker}, \leandecl{MatrixMultiplication.\allowbreak{}Tensor.\allowbreak{}directSum}, \leandecl{MatrixMultiplication.\allowbreak{}Tensor.\allowbreak{}partialSum}, \leandecl{MatrixMultiplication.\allowbreak{}Tensor.\allowbreak{}map\_\allowbreak{}one}, \leandecl{MatrixMultiplication.\allowbreak{}Tensor.\allowbreak{}map\_\allowbreak{}map}.
\textit{Formalization note:} Checked coordinate arrays and factor-map actions. Finite families, powers, and permutation/reindexing laws are checked in their dedicated support nodes.

\begin{proof}
\proves{aux:coefficient-tensors}
\leanok
\textit{Checked proof.} The checked implementation uses finite sums for the factor-map action. Reorder the six sums to prove map composition; use identity matrices for reflexivity and composed witnesses for transitivity.
\end{proof}

\renewcommand{\thetheorem}{A16}
\begin{definition}[Implemented diagonal and multiplication tensors]
\label{aux:explicit-tensor-definitions}
\lean{MatrixMultiplication.Tensor.diagonal,MatrixMultiplication.Tensor.matrixMul,MatrixMultiplication.Tensor.polynomialMul}
\leanok
\uses{aux:coefficient-tensors}
The diagonal tensor, matrix multiplication tensor, and polynomial multiplication tensor are defined by their coefficient formulas on finite index sets.
\end{definition}

\noindent\textit{Auxiliary foundation. Source:} \texttt{material/main.tex:324}.
\textit{Status:} \texttt{formalized}. Main proof path.
\textit{Checked Lean declarations:} \leandecl{MatrixMultiplication.\allowbreak{}Tensor.\allowbreak{}diagonal}, \leandecl{MatrixMultiplication.\allowbreak{}Tensor.\allowbreak{}matrixMul}, \leandecl{MatrixMultiplication.\allowbreak{}Tensor.\allowbreak{}polynomialMul}.
\textit{Formalization note:} Checked coordinate definitions only; no rank identities or bilinear-structure equivalences are claimed.

\renewcommand{\thetheorem}{A17}
\begin{lemma}[Binary partial-sum projections]
\label{aux:binary-partial-sum-projections}
\lean{MatrixMultiplication.Tensor.partialSum_restricts_left,MatrixMultiplication.Tensor.partialSum_restricts_right}
\leanok
For two tensors with the same first index set, their binary partial sum restricts to either summand using the identity map on the first factor.
\end{lemma}

\noindent\textit{Auxiliary foundation. Source:} \texttt{material/main.tex:540}.
\textit{Status:} \texttt{formalized}. Additional manuscript coverage.
\textit{Checked Lean declarations:} \leandecl{MatrixMultiplication.\allowbreak{}Tensor.\allowbreak{}partialSum\_\allowbreak{}restricts\_\allowbreak{}left}, \leandecl{MatrixMultiplication.\allowbreak{}Tensor.\allowbreak{}partialSum\_\allowbreak{}restricts\_\allowbreak{}right}.
\textit{Formalization note:} This proves item 1 of diamondplus-basic in the binary implementation, not all six items.

\begin{proof}
\proves{aux:binary-partial-sum-projections}
\leanok
\uses{aux:coefficient-tensors}
\textit{Checked proof.} Use identity matrices on the first factor and the coordinate projections on the second and third factors.
\end{proof}

\renewcommand{\thetheorem}{A18}
\begin{definition}[Implemented numerical growth hypotheses]
\label{aux:numerical-growth-hypotheses}
\lean{MatrixMultiplication.GrowthHypotheses}
\leanok
For $P:\mathbb N\times\mathbb N\to\mathbb R$, the growth hypotheses assert symmetry on positive indices, $P(1,a)=a$, the concavity inequality $pP(a,b)+P(a,b+p+1)\le(p+1)P(a,b+1)$ for $a,b\ge1$ and $p\ge0$, and $3P(a,a)\le P(a,4a-1)$ for $a\ge1$.
\end{definition}

\noindent\textit{Auxiliary foundation. Source:} \texttt{material/main.tex:1126}.
\textit{Status:} \texttt{formalized}. Main proof path.
\textit{Checked Lean declarations:} \leandecl{MatrixMultiplication.\allowbreak{}GrowthHypotheses}.
\textit{Formalization note:} Checked in Omega/Analysis/Growth.lean. This is a structure of hypotheses, not an axiom asserting their existence for polynomial multiplication value.

\renewcommand{\thetheorem}{A19}
\begin{lemma}[Checked abstract diagonal recurrence]
\label{aux:numerical-recurrence}
\lean{MatrixMultiplication.growthRatio,MatrixMultiplication.GrowthHypotheses.off_diagonal,MatrixMultiplication.GrowthHypotheses.diagonal_step,MatrixMultiplication.GrowthHypotheses.diagonal_recurrence_mul,MatrixMultiplication.GrowthHypotheses.diagonal_recurrence}
\leanok
If $P$ satisfies the numerical growth hypotheses, then for every $a\ge1$, $\frac{(3a+1)(3a+2)}{3a(3a-1)}P(a,a)\le P(a+1,a+1)$.
\end{lemma}

\noindent\textit{Auxiliary foundation. Source:} \texttt{material/main.tex:1153}.
\textit{Status:} \texttt{formalized}. Main proof path.
\textit{Checked Lean declarations:} \leandecl{MatrixMultiplication.\allowbreak{}growthRatio}, \leandecl{MatrixMultiplication.\allowbreak{}GrowthHypotheses.\allowbreak{}off\_\allowbreak{}diagonal}, \leandecl{MatrixMultiplication.\allowbreak{}GrowthHypotheses.\allowbreak{}diagonal\_\allowbreak{}step}, \leandecl{MatrixMultiplication.\allowbreak{}GrowthHypotheses.\allowbreak{}diagonal\_\allowbreak{}recurrence\_\allowbreak{}mul}, \leandecl{MatrixMultiplication.\allowbreak{}GrowthHypotheses.\allowbreak{}diagonal\_\allowbreak{}recurrence}.
\textit{Formalization note:} The generic numerical theorem is checked. Its application to the manuscript P depends on aux:polynomial-value-hypotheses.

\begin{proof}
\proves{aux:numerical-recurrence}
\leanok
\uses{aux:numerical-growth-hypotheses}
\textit{Checked proof.} Apply concavity twice at the parameters used in the manuscript, then use tripling and symmetry. The checked proof clears positive denominators and combines the resulting inequalities.
\end{proof}

\renewcommand{\thetheorem}{A20}
\begin{lemma}[Checked abstract four-thirds growth]
\label{aux:numerical-elementary-growth}
\lean{MatrixMultiplication.growthRatio_pos,MatrixMultiplication.growthRatio_cube_mul,MatrixMultiplication.GrowthHypotheses.diagonal_nonneg_and_cube,MatrixMultiplication.GrowthHypotheses.diagonal_rpow_lower_bound}
\leanok
If $P$ satisfies the numerical growth hypotheses and $a\ge1$, then $P(a,a)\ge a^{4/3}$.
\end{lemma}

\noindent\textit{Auxiliary foundation. Source:} \texttt{material/main.tex:1168}.
\textit{Status:} \texttt{formalized}. Main proof path.
\textit{Checked Lean declarations:} \leandecl{MatrixMultiplication.\allowbreak{}growthRatio\_\allowbreak{}pos}, \leandecl{MatrixMultiplication.\allowbreak{}growthRatio\_\allowbreak{}cube\_\allowbreak{}mul}, \leandecl{MatrixMultiplication.\allowbreak{}GrowthHypotheses.\allowbreak{}diagonal\_\allowbreak{}nonneg\_\allowbreak{}and\_\allowbreak{}cube}, \leandecl{MatrixMultiplication.\allowbreak{}GrowthHypotheses.\allowbreak{}diagonal\_\allowbreak{}rpow\_\allowbreak{}lower\_\allowbreak{}bound}.

\begin{proof}
\proves{aux:numerical-elementary-growth}
\leanok
\uses{aux:numerical-recurrence}
\textit{Checked proof.} A polynomial identity proves the cubed recurrence-factor inequality. Inductively obtain $0\le P(a,a)$ and $a^4\le P(a,a)^3$, then apply monotonicity of real powers.
\end{proof}

\renewcommand{\thetheorem}{A21}
\begin{lemma}[Checked numerical nine-quarters conclusion]
\label{aux:numerical-endpoint}
\lean{MatrixMultiplication.exponent_le_nine_quarters_of_growth}
\leanok
Suppose $P$ satisfies the numerical growth hypotheses, $w>0$, and $P(a,a)^{w/3}\le2a-1$ for every positive integer $a$. Then $w\le9/4$.
\end{lemma}

\noindent\textit{Auxiliary foundation. Source:} \texttt{material/main.tex:1182}.
\textit{Status:} \texttt{formalized}. Main proof path.
\textit{Checked Lean declarations:} \leandecl{MatrixMultiplication.\allowbreak{}exponent\_\allowbreak{}le\_\allowbreak{}nine\_\allowbreak{}quarters\_\allowbreak{}of\_\allowbreak{}growth}.
\textit{Formalization note:} This generic numerical implication is checked independently. The concrete polynomial-value hypotheses and upper bound are now supplied, and the final field-uniform theorem applies it without additional assumptions.

\begin{proof}
\proves{aux:numerical-endpoint}
\leanok
\uses{aux:numerical-elementary-growth,aux:power-growth-comparison}
\textit{Checked proof.} Raise the elementary lower bound to $w/3$ and combine it with the assumed upper bound. The power-comparison lemma gives $4w/9\le1$.
\end{proof}

\renewcommand{\thetheorem}{A22}
\begin{lemma}[Implemented polynomial degeneration witnesses]
\label{aux:polynomial-degeneration-definition}
\lean{MatrixMultiplication.Tensor.liftPolynomial,MatrixMultiplication.Tensor.coefficient,MatrixMultiplication.Tensor.DegenerationWitness,MatrixMultiplication.Tensor.Degenerates,MatrixMultiplication.Tensor.Restricts.degenerates,MatrixMultiplication.Tensor.Degenerates.refl,MatrixMultiplication.Tensor.DegenerationWitness.leadingCoefficient,MatrixMultiplication.Tensor.DegenerationWitness.lowerCoefficient,MatrixMultiplication.Tensor.zero_degenerates_iff}
\leanok
A polynomial degeneration witness records three polynomial factor maps, an order $d$, and a remainder, whose transformed tensor is $X^dS+X^{d+1}E$. Restriction gives an order-zero witness, degeneration is reflexive, the order coefficient is $S$, and all lower coefficients vanish. The zero tensor degenerates only to zero.
\end{lemma}

\noindent\textit{Auxiliary foundation. Source:} \texttt{material/main.tex:312}.
\textit{Status:} \texttt{formalized}. Main proof path.
\textit{Checked Lean declarations:} \leandecl{MatrixMultiplication.\allowbreak{}Tensor.\allowbreak{}liftPolynomial}, \leandecl{MatrixMultiplication.\allowbreak{}Tensor.\allowbreak{}coefficient}, \leandecl{MatrixMultiplication.\allowbreak{}Tensor.\allowbreak{}DegenerationWitness}, \leandecl{MatrixMultiplication.\allowbreak{}Tensor.\allowbreak{}Degenerates}, \leandecl{MatrixMultiplication.\allowbreak{}Tensor.\allowbreak{}Restricts.\allowbreak{}degenerates}, \leandecl{MatrixMultiplication.\allowbreak{}Tensor.\allowbreak{}Degenerates.\allowbreak{}refl}, \leandecl{MatrixMultiplication.\allowbreak{}Tensor.\allowbreak{}DegenerationWitness.\allowbreak{}leadingCoefficient}, \leandecl{MatrixMultiplication.\allowbreak{}Tensor.\allowbreak{}DegenerationWitness.\allowbreak{}lowerCoefficient}, \leandecl{MatrixMultiplication.\allowbreak{}Tensor.\allowbreak{}zero\_\allowbreak{}degenerates\_\allowbreak{}iff}.
\textit{Formalization note:} Checked polynomial definition and elementary properties. Composition and operation compatibility are checked in their support node; equivalence with rational-function specialization is proved in aux:rational-degeneration-equivalence.

\begin{proof}
\proves{aux:polynomial-degeneration-definition}
\leanok
\uses{aux:coefficient-tensors}
\textit{Checked proof.} Use constant polynomial matrices for restriction. Extract coefficients from the witness identity, using the vanishing of coefficients below the displayed powers.
\end{proof}

\renewcommand{\thetheorem}{A23}
\begin{lemma}[Implemented ordinary rank and monotonicity]
\label{aux:ordinary-rank-definition}
\lean{MatrixMultiplication.Tensor.rank,MatrixMultiplication.Tensor.diagonal_restricts_coordinates,MatrixMultiplication.Tensor.diagonal_restricts_of_rank,MatrixMultiplication.Tensor.rank_le_of_diagonal_restricts,MatrixMultiplication.Tensor.rank_mono,MatrixMultiplication.Tensor.rank_eq_of_equivalent,MatrixMultiplication.Tensor.rank_reindex,MatrixMultiplication.Tensor.rank_zero,MatrixMultiplication.Tensor.rank_eq_zero_iff,MatrixMultiplication.Tensor.rank_pure_le_one}
\leanok
Tensor rank is the least $r$ such that $\langle r\rangle$ restricts to $T$. The minimum exists, is attained, and is monotone under restriction and invariant under equivalence and reindexing. A tensor has rank zero exactly when it is zero, and a pure tensor has rank at most one.
\end{lemma}

\noindent\textit{Auxiliary foundation. Source:} \texttt{material/main.tex:324}.
\textit{Status:} \texttt{formalized}. Main proof path.
\textit{Checked Lean declarations:} \leandecl{MatrixMultiplication.\allowbreak{}Tensor.\allowbreak{}rank}, \leandecl{MatrixMultiplication.\allowbreak{}Tensor.\allowbreak{}diagonal\_\allowbreak{}restricts\_\allowbreak{}coordinates}, \leandecl{MatrixMultiplication.\allowbreak{}Tensor.\allowbreak{}diagonal\_\allowbreak{}restricts\_\allowbreak{}of\_\allowbreak{}rank}, \leandecl{MatrixMultiplication.\allowbreak{}Tensor.\allowbreak{}rank\_\allowbreak{}le\_\allowbreak{}of\_\allowbreak{}diagonal\_\allowbreak{}restricts}, \leandecl{MatrixMultiplication.\allowbreak{}Tensor.\allowbreak{}rank\_\allowbreak{}mono}, \leandecl{MatrixMultiplication.\allowbreak{}Tensor.\allowbreak{}rank\_\allowbreak{}eq\_\allowbreak{}of\_\allowbreak{}equivalent}, \leandecl{MatrixMultiplication.\allowbreak{}Tensor.\allowbreak{}rank\_\allowbreak{}reindex}, \leandecl{MatrixMultiplication.\allowbreak{}Tensor.\allowbreak{}rank\_\allowbreak{}zero}, \leandecl{MatrixMultiplication.\allowbreak{}Tensor.\allowbreak{}rank\_\allowbreak{}eq\_\allowbreak{}zero\_\allowbreak{}iff}, \leandecl{MatrixMultiplication.\allowbreak{}Tensor.\allowbreak{}rank\_\allowbreak{}pure\_\allowbreak{}le\_\allowbreak{}one}.

\begin{proof}
\proves{aux:ordinary-rank-definition}
\leanok
\uses{aux:coefficient-tensors,aux:explicit-tensor-definitions}
\textit{Checked proof.} The coordinate expansion gives a restriction of a finite diagonal tensor. Define rank by the least such natural number, then use composition of restrictions and minimality. Prove the zero and pure-tensor cases by explicit maps.
\end{proof}

\renewcommand{\thetheorem}{A24}
\begin{lemma}[Implemented asymptotic rank definition]
\label{aux:asymptotic-rank-definition}
\lean{MatrixMultiplication.Tensor.kroneckerPower,MatrixMultiplication.Tensor.rankRootValues,MatrixMultiplication.Tensor.asymptoticRank,MatrixMultiplication.Tensor.rank_kroneckerPower_one,MatrixMultiplication.Tensor.rankRootValues_nonempty,MatrixMultiplication.Tensor.rankRootValues_bddBelow,MatrixMultiplication.Tensor.asymptoticRank_nonneg,MatrixMultiplication.Tensor.asymptoticRank_le_root,MatrixMultiplication.Tensor.asymptoticRank_le_rank,MatrixMultiplication.Tensor.asymptoticRank_zero}
\leanok
Kronecker powers are indexed by finite tuples. Define $\widetilde R(T)=\inf_{n\ge1}R(T^{\otimes n})^{1/n}$. The set of rank roots is nonempty and bounded below by zero, so $0\le\widetilde R(T)\le R(T^{\otimes n})^{1/n}$ for $n\ge1$. In particular $\widetilde R(T)\le R(T)$ and $\widetilde R(0)=0$.
\end{lemma}

\noindent\textit{Auxiliary foundation. Source:} \texttt{material/main.tex:324}.
\textit{Status:} \texttt{formalized}. Main proof path.
\textit{Checked Lean declarations:} \leandecl{MatrixMultiplication.\allowbreak{}Tensor.\allowbreak{}kroneckerPower}, \leandecl{MatrixMultiplication.\allowbreak{}Tensor.\allowbreak{}rankRootValues}, \leandecl{MatrixMultiplication.\allowbreak{}Tensor.\allowbreak{}asymptoticRank}, \leandecl{MatrixMultiplication.\allowbreak{}Tensor.\allowbreak{}rank\_\allowbreak{}kroneckerPower\_\allowbreak{}one}, \leandecl{MatrixMultiplication.\allowbreak{}Tensor.\allowbreak{}rankRootValues\_\allowbreak{}nonempty}, \leandecl{MatrixMultiplication.\allowbreak{}Tensor.\allowbreak{}rankRootValues\_\allowbreak{}bddBelow}, \leandecl{MatrixMultiplication.\allowbreak{}Tensor.\allowbreak{}asymptoticRank\_\allowbreak{}nonneg}, \leandecl{MatrixMultiplication.\allowbreak{}Tensor.\allowbreak{}asymptoticRank\_\allowbreak{}le\_\allowbreak{}root}, \leandecl{MatrixMultiplication.\allowbreak{}Tensor.\allowbreak{}asymptoticRank\_\allowbreak{}le\_\allowbreak{}rank}, \leandecl{MatrixMultiplication.\allowbreak{}Tensor.\allowbreak{}asymptoticRank\_\allowbreak{}zero}.
\textit{Formalization note:} The infimum definition and elementary bounds are checked here. Rank-root convergence, degeneration monotonicity, and the other asymptotic-rank operation laws are checked in separate support nodes.

\begin{proof}
\proves{aux:asymptotic-rank-definition}
\leanok
\uses{aux:ordinary-rank-definition}
\textit{Checked proof.} Use the first tensor power for nonemptiness and the upper bound by ordinary rank. Real powers of nonnegative ranks provide the lower bound for the infimum.
\end{proof}

\renewcommand{\thetheorem}{A25}
\begin{lemma}[Rank bounds for sums and products]
\label{aux:rank-operation-bounds}
\lean{MatrixMultiplication.Tensor.rank_le_of_sum_pure,MatrixMultiplication.Tensor.exists_rank_decomposition,MatrixMultiplication.Tensor.rank_add_le,MatrixMultiplication.Tensor.rank_sum_le,MatrixMultiplication.Tensor.rank_map_le,MatrixMultiplication.Tensor.rank_kronecker_le,MatrixMultiplication.Tensor.rank_directSum_le}
\leanok
A rank-minimizing diagonal restriction gives a decomposition into pure tensors. Rank is subadditive for finite sums in common spaces and for ordinary direct sums, submultiplicative for Kronecker products, and cannot increase under factor maps.
\end{lemma}

\noindent\textit{Auxiliary foundation. Source:} \texttt{material/main.tex:354}.
\textit{Status:} \texttt{formalized}. Main proof path.
\textit{Checked Lean declarations:} \leandecl{MatrixMultiplication.\allowbreak{}Tensor.\allowbreak{}rank\_\allowbreak{}le\_\allowbreak{}of\_\allowbreak{}sum\_\allowbreak{}pure}, \leandecl{MatrixMultiplication.\allowbreak{}Tensor.\allowbreak{}exists\_\allowbreak{}rank\_\allowbreak{}decomposition}, \leandecl{MatrixMultiplication.\allowbreak{}Tensor.\allowbreak{}rank\_\allowbreak{}add\_\allowbreak{}le}, \leandecl{MatrixMultiplication.\allowbreak{}Tensor.\allowbreak{}rank\_\allowbreak{}sum\_\allowbreak{}le}, \leandecl{MatrixMultiplication.\allowbreak{}Tensor.\allowbreak{}rank\_\allowbreak{}map\_\allowbreak{}le}, \leandecl{MatrixMultiplication.\allowbreak{}Tensor.\allowbreak{}rank\_\allowbreak{}kronecker\_\allowbreak{}le}, \leandecl{MatrixMultiplication.\allowbreak{}Tensor.\allowbreak{}rank\_\allowbreak{}directSum\_\allowbreak{}le}.
\textit{Formalization note:} Checked over every commutative semiring; these are ordinary-rank bounds.

\begin{proof}
\proves{aux:rank-operation-bounds}
\leanok
\uses{aux:ordinary-rank-definition}
\textit{Checked proof.} Concatenate pure-tensor decompositions for addition and multiply their terms for Kronecker products. For direct sums, embed each summand with zero padding and apply subadditivity.
\end{proof}

\renewcommand{\thetheorem}{A26}
\begin{lemma}[Restrictions and binary tensor operations]
\label{aux:restriction-operations}
\lean{MatrixMultiplication.Tensor.map_kronecker,MatrixMultiplication.Tensor.Restricts.kronecker,MatrixMultiplication.Tensor.directSumMap,MatrixMultiplication.Tensor.map_directSum,MatrixMultiplication.Tensor.Restricts.directSum,MatrixMultiplication.Tensor.restricts_directSum_zero,MatrixMultiplication.Tensor.restricts_zero_directSum}
\leanok
Factor maps commute with Kronecker products and ordinary direct sums. Consequently restriction respects both operations, and zero padding a tensor is a restriction.
\end{lemma}

\noindent\textit{Auxiliary foundation. Source:} \texttt{material/main.tex:299}.
\textit{Status:} \texttt{formalized}. Main proof path.
\textit{Checked Lean declarations:} \leandecl{MatrixMultiplication.\allowbreak{}Tensor.\allowbreak{}map\_\allowbreak{}kronecker}, \leandecl{MatrixMultiplication.\allowbreak{}Tensor.\allowbreak{}Restricts.\allowbreak{}kronecker}, \leandecl{MatrixMultiplication.\allowbreak{}Tensor.\allowbreak{}directSumMap}, \leandecl{MatrixMultiplication.\allowbreak{}Tensor.\allowbreak{}map\_\allowbreak{}directSum}, \leandecl{MatrixMultiplication.\allowbreak{}Tensor.\allowbreak{}Restricts.\allowbreak{}directSum}, \leandecl{MatrixMultiplication.\allowbreak{}Tensor.\allowbreak{}restricts\_\allowbreak{}directSum\_\allowbreak{}zero}, \leandecl{MatrixMultiplication.\allowbreak{}Tensor.\allowbreak{}restricts\_\allowbreak{}zero\_\allowbreak{}directSum}.

\begin{proof}
\proves{aux:restriction-operations}
\leanok
\uses{aux:coefficient-tensors}
\textit{Checked proof.} Use Kronecker products of the factor matrices for tensor products and block-diagonal matrices for direct sums. Expand the finite coordinate sums.
\end{proof}

\renewcommand{\thetheorem}{A27}
\begin{lemma}[Tensor-power identities and rank bounds]
\label{aux:power-rank-submultiplicativity}
\lean{MatrixMultiplication.Tensor.kroneckerPower_add_reindex,MatrixMultiplication.Tensor.rank_kroneckerPower_add_le,MatrixMultiplication.Tensor.rank_kroneckerPower_zero_le,MatrixMultiplication.Tensor.rank_kroneckerPower_le}
\leanok
Splitting finite tuples identifies $T^{\otimes(m+n)}$ with $T^{\otimes m}\otimes T^{\otimes n}$. Hence $R(T^{\otimes(m+n)})\le R(T^{\otimes m})R(T^{\otimes n})$ and $R(T^{\otimes n})\le R(T)^n$. The zeroth power has rank at most one.
\end{lemma}

\noindent\textit{Auxiliary foundation. Source:} \texttt{material/main.tex:354}.
\textit{Status:} \texttt{formalized}. Main proof path.
\textit{Checked Lean declarations:} \leandecl{MatrixMultiplication.\allowbreak{}Tensor.\allowbreak{}kroneckerPower\_\allowbreak{}add\_\allowbreak{}reindex}, \leandecl{MatrixMultiplication.\allowbreak{}Tensor.\allowbreak{}rank\_\allowbreak{}kroneckerPower\_\allowbreak{}add\_\allowbreak{}le}, \leandecl{MatrixMultiplication.\allowbreak{}Tensor.\allowbreak{}rank\_\allowbreak{}kroneckerPower\_\allowbreak{}zero\_\allowbreak{}le}, \leandecl{MatrixMultiplication.\allowbreak{}Tensor.\allowbreak{}rank\_\allowbreak{}kroneckerPower\_\allowbreak{}le}.

\begin{proof}
\proves{aux:power-rank-submultiplicativity}
\leanok
\uses{aux:asymptotic-rank-definition,aux:rank-operation-bounds}
\textit{Checked proof.} Reindex by concatenation and splitting of finite tuples, use invariance of rank under reindexing, and apply the binary product bound. Induct for the bound by $R(T)^n$.
\end{proof}

\renewcommand{\thetheorem}{A28}
\begin{lemma}[Composition and operations for polynomial degeneration]
\label{aux:polynomial-degeneration-operations}
\lean{MatrixMultiplication.Tensor.matrixPower,MatrixMultiplication.Tensor.map_kroneckerPower,MatrixMultiplication.Tensor.DegenerationWitness.kroneckerPower,MatrixMultiplication.Tensor.DegenerationWitness.kroneckerPower_order,MatrixMultiplication.Tensor.Degenerates.kroneckerPower,MatrixMultiplication.Tensor.DegenerationWitness.trans,MatrixMultiplication.Tensor.Degenerates.trans,MatrixMultiplication.Tensor.DegenerationWitness.kronecker,MatrixMultiplication.Tensor.Degenerates.kronecker,MatrixMultiplication.Tensor.DegenerationWitness.shift,MatrixMultiplication.Tensor.DegenerationWitness.directSum,MatrixMultiplication.Tensor.Degenerates.directSum}
\leanok
Polynomial degeneration is transitive and respects tensor powers, binary Kronecker products, and ordinary direct sums. A witness of order $d$ gives a witness of order $dn$ for the $n$th tensor powers. Binary product witnesses have order $d+e$; scalar shifts align direct-sum orders.
\end{lemma}

\noindent\textit{Auxiliary foundation. Source:} \texttt{material/main.tex:312}.
\textit{Status:} \texttt{formalized}. Main proof path.
\textit{Checked Lean declarations:} \leandecl{MatrixMultiplication.\allowbreak{}Tensor.\allowbreak{}matrixPower}, \leandecl{MatrixMultiplication.\allowbreak{}Tensor.\allowbreak{}map\_\allowbreak{}kroneckerPower}, \leandecl{MatrixMultiplication.\allowbreak{}Tensor.\allowbreak{}DegenerationWitness.\allowbreak{}kroneckerPower}, \leandecl{MatrixMultiplication.\allowbreak{}Tensor.\allowbreak{}DegenerationWitness.\allowbreak{}kroneckerPower\_\allowbreak{}order}, \leandecl{MatrixMultiplication.\allowbreak{}Tensor.\allowbreak{}Degenerates.\allowbreak{}kroneckerPower}, \leandecl{MatrixMultiplication.\allowbreak{}Tensor.\allowbreak{}DegenerationWitness.\allowbreak{}trans}, \leandecl{MatrixMultiplication.\allowbreak{}Tensor.\allowbreak{}Degenerates.\allowbreak{}trans}, \leandecl{MatrixMultiplication.\allowbreak{}Tensor.\allowbreak{}DegenerationWitness.\allowbreak{}kronecker}, \leandecl{MatrixMultiplication.\allowbreak{}Tensor.\allowbreak{}Degenerates.\allowbreak{}kronecker}, \leandecl{MatrixMultiplication.\allowbreak{}Tensor.\allowbreak{}DegenerationWitness.\allowbreak{}shift}, \leandecl{MatrixMultiplication.\allowbreak{}Tensor.\allowbreak{}DegenerationWitness.\allowbreak{}directSum}, \leandecl{MatrixMultiplication.\allowbreak{}Tensor.\allowbreak{}Degenerates.\allowbreak{}directSum}.
\textit{Formalization note:} All witness constructions are checked over commutative semirings. Over fields the polynomial relation agrees with rational-function specialization by aux:rational-degeneration-equivalence.

\begin{proof}
\proves{aux:polynomial-degeneration-operations}
\leanok
\uses{aux:polynomial-degeneration-definition,aux:asymptotic-rank-definition,aux:restriction-operations}
\textit{Checked proof.} For composition with a witness of order $e$, substitute $X^{e+1}$ in the first witness before composing factor maps. Expand products with their remainders. For direct sums, shift orders to a common value and use block-diagonal maps.
\end{proof}

\renewcommand{\thetheorem}{A29}
\begin{lemma}[Rank bound for polynomial coefficient extraction]
\label{aux:coefficient-rank-bound}
\manuscriptproof
\lean{MatrixMultiplication.Tensor.matrixCoefficient,MatrixMultiplication.Tensor.coefficient_map_liftPolynomial,MatrixMultiplication.Tensor.rank_coefficient_map_le,MatrixMultiplication.Tensor.DegenerationWitness.rank_le,MatrixMultiplication.Tensor.rank_coefficient_map_le_choose,MatrixMultiplication.Tensor.DegenerationWitness.rank_le_choose,MatrixMultiplication.Tensor.DegenerationWitness.rank_kroneckerPower_le_choose}
\leanok
The coefficient of order $d$ of a polynomial restriction is a sum of $\binom{d+2}{2}$ ordinary restrictions, so its rank is at most $\binom{d+2}{2}R(T)$. A witness of order $d$ gives $R(S^{\otimes N})\le\binom{dN+2}{2}R(T^{\otimes N})$. The coarser quadratic bounds remain available.
\end{lemma}

\noindent\textit{Auxiliary foundation. Source:} \texttt{material/main.tex:370}.
\textit{Status:} \texttt{formalized}. Main proof path.
\textit{Proof supplied in manuscript.}
\textit{Checked Lean declarations:} \leandecl{MatrixMultiplication.\allowbreak{}Tensor.\allowbreak{}matrixCoefficient}, \leandecl{MatrixMultiplication.\allowbreak{}Tensor.\allowbreak{}coefficient\_\allowbreak{}map\_\allowbreak{}liftPolynomial}, \leandecl{MatrixMultiplication.\allowbreak{}Tensor.\allowbreak{}rank\_\allowbreak{}coefficient\_\allowbreak{}map\_\allowbreak{}le}, \leandecl{MatrixMultiplication.\allowbreak{}Tensor.\allowbreak{}DegenerationWitness.\allowbreak{}rank\_\allowbreak{}le}, \leandecl{MatrixMultiplication.\allowbreak{}Tensor.\allowbreak{}rank\_\allowbreak{}coefficient\_\allowbreak{}map\_\allowbreak{}le\_\allowbreak{}choose}, \leandecl{MatrixMultiplication.\allowbreak{}Tensor.\allowbreak{}DegenerationWitness.\allowbreak{}rank\_\allowbreak{}le\_\allowbreak{}choose}, \leandecl{MatrixMultiplication.\allowbreak{}Tensor.\allowbreak{}DegenerationWitness.\allowbreak{}rank\_\allowbreak{}kroneckerPower\_\allowbreak{}le\_\allowbreak{}choose}.
\textit{Formalization note:} Uses formal coefficients only, with no interpolation or assumption that the coefficient field is infinite.

\begin{proof}
\proves{aux:coefficient-rank-bound}
\leanok
\uses{aux:polynomial-degeneration-definition,aux:rank-operation-bounds,aux:polynomial-degeneration-operations}
\textit{Checked proof.} Expand into two antidiagonal sums and count their total cardinality exactly. Apply rank subadditivity and restriction monotonicity, then apply the same estimate to the tensor-power witness.
\end{proof}

\renewcommand{\thetheorem}{A30}
\begin{lemma}[Polynomial overhead disappears in rank-root limits]
\label{aux:polynomial-overhead-limit}
\lean{MatrixMultiplication.affine_root_tendsto_one,MatrixMultiplication.polynomial_overhead_root_tendsto_one,MatrixMultiplication.submultiplicative_limit_le_of_polynomial_bound}
\leanok
For fixed nonnegative integers $d,k$, $((dn+1)^k)^{1/n}\to1$. If nonnegative submultiplicative real sequences satisfy $v_n\le(dn+1)^ku_n$ for $n\ge1$, then their rank-root limits satisfy $\lim v_n^{1/n}\le\lim u_n^{1/n}$.
\end{lemma}

\noindent\textit{Auxiliary foundation. Source:} \texttt{material/main.tex:361}.
\textit{Status:} \texttt{formalized}. Main proof path.
\textit{Checked Lean declarations:} \leandecl{MatrixMultiplication.\allowbreak{}affine\_\allowbreak{}root\_\allowbreak{}tendsto\_\allowbreak{}one}, \leandecl{MatrixMultiplication.\allowbreak{}polynomial\_\allowbreak{}overhead\_\allowbreak{}root\_\allowbreak{}tendsto\_\allowbreak{}one}, \leandecl{MatrixMultiplication.\allowbreak{}submultiplicative\_\allowbreak{}limit\_\allowbreak{}le\_\allowbreak{}of\_\allowbreak{}polynomial\_\allowbreak{}bound}.

\begin{proof}
\proves{aux:polynomial-overhead-limit}
\leanok
\textit{Checked proof.} Use the real-power asymptotic limit for an affine function, then a fixed integer power. Apply the submultiplicative-sequence limit theorem and take limits in the root inequality.
\end{proof}

\renewcommand{\thetheorem}{A31}
\begin{lemma}[Asymptotic rank is the limit of rank roots]
\label{aux:asymptotic-rank-limit}
\lean{MatrixMultiplication.Tensor.rankPowerSequence_submultiplicative,MatrixMultiplication.Tensor.asymptoticRank_eq_submultiplicative_limit,MatrixMultiplication.Tensor.asymptoticRank_eq_limit,MatrixMultiplication.Tensor.asymptoticRank_le_of_power_rank_le}
\leanok
The infimum defining asymptotic rank equals the limit: $\widetilde R(T)=\lim_{n\to\infty}R(T^{\otimes n})^{1/n}$. In particular, a bound $R(S^{\otimes n})\le(dn+1)^kR(T^{\otimes n})$ for all $n\ge1$ implies $\widetilde R(S)\le\widetilde R(T)$.
\end{lemma}

\noindent\textit{Auxiliary foundation. Source:} \texttt{material/main.tex:354}.
\textit{Status:} \texttt{formalized}. Main proof path.
\textit{Checked Lean declarations:} \leandecl{MatrixMultiplication.\allowbreak{}Tensor.\allowbreak{}rankPowerSequence\_\allowbreak{}submultiplicative}, \leandecl{MatrixMultiplication.\allowbreak{}Tensor.\allowbreak{}asymptoticRank\_\allowbreak{}eq\_\allowbreak{}submultiplicative\_\allowbreak{}limit}, \leandecl{MatrixMultiplication.\allowbreak{}Tensor.\allowbreak{}asymptoticRank\_\allowbreak{}eq\_\allowbreak{}limit}, \leandecl{MatrixMultiplication.\allowbreak{}Tensor.\allowbreak{}asymptoticRank\_\allowbreak{}le\_\allowbreak{}of\_\allowbreak{}power\_\allowbreak{}rank\_\allowbreak{}le}.

\begin{proof}
\proves{aux:asymptotic-rank-limit}
\leanok
\uses{aux:power-rank-submultiplicativity,aux:polynomial-overhead-limit}
\textit{Checked proof.} Identify the defining infimum with the limit in mathlib for a nonnegative submultiplicative sequence. Apply the polynomial-overhead comparison theorem to the two rank sequences.
\end{proof}

\renewcommand{\thetheorem}{A32}
\begin{lemma}[Degeneration cannot increase asymptotic rank]
\label{aux:asymptotic-rank-degeneration}
\manuscriptproof
\lean{MatrixMultiplication.Tensor.DegenerationWitness.rank_kroneckerPower_le,MatrixMultiplication.Tensor.asymptoticRank_mono_degeneration,MatrixMultiplication.Tensor.asymptoticRank_mono,MatrixMultiplication.Tensor.asymptoticRank_eq_of_equivalent}
\leanok
If $T\trianglerighteq S$, then $\widetilde R(S)\le\widetilde R(T)$. More explicitly, a polynomial witness of order $d$ gives $R(S^{\otimes n})\le(dn+1)^2R(T^{\otimes n})$ for every $n\ge0$. Asymptotic rank is therefore also monotone under restriction and invariant under equivalence.
\end{lemma}

\noindent\textit{Auxiliary foundation. Source:} \texttt{material/main.tex:361}.
\textit{Status:} \texttt{formalized}. Main proof path.
\textit{Proof supplied in manuscript.}
\textit{Checked Lean declarations:} \leandecl{MatrixMultiplication.\allowbreak{}Tensor.\allowbreak{}DegenerationWitness.\allowbreak{}rank\_\allowbreak{}kroneckerPower\_\allowbreak{}le}, \leandecl{MatrixMultiplication.\allowbreak{}Tensor.\allowbreak{}asymptoticRank\_\allowbreak{}mono\_\allowbreak{}degeneration}, \leandecl{MatrixMultiplication.\allowbreak{}Tensor.\allowbreak{}asymptoticRank\_\allowbreak{}mono}, \leandecl{MatrixMultiplication.\allowbreak{}Tensor.\allowbreak{}asymptoticRank\_\allowbreak{}eq\_\allowbreak{}of\_\allowbreak{}equivalent}.
\textit{Formalization note:} Checked for the polynomial-witness relation over all commutative semirings, hence over every field.

\begin{proof}
\proves{aux:asymptotic-rank-degeneration}
\leanok
\uses{aux:coefficient-rank-bound,aux:polynomial-degeneration-operations,aux:asymptotic-rank-limit}
\textit{Checked proof.} Take tensor powers of the witness, whose order becomes $dn$, and extract its leading coefficient. Apply the coefficient-rank bound and then remove the polynomial overhead in the rank-root limit.
\end{proof}

\renewcommand{\thetheorem}{A33}
\begin{lemma}[Border rank and its asymptotic lower bound]
\label{aux:border-rank-bound}
\manuscriptproof
\lean{MatrixMultiplication.Tensor.exists_diagonal_degenerates,MatrixMultiplication.Tensor.borderRank,MatrixMultiplication.Tensor.diagonal_degenerates_of_borderRank,MatrixMultiplication.Tensor.borderRank_le_of_diagonal_degenerates,MatrixMultiplication.Tensor.borderRank_le_rank,MatrixMultiplication.Tensor.borderRank_mono_degeneration,MatrixMultiplication.Tensor.asymptoticRank_le_borderRank}
\leanok
Border rank is the least size of a diagonal tensor degenerating to $T$, and this minimum is attained. It is monotone under degeneration and satisfies $\widetilde R(T)\le\underline R(T)\le R(T)$.
\end{lemma}

\noindent\textit{Auxiliary foundation. Source:} \texttt{material/main.tex:361}.
\textit{Status:} \texttt{formalized}. Main proof path.
\textit{Proof supplied in manuscript.}
\textit{Checked Lean declarations:} \leandecl{MatrixMultiplication.\allowbreak{}Tensor.\allowbreak{}exists\_\allowbreak{}diagonal\_\allowbreak{}degenerates}, \leandecl{MatrixMultiplication.\allowbreak{}Tensor.\allowbreak{}borderRank}, \leandecl{MatrixMultiplication.\allowbreak{}Tensor.\allowbreak{}diagonal\_\allowbreak{}degenerates\_\allowbreak{}of\_\allowbreak{}borderRank}, \leandecl{MatrixMultiplication.\allowbreak{}Tensor.\allowbreak{}borderRank\_\allowbreak{}le\_\allowbreak{}of\_\allowbreak{}diagonal\_\allowbreak{}degenerates}, \leandecl{MatrixMultiplication.\allowbreak{}Tensor.\allowbreak{}borderRank\_\allowbreak{}le\_\allowbreak{}rank}, \leandecl{MatrixMultiplication.\allowbreak{}Tensor.\allowbreak{}borderRank\_\allowbreak{}mono\_\allowbreak{}degeneration}, \leandecl{MatrixMultiplication.\allowbreak{}Tensor.\allowbreak{}asymptoticRank\_\allowbreak{}le\_\allowbreak{}borderRank}.
\textit{Formalization note:} Border rank uses polynomial witnesses, proved equivalent to rational-function degeneration. The subrank variants are checked in the separate subrank-invariants node.

\begin{proof}
\proves{aux:border-rank-bound}
\leanok
\uses{aux:ordinary-rank-definition,aux:polynomial-degeneration-operations,aux:asymptotic-rank-degeneration}
\textit{Checked proof.} An ordinary rank decomposition supplies a diagonal degeneration source. Take the least such size, use transitivity for border-rank monotonicity, and apply asymptotic-rank monotonicity to a minimizing source.
\end{proof}

\renewcommand{\thetheorem}{A34}
\begin{definition}[Implemented matrix multiplication and exponent definitions]
\label{aux:matrix-multiplication-exponent}
\lean{MatrixMultiplication.Tensor.matrixMul,MatrixMultiplication.omega}
\leanok
\uses{aux:explicit-tensor-definitions,aux:asymptotic-rank-definition}
For positive integers $\ell,m,n$, define $\langle\ell,m,n\rangle=\sum_{i,j,k}e_{ij}\otimes e_{jk}\otimes e_{ik}$. For every field $F$, set $\omega_F=\log_2\widetilde R_F(\langle2,2,2\rangle)$.
\end{definition}

\noindent\textit{Auxiliary foundation. Source:} \texttt{material/main.tex:404}.
\textit{Status:} \texttt{formalized}. Main proof path.
\textit{Checked Lean declarations:} \leandecl{MatrixMultiplication.\allowbreak{}Tensor.\allowbreak{}matrixMul}, \leandecl{MatrixMultiplication.\allowbreak{}omega}.
\textit{Formalization note:} The coordinate tensor and rank-based exponent are checked. The special-shape equivalences are checked in the full matrix-multiplication definition node, and the bilinear interpretation has its own correspondence node.

\renewcommand{\thetheorem}{A35}
\begin{lemma}[Matrix rank of a Kronecker product]
\label{aux:matrix-rank-kronecker}
\lean{MatrixMultiplication.matrix_rank_kronecker}
\leanok
For arbitrary finite rectangular matrices $A,B$ over a field, $\operatorname{rank}(A\otimes B)=\operatorname{rank}(A)\operatorname{rank}(B)$.
\end{lemma}

\noindent\textit{Auxiliary foundation. Source:} \texttt{material/main.tex:393}.
\textit{Status:} \texttt{formalized}. Main proof path.
\textit{Checked Lean declarations:} \leandecl{MatrixMultiplication.\allowbreak{}matrix\_\allowbreak{}rank\_\allowbreak{}kronecker}.
\textit{Formalization note:} No characteristic or field-cardinality hypothesis is used; rectangular and empty shapes are included.

\begin{proof}
\proves{aux:matrix-rank-kronecker}
\leanok
\textit{Checked proof.} Identify the Kronecker matrix map with the tensor product of the two matrix maps using finite-function coordinate equivalences. The image is the tensor product of their images, whose dimensions multiply.
\end{proof}

\renewcommand{\thetheorem}{A36}
\begin{lemma}[Implemented flattenings and their product ranks]
\label{aux:flattening-operations}
\lean{MatrixMultiplication.Tensor.flatteningFirst,MatrixMultiplication.Tensor.flatteningSecond,MatrixMultiplication.Tensor.flatteningThird,MatrixMultiplication.Tensor.flatteningRank,MatrixMultiplication.Tensor.flatteningRank_le_rank,MatrixMultiplication.Tensor.flatteningRank_reindex,MatrixMultiplication.Tensor.flatteningFirst_kronecker,MatrixMultiplication.Tensor.flatteningSecond_kronecker,MatrixMultiplication.Tensor.flatteningThird_kronecker,MatrixMultiplication.Tensor.flatteningRank_kronecker}
\leanok
Each tensor has three flattening matrices, with one factor indexing columns and the other two indexing rows. For each factor $i$, $R^{(i)}(T)\le R(T)$. These ranks are invariant under coordinate reindexing and satisfy $R^{(i)}(T\otimes S)=R^{(i)}(T)R^{(i)}(S)$.
\end{lemma}

\noindent\textit{Auxiliary foundation. Source:} \texttt{material/main.tex:393}.
\textit{Status:} \texttt{formalized}. Main proof path.
\textit{Checked Lean declarations:} \leandecl{MatrixMultiplication.\allowbreak{}Tensor.\allowbreak{}flatteningFirst}, \leandecl{MatrixMultiplication.\allowbreak{}Tensor.\allowbreak{}flatteningSecond}, \leandecl{MatrixMultiplication.\allowbreak{}Tensor.\allowbreak{}flatteningThird}, \leandecl{MatrixMultiplication.\allowbreak{}Tensor.\allowbreak{}flatteningRank}, \leandecl{MatrixMultiplication.\allowbreak{}Tensor.\allowbreak{}flatteningRank\_\allowbreak{}le\_\allowbreak{}rank}, \leandecl{MatrixMultiplication.\allowbreak{}Tensor.\allowbreak{}flatteningRank\_\allowbreak{}reindex}, \leandecl{MatrixMultiplication.\allowbreak{}Tensor.\allowbreak{}flatteningFirst\_\allowbreak{}kronecker}, \leandecl{MatrixMultiplication.\allowbreak{}Tensor.\allowbreak{}flatteningSecond\_\allowbreak{}kronecker}, \leandecl{MatrixMultiplication.\allowbreak{}Tensor.\allowbreak{}flatteningThird\_\allowbreak{}kronecker}, \leandecl{MatrixMultiplication.\allowbreak{}Tensor.\allowbreak{}flatteningRank\_\allowbreak{}kronecker}.
\textit{Formalization note:} The Lean factor indices 0, 1, 2 correspond to the paper indices 1, 2, 3. Restriction and degeneration monotonicity are tracked separately in the full flattening-calculus node.

\begin{proof}
\proves{aux:flattening-operations}
\leanok
\uses{aux:rank-operation-bounds,aux:matrix-rank-kronecker}
\textit{Checked proof.} A rank decomposition factors each flattening through a space of dimension $R(T)$. A tensor-product flattening is a coordinate reindexing of the matrix Kronecker product, so its rank is given by matrix-rank multiplicativity.
\end{proof}

\renewcommand{\thetheorem}{A37}
\begin{lemma}[Flattening ranks of tensor powers]
\label{aux:flattening-powers}
\lean{MatrixMultiplication.Tensor.flatteningRank_kroneckerPower_zero,MatrixMultiplication.Tensor.flatteningRank_kroneckerPower_one,MatrixMultiplication.Tensor.flatteningRank_kroneckerPower,MatrixMultiplication.Tensor.flatteningRank_le_borderRank}
\leanok
For every $n\ge0$ and each factor $i$, $R^{(i)}(T^{\otimes n})=R^{(i)}(T)^n$. In particular, every flattening of the zeroth power has rank one, including when a factor of $T$ has dimension zero.
\end{lemma}

\noindent\textit{Auxiliary foundation. Source:} \texttt{material/main.tex:393}.
\textit{Status:} \texttt{formalized}. Main proof path.
\textit{Checked Lean declarations:} \leandecl{MatrixMultiplication.\allowbreak{}Tensor.\allowbreak{}flatteningRank\_\allowbreak{}kroneckerPower\_\allowbreak{}zero}, \leandecl{MatrixMultiplication.\allowbreak{}Tensor.\allowbreak{}flatteningRank\_\allowbreak{}kroneckerPower\_\allowbreak{}one}, \leandecl{MatrixMultiplication.\allowbreak{}Tensor.\allowbreak{}flatteningRank\_\allowbreak{}kroneckerPower}, \leandecl{MatrixMultiplication.\allowbreak{}Tensor.\allowbreak{}flatteningRank\_\allowbreak{}le\_\allowbreak{}borderRank}.

\begin{proof}
\proves{aux:flattening-powers}
\leanok
\uses{aux:flattening-operations,aux:power-rank-submultiplicativity}
\textit{Checked proof.} At power zero the coordinate types are singleton types, giving a rank-one matrix. At power one remove singleton tuples. For the induction step split tuples, use invariance under reindexing, and apply the binary product identity.
\end{proof}

\renewcommand{\thetheorem}{A38}
\begin{lemma}[Exact rank invariants of diagonal tensors]
\label{aux:diagonal-rank-normalization}
\lean{MatrixMultiplication.Tensor.flatteningFirst_rank_diagonal,MatrixMultiplication.Tensor.flatteningSecond_rank_diagonal,MatrixMultiplication.Tensor.flatteningThird_rank_diagonal,MatrixMultiplication.Tensor.flatteningRank_diagonal,MatrixMultiplication.Tensor.rank_diagonal,MatrixMultiplication.Tensor.borderRank_diagonal,MatrixMultiplication.Tensor.asymptoticRank_diagonal}
\leanok
For every integer $r\ge0$, all three flattening ranks of $\langle r\rangle$ equal $r$, and $R(\langle r\rangle)=\underline R(\langle r\rangle)=\widetilde R(\langle r\rangle)=r$.
\end{lemma}

\noindent\textit{Auxiliary foundation. Source:} \texttt{material/main.tex:324}.
\textit{Status:} \texttt{formalized}. Main proof path.
\textit{Checked Lean declarations:} \leandecl{MatrixMultiplication.\allowbreak{}Tensor.\allowbreak{}flatteningFirst\_\allowbreak{}rank\_\allowbreak{}diagonal}, \leandecl{MatrixMultiplication.\allowbreak{}Tensor.\allowbreak{}flatteningSecond\_\allowbreak{}rank\_\allowbreak{}diagonal}, \leandecl{MatrixMultiplication.\allowbreak{}Tensor.\allowbreak{}flatteningThird\_\allowbreak{}rank\_\allowbreak{}diagonal}, \leandecl{MatrixMultiplication.\allowbreak{}Tensor.\allowbreak{}flatteningRank\_\allowbreak{}diagonal}, \leandecl{MatrixMultiplication.\allowbreak{}Tensor.\allowbreak{}rank\_\allowbreak{}diagonal}, \leandecl{MatrixMultiplication.\allowbreak{}Tensor.\allowbreak{}borderRank\_\allowbreak{}diagonal}, \leandecl{MatrixMultiplication.\allowbreak{}Tensor.\allowbreak{}asymptoticRank\_\allowbreak{}diagonal}.
\textit{Formalization note:} Checked over every field, including r = 0. Border rank uses the implemented polynomial-witness definition.

\begin{proof}
\proves{aux:diagonal-rank-normalization}
\leanok
\uses{lem:flattening-lb,aux:ordinary-rank-definition,aux:border-rank-bound}
\textit{Checked proof.} Select the diagonal rows of each flattening to obtain an identity submatrix. The resulting lower bounds match the elementary upper bound on the rank of a diagonal tensor.
\end{proof}

\renewcommand{\thetheorem}{A39}
\begin{lemma}[First flattening rank of matrix multiplication]
\label{aux:matrix-multiplication-flattening}
\lean{MatrixMultiplication.Tensor.flatteningFirst_rank_matrixMul,MatrixMultiplication.Tensor.flatteningRank_matrixMul_zero}
\leanok
For integers $\ell,m\ge0$ and $n\ge1$, the first flattening of $\langle\ell,m,n\rangle$ has rank $\ell m$.
\end{lemma}

\noindent\textit{Auxiliary foundation. Source:} \texttt{material/main.tex:404}.
\textit{Status:} \texttt{formalized}. Main proof path.
\textit{Checked Lean declarations:} \leandecl{MatrixMultiplication.\allowbreak{}Tensor.\allowbreak{}flatteningFirst\_\allowbreak{}rank\_\allowbreak{}matrixMul}, \leandecl{MatrixMultiplication.\allowbreak{}Tensor.\allowbreak{}flatteningRank\_\allowbreak{}matrixMul\_\allowbreak{}zero}.
\textit{Formalization note:} This node proves the first flattening formula; no formulas for the other two matrix-multiplication flattenings are asserted here.

\begin{proof}
\proves{aux:matrix-multiplication-flattening}
\leanok
\uses{aux:flattening-operations,aux:explicit-tensor-definitions}
\textit{Checked proof.} Fix one index in the nonempty third matrix dimension. The corresponding rows give an identity submatrix of size $\ell m$, matching the number of columns.
\end{proof}

\renewcommand{\thetheorem}{A40}
\begin{lemma}[The lower bound two for the exponent]
\label{aux:omega-lower-bound}
\lean{MatrixMultiplication.four_le_asymptoticRank_matrixMul_two,MatrixMultiplication.two_le_omega}
\leanok
For every field $F$, $\widetilde R_F(\langle2,2,2\rangle)\ge4$, hence $\omega_F\ge2$.
\end{lemma}

\noindent\textit{Auxiliary foundation. Source:} \texttt{material/main.tex:404}.
\textit{Status:} \texttt{formalized}. Main proof path.
\textit{Checked Lean declarations:} \leandecl{MatrixMultiplication.\allowbreak{}four\_\allowbreak{}le\_\allowbreak{}asymptoticRank\_\allowbreak{}matrixMul\_\allowbreak{}two}, \leandecl{MatrixMultiplication.\allowbreak{}two\_\allowbreak{}le\_\allowbreak{}omega}.
\textit{Formalization note:} Only the field assumption is required. The upper bound is also checked in the full omega-bounds node.

\begin{proof}
\proves{aux:omega-lower-bound}
\leanok
\uses{aux:matrix-multiplication-exponent,aux:matrix-multiplication-flattening,lem:flattening-lb}
\textit{Checked proof.} Apply the first-flattening lower bound to the two-by-two matrix-multiplication tensor, whose first flattening has rank four. The base-two logarithm is increasing on positive reals.
\end{proof}

\renewcommand{\thetheorem}{A41}
\begin{lemma}[Asymptotic rank of powers and products]
\label{aux:asymptotic-rank-powers-products}
\lean{MatrixMultiplication.Tensor.asymptoticRank_pow,MatrixMultiplication.Tensor.asymptoticRank_kronecker_le,MatrixMultiplication.Tensor.rank_kroneckerPower_kroneckerPower,MatrixMultiplication.Tensor.rank_kroneckerPower_kronecker}
\leanok
Over every field, $\widetilde R(T^{\otimes k})=\widetilde R(T)^k$ for all integers $k\ge0$, and $\widetilde R(T\otimes S)\le\widetilde R(T)\widetilde R(S)$.
\end{lemma}

\noindent\textit{Auxiliary foundation. Source:} \texttt{material/main.tex:354}.
\textit{Status:} \texttt{formalized}. Main proof path.
\textit{Checked Lean declarations:} \leandecl{MatrixMultiplication.\allowbreak{}Tensor.\allowbreak{}asymptoticRank\_\allowbreak{}pow}, \leandecl{MatrixMultiplication.\allowbreak{}Tensor.\allowbreak{}asymptoticRank\_\allowbreak{}kronecker\_\allowbreak{}le}, \leandecl{MatrixMultiplication.\allowbreak{}Tensor.\allowbreak{}rank\_\allowbreak{}kroneckerPower\_\allowbreak{}kroneckerPower}, \leandecl{MatrixMultiplication.\allowbreak{}Tensor.\allowbreak{}rank\_\allowbreak{}kroneckerPower\_\allowbreak{}kronecker}.
\textit{Formalization note:} Checked including power zero. The power law holds over nontrivial commutative semirings; the product bound holds over arbitrary commutative semirings.

\begin{proof}
\proves{aux:asymptotic-rank-powers-products}
\leanok
\uses{aux:asymptotic-rank-limit,aux:power-rank-submultiplicativity,aux:rank-operation-bounds}
\textit{Checked proof.} Regroup power coordinates explicitly. For $k>0$, apply the root-limit theorem along the subsequence $Nk$; handle the empty power separately. For products, take limits in the ordinary rank product bound.
\end{proof}

\renewcommand{\thetheorem}{A42}
\begin{lemma}[Asymptotic rank of sums]
\label{aux:asymptotic-rank-sums}
\lean{MatrixMultiplication.Tensor.asymptoticRank_directSum_le,MatrixMultiplication.Tensor.asymptoticRank_add_le,MatrixMultiplication.Tensor.exists_rank_power_exponential_bound,MatrixMultiplication.Tensor.rank_kroneckerPower_add_le_sum,MatrixMultiplication.exists_pos_exponential_bound,MatrixMultiplication.root_limit_le_of_exponential_bound}
\leanok
Asymptotic rank is subadditive: $\widetilde R(T\oplus S)\le\widetilde R(T)+\widetilde R(S)$. For tensors with the same coordinate sets, also $\widetilde R(T+S)\le\widetilde R(T)+\widetilde R(S)$.
\end{lemma}

\noindent\textit{Auxiliary foundation. Source:} \texttt{material/main.tex:354}.
\textit{Status:} \texttt{formalized}. Additional manuscript coverage.
\textit{Checked Lean declarations:} \leandecl{MatrixMultiplication.\allowbreak{}Tensor.\allowbreak{}asymptoticRank\_\allowbreak{}directSum\_\allowbreak{}le}, \leandecl{MatrixMultiplication.\allowbreak{}Tensor.\allowbreak{}asymptoticRank\_\allowbreak{}add\_\allowbreak{}le}, \leandecl{MatrixMultiplication.\allowbreak{}Tensor.\allowbreak{}exists\_\allowbreak{}rank\_\allowbreak{}power\_\allowbreak{}exponential\_\allowbreak{}bound}, \leandecl{MatrixMultiplication.\allowbreak{}Tensor.\allowbreak{}rank\_\allowbreak{}kroneckerPower\_\allowbreak{}add\_\allowbreak{}le\_\allowbreak{}sum}, \leandecl{MatrixMultiplication.\allowbreak{}exists\_\allowbreak{}pos\_\allowbreak{}exponential\_\allowbreak{}bound}, \leandecl{MatrixMultiplication.\allowbreak{}root\_\allowbreak{}limit\_\allowbreak{}le\_\allowbreak{}of\_\allowbreak{}exponential\_\allowbreak{}bound}.
\textit{Formalization note:} Checked over arbitrary commutative semirings, including zero tensors. Uniform exponential bounds include the zeroth power; the limiting argument uses positive powers.

\begin{proof}
\proves{aux:asymptotic-rank-sums}
\leanok
\uses{aux:asymptotic-rank-limit,aux:rank-operation-bounds,aux:asymptotic-rank-degeneration}
\textit{Checked proof.} For rates $a>\widetilde R(T)$ and $b>\widetilde R(S)$, bound every power rank by $C_Ta^k$ and $C_Sb^k$. Expand a power of $T+S$ over subsets of its coordinates and separate each subset from its complement. The resulting sum is bounded by $C_TC_S(a+b)^N$. Take roots and limits, then let the rates decrease. Obtain direct sums by embedding into disjoint coordinate blocks.
\end{proof}

\renewcommand{\thetheorem}{A43}
\begin{lemma}[Matrix multiplication dimension calculus]
\label{aux:matrix-multiplication-calculus}
\lean{MatrixMultiplication.Tensor.matrixMul_restricts_of_le,MatrixMultiplication.Tensor.asymptoticRank_matrixMul_mono,MatrixMultiplication.Tensor.matrixMul_kronecker_equivalent,MatrixMultiplication.Tensor.matrixMul_kroneckerPower_equivalent,MatrixMultiplication.Tensor.rank_matrixMul_pow,MatrixMultiplication.Tensor.asymptoticRank_matrixMul_pow,MatrixMultiplication.Tensor.rank_matrixMul_le,MatrixMultiplication.Tensor.matrixMul_one_equivalent_diagonal,MatrixMultiplication.Tensor.asymptoticRank_matrixMul_one}
\leanok
Decreasing matrix dimensions gives a restriction. Moreover, $\langle\ell,m,n\rangle\otimes\langle\ell',m',n'\rangle\cong\langle\ell\ell',mm',nn'\rangle$ and $\langle\ell,m,n\rangle^{\otimes k}\cong\langle\ell^k,m^k,n^k\rangle$. The elementary bound is $R(\langle\ell,m,n\rangle)\le\ell mn$, and $\widetilde R(\langle1,1,1\rangle)=1$.
\end{lemma}

\noindent\textit{Auxiliary foundation. Source:} \texttt{material/main.tex:433}.
\textit{Status:} \texttt{formalized}. Main proof path.
\textit{Checked Lean declarations:} \leandecl{MatrixMultiplication.\allowbreak{}Tensor.\allowbreak{}matrixMul\_\allowbreak{}restricts\_\allowbreak{}of\_\allowbreak{}le}, \leandecl{MatrixMultiplication.\allowbreak{}Tensor.\allowbreak{}asymptoticRank\_\allowbreak{}matrixMul\_\allowbreak{}mono}, \leandecl{MatrixMultiplication.\allowbreak{}Tensor.\allowbreak{}matrixMul\_\allowbreak{}kronecker\_\allowbreak{}equivalent}, \leandecl{MatrixMultiplication.\allowbreak{}Tensor.\allowbreak{}matrixMul\_\allowbreak{}kroneckerPower\_\allowbreak{}equivalent}, \leandecl{MatrixMultiplication.\allowbreak{}Tensor.\allowbreak{}rank\_\allowbreak{}matrixMul\_\allowbreak{}pow}, \leandecl{MatrixMultiplication.\allowbreak{}Tensor.\allowbreak{}asymptoticRank\_\allowbreak{}matrixMul\_\allowbreak{}pow}, \leandecl{MatrixMultiplication.\allowbreak{}Tensor.\allowbreak{}rank\_\allowbreak{}matrixMul\_\allowbreak{}le}, \leandecl{MatrixMultiplication.\allowbreak{}Tensor.\allowbreak{}matrixMul\_\allowbreak{}one\_\allowbreak{}equivalent\_\allowbreak{}diagonal}, \leandecl{MatrixMultiplication.\allowbreak{}Tensor.\allowbreak{}asymptoticRank\_\allowbreak{}matrixMul\_\allowbreak{}one}.
\textit{Formalization note:} Checked for all natural matrix dimensions and powers, including zero. Coordinate restrictions and equivalences use commutative semirings; scalar asymptotic-rank normalization is stated over fields. The three special-shape presentations are checked in Definition 2.4.

\begin{proof}
\proves{aux:matrix-multiplication-calculus}
\leanok
\uses{aux:explicit-tensor-definitions,aux:rank-operation-bounds,aux:asymptotic-rank-degeneration,aux:diagonal-rank-normalization}
\textit{Checked proof.} Use coordinate inclusions for restrictions, finite-product and tuple equivalences for products and powers, and the defining sum of pure tensors for the rank bound. Identify scalar multiplication with the diagonal tensor of size one.
\end{proof}

\renewcommand{\thetheorem}{A44}
\begin{lemma}[Monotone functions preserving powers]
\label{aux:monotone-power-functions}
\lean{MatrixMultiplication.log_le_logb_mul_of_monotone_pow,MatrixMultiplication.eq_rpow_logb_of_monotone_pow}
\leanok
Let $f:\mathbb N\to\mathbb R$ be monotone, satisfy $f(1)=1$, and preserve natural powers: $f(a^k)=f(a)^k$. Then $f(n)=n^{\log_2 f(2)}$ for every $n\ge1$.
\end{lemma}

\noindent\textit{Auxiliary foundation. Source:} \texttt{material/main.tex:433}.
\textit{Status:} \texttt{formalized}. Main proof path.
\textit{Checked Lean declarations:} \leandecl{MatrixMultiplication.\allowbreak{}log\_\allowbreak{}le\_\allowbreak{}logb\_\allowbreak{}mul\_\allowbreak{}of\_\allowbreak{}monotone\_\allowbreak{}pow}, \leandecl{MatrixMultiplication.\allowbreak{}eq\_\allowbreak{}rpow\_\allowbreak{}logb\_\allowbreak{}of\_\allowbreak{}monotone\_\allowbreak{}pow}.
\textit{Formalization note:} Checked general analytic helper. Monotonicity and f(1)=1 ensure all logarithms at positive integer arguments are defined. No lower bound greater than one on f(2) is required.

\begin{proof}
\proves{aux:monotone-power-functions}
\leanok
\textit{Checked proof.} Compare $a^k$ with $b^{\lceil\log_b(a^k)\rceil}$ for $a\ge1$ and $b\ge2$. Taking logarithms bounds $k\log f(a)$ by $k\log_b(a)\log f(b)$ plus a fixed rounding error. The Archimedean property removes that error. Interchange the two bases to get equality, and handle $n=1$ directly.
\end{proof}

\renewcommand{\thetheorem}{A45}
\begin{lemma}[Repeated direct-sum blocks]
\label{aux:direct-sum-copies}
\lean{MatrixMultiplication.Tensor.directSumCopies,MatrixMultiplication.Tensor.directSumCopies_apply,MatrixMultiplication.Tensor.kroneckerPower_diagonal_reindex,MatrixMultiplication.Tensor.kroneckerPower_directSumCopies_equivalent,MatrixMultiplication.Tensor.rank_kroneckerPower_directSumCopies,MatrixMultiplication.Tensor.directSumCopies_restricts,MatrixMultiplication.Tensor.asymptoticRank_le_directSumCopies,MatrixMultiplication.Tensor.asymptoticRank_directSumCopies_le}
\leanok
The ordinary direct sum of $a$ identical tensors is represented by $T^{\oplus a}=\langle a\rangle\otimes T$, with coefficients zero unless the three block labels agree. Its powers satisfy $(T^{\oplus a})^{\otimes n}\cong(T^{\otimes n})^{\oplus a^n}$. Over a field, $\widetilde R(T^{\oplus a})\le a\widetilde R(T)$; when $a\ge1$, selecting one block gives the reverse lower bound $\widetilde R(T)\le\widetilde R(T^{\oplus a})$.
\end{lemma}

\noindent\textit{Auxiliary foundation. Source:} \texttt{material/main.tex:492}.
\textit{Status:} \texttt{formalized}. Main proof path.
\textit{Checked Lean declarations:} \leandecl{MatrixMultiplication.\allowbreak{}Tensor.\allowbreak{}directSumCopies}, \leandecl{MatrixMultiplication.\allowbreak{}Tensor.\allowbreak{}directSumCopies\_\allowbreak{}apply}, \leandecl{MatrixMultiplication.\allowbreak{}Tensor.\allowbreak{}kroneckerPower\_\allowbreak{}diagonal\_\allowbreak{}reindex}, \leandecl{MatrixMultiplication.\allowbreak{}Tensor.\allowbreak{}kroneckerPower\_\allowbreak{}directSumCopies\_\allowbreak{}equivalent}, \leandecl{MatrixMultiplication.\allowbreak{}Tensor.\allowbreak{}rank\_\allowbreak{}kroneckerPower\_\allowbreak{}directSumCopies}, \leandecl{MatrixMultiplication.\allowbreak{}Tensor.\allowbreak{}directSumCopies\_\allowbreak{}restricts}, \leandecl{MatrixMultiplication.\allowbreak{}Tensor.\allowbreak{}asymptoticRank\_\allowbreak{}le\_\allowbreak{}directSumCopies}, \leandecl{MatrixMultiplication.\allowbreak{}Tensor.\allowbreak{}asymptoticRank\_\allowbreak{}directSumCopies\_\allowbreak{}le}.
\textit{Formalization note:} Checked including zero copies and zeroth powers for the power identities and upper bound. Block projection requires a positive number of copies. The coefficient lemma verifies that the diagonal-product representation is exactly the ordinary repeated direct sum.

\begin{proof}
\proves{aux:direct-sum-copies}
\leanok
\uses{aux:explicit-tensor-definitions,aux:asymptotic-rank-powers-products,aux:diagonal-rank-normalization,aux:asymptotic-rank-degeneration}
\textit{Checked proof.} Encode words of block labels as a single finite label and split tuples of block-and-basis indices. Projection to one block is a restriction. The upper bound follows from product submultiplicativity and the normalization $\widetilde R(\langle a\rangle)=a$.
\end{proof}

\renewcommand{\thetheorem}{A46}
\begin{lemma}[Matrix multiplication capacity comparison]
\label{aux:matrix-multiplication-capacity}
\lean{MatrixMultiplication.Tensor.diagonal_restricts_of_le,MatrixMultiplication.Tensor.diagonal_restricts_of_rank_le,MatrixMultiplication.Tensor.asymptoticRank_matrixMul_kronecker_eq,MatrixMultiplication.Tensor.asymptoticRank_matrixMul_power_kronecker_eq,MatrixMultiplication.Tensor.matrixMul_capacity_inequality}
\leanok
Fix $b\ge1$ and put $q=\widetilde R(\langle2,2,2\rangle)$ and $s=\widetilde R(\langle b,b,b\rangle^{\oplus a})$. For natural $n,m$, if $R(\langle2,2,2\rangle^{\otimes n})\le a^m$, then $q^n(b^\omega)^m\le s^m$.
\end{lemma}

\noindent\textit{Auxiliary foundation. Source:} \texttt{material/main.tex:492}.
\textit{Status:} \texttt{formalized}. Main proof path.
\textit{Checked Lean declarations:} \leandecl{MatrixMultiplication.\allowbreak{}Tensor.\allowbreak{}diagonal\_\allowbreak{}restricts\_\allowbreak{}of\_\allowbreak{}le}, \leandecl{MatrixMultiplication.\allowbreak{}Tensor.\allowbreak{}diagonal\_\allowbreak{}restricts\_\allowbreak{}of\_\allowbreak{}rank\_\allowbreak{}le}, \leandecl{MatrixMultiplication.\allowbreak{}Tensor.\allowbreak{}asymptoticRank\_\allowbreak{}matrixMul\_\allowbreak{}kronecker\_\allowbreak{}eq}, \leandecl{MatrixMultiplication.\allowbreak{}Tensor.\allowbreak{}asymptoticRank\_\allowbreak{}matrixMul\_\allowbreak{}power\_\allowbreak{}kronecker\_\allowbreak{}eq}, \leandecl{MatrixMultiplication.\allowbreak{}Tensor.\allowbreak{}matrixMul\_\allowbreak{}capacity\_\allowbreak{}inequality}.
\textit{Formalization note:} Checked over every field. The exponents n,m and copy count a may be zero; only the matrix dimension b must be positive. The auxiliary diagonal restriction and rank-attainment lemmas work over commutative semirings.

\begin{proof}
\proves{aux:matrix-multiplication-capacity}
\leanok
\uses{aux:ordinary-rank-definition,aux:restriction-operations,aux:direct-sum-copies,aux:matrix-multiplication-calculus,eq:omega-log-n,aux:asymptotic-rank-powers-products}
\textit{Checked proof.} A rank bound gives a restriction from the diagonal tensor of size $a^m$. Tensor it with $\langle b,b,b\rangle^{\otimes m}$ and identify the source as the $m$th power of the repeated block tensor. Matrix dimensions multiply under tensor products, so the square-size identity computes the target asymptotic rank exactly. Apply restriction monotonicity and power multiplicativity.
\end{proof}

\renewcommand{\thetheorem}{A47}
\begin{lemma}[Integer capacities and root growth]
\label{aux:integer-capacity-growth}
\lean{MatrixMultiplication.le_of_root_limit_capacity}
\leanok
Let $r_n\ge1$ be natural numbers with $r_n^{1/n}\to q>1$, let $a\ge2$ be an integer, and let $p\ge1$. If $r_n\le a^m$ always implies $q^n\le p^m$, then $a\le p$.
\end{lemma}

\noindent\textit{Auxiliary foundation. Source:} \texttt{material/main.tex:492}.
\textit{Status:} \texttt{formalized}. Main proof path.
\textit{Checked Lean declarations:} \leandecl{MatrixMultiplication.\allowbreak{}le\_\allowbreak{}of\_\allowbreak{}root\_\allowbreak{}limit\_\allowbreak{}capacity}.
\textit{Formalization note:} Checked analytic helper with all logarithm and division hypotheses explicit. The ceiling is implemented as Nat.clog, so every exponent is an integer. The limit ignores n=0.

\begin{proof}
\proves{aux:integer-capacity-growth}
\leanok
\textit{Checked proof.} Take $m=\lceil\log_a r_n\rceil$. Taking logarithms of the capacity inequality and using the ceiling bound gives $\log q\le(\log(r_n^{1/n})/\log a+1/n)\log p$ for positive $n$. Pass to the limit and cancel the positive logarithms of $a$ and $q$.
\end{proof}

\renewcommand{\thetheorem}{A48}
\begin{lemma}[Permuting the three tensor factors]
\label{aux:factor-permutations}
\lean{MatrixMultiplication.Tensor.rotate,MatrixMultiplication.Tensor.swapFirstTwo,MatrixMultiplication.Tensor.rank_rotate,MatrixMultiplication.Tensor.rank_swapFirstTwo,MatrixMultiplication.Tensor.asymptoticRank_rotate,MatrixMultiplication.Tensor.asymptoticRank_swapFirstTwo,MatrixMultiplication.Tensor.FactorPermutation,MatrixMultiplication.Tensor.FactorPermutation.equivPerm,MatrixMultiplication.Tensor.permuteS3}
\leanok
The six permutations of three tensor factors preserve ordinary rank and asymptotic rank. A three-cycle and a transposition generate the six explicit orders, identified bijectively with $\mathfrak S_3$.
\end{lemma}

\noindent\textit{Auxiliary foundation. Source:} \texttt{material/main.tex:662}.
\textit{Status:} \texttt{formalized}. Main proof path.
\textit{Checked Lean declarations:} \leandecl{MatrixMultiplication.\allowbreak{}Tensor.\allowbreak{}rotate}, \leandecl{MatrixMultiplication.\allowbreak{}Tensor.\allowbreak{}swapFirstTwo}, \leandecl{MatrixMultiplication.\allowbreak{}Tensor.\allowbreak{}rank\_\allowbreak{}rotate}, \leandecl{MatrixMultiplication.\allowbreak{}Tensor.\allowbreak{}rank\_\allowbreak{}swapFirstTwo}, \leandecl{MatrixMultiplication.\allowbreak{}Tensor.\allowbreak{}asymptoticRank\_\allowbreak{}rotate}, \leandecl{MatrixMultiplication.\allowbreak{}Tensor.\allowbreak{}asymptoticRank\_\allowbreak{}swapFirstTwo}, \leandecl{MatrixMultiplication.\allowbreak{}Tensor.\allowbreak{}FactorPermutation}, \leandecl{MatrixMultiplication.\allowbreak{}Tensor.\allowbreak{}FactorPermutation.\allowbreak{}equivPerm}, \leandecl{MatrixMultiplication.\allowbreak{}Tensor.\allowbreak{}permuteS3}.
\textit{Formalization note:} Checked over commutative semirings. The S3 interface permutes the coordinate types as well as the coefficients; its six-case representation is proved bijective to Equiv.Perm (Fin 3).

\begin{proof}
\proves{aux:factor-permutations}
\leanok
\uses{aux:rank-operation-bounds,aux:asymptotic-rank-definition}
\textit{Checked proof.} Permute each pure tensor in an attained rank decomposition to obtain one rank inequality; apply the inverse permutation for the other. Tensor powers commute with these permutations, so the asymptotic-rank equality follows from its definition. Enumerate the six orders and identify them with permutations of three slots.
\end{proof}

\renewcommand{\thetheorem}{A49}
\begin{lemma}[Asymptotic rank of symmetrization]
\label{aux:symmetrization-rank-bound}
\lean{MatrixMultiplication.Tensor.asymptoticRank_cyclicSymmetrize_le,MatrixMultiplication.Tensor.asymptoticRank_symmetrize_le}
\leanok
For every tensor over a field, $\widetilde R(T^{\otimes\mathfrak S_3})\le\widetilde R(T)^6$. The cyclic product of three rotations has asymptotic rank at most $\widetilde R(T)^3$.
\end{lemma}

\noindent\textit{Auxiliary foundation. Source:} \texttt{material/main.tex:662}.
\textit{Status:} \texttt{formalized}. Main proof path.
\textit{Checked Lean declarations:} \leandecl{MatrixMultiplication.\allowbreak{}Tensor.\allowbreak{}asymptoticRank\_\allowbreak{}cyclicSymmetrize\_\allowbreak{}le}, \leandecl{MatrixMultiplication.\allowbreak{}Tensor.\allowbreak{}asymptoticRank\_\allowbreak{}symmetrize\_\allowbreak{}le}.
\textit{Formalization note:} Checked over every commutative semiring, including zero tensors. This estimate does not use any fractional-value identity.

\begin{proof}
\proves{aux:symmetrization-rank-bound}
\leanok
\uses{def:symmetrization,aux:factor-permutations,aux:asymptotic-rank-powers-products}
\textit{Checked proof.} Apply product submultiplicativity to the three even and three odd permutations, using permutation invariance of asymptotic rank.
\end{proof}

\renewcommand{\thetheorem}{A50}
\begin{lemma}[Matrices as tensors with a unit factor]
\label{aux:matrix-tensor-rank}
\lean{MatrixMultiplication.Tensor.matrixTensor,MatrixMultiplication.Tensor.exists_matrix_rank_decomposition,MatrixMultiplication.Tensor.rank_matrixTensor,MatrixMultiplication.Tensor.asymptoticRank_matrixTensor}
\leanok
Regard a matrix $A$ as a three-tensor with one unit factor. Its ordinary tensor rank and asymptotic rank both equal the matrix rank of $A$.
\end{lemma}

\noindent\textit{Auxiliary foundation. Source:} \texttt{material/main.tex:393}.
\textit{Status:} \texttt{formalized}. Main proof path.
\textit{Checked Lean declarations:} \leandecl{MatrixMultiplication.\allowbreak{}Tensor.\allowbreak{}matrixTensor}, \leandecl{MatrixMultiplication.\allowbreak{}Tensor.\allowbreak{}exists\_\allowbreak{}matrix\_\allowbreak{}rank\_\allowbreak{}decomposition}, \leandecl{MatrixMultiplication.\allowbreak{}Tensor.\allowbreak{}rank\_\allowbreak{}matrixTensor}, \leandecl{MatrixMultiplication.\allowbreak{}Tensor.\allowbreak{}asymptoticRank\_\allowbreak{}matrixTensor}.
\textit{Formalization note:} Checked over every field. This supplies an algebraic route to flattening monotonicity without a separate minors or algebraic-closure argument.

\begin{proof}
\proves{aux:matrix-tensor-rank}
\leanok
\uses{aux:rank-operation-bounds,lem:flattening-lb}
\textit{Checked proof.} Choose a basis of the image to decompose the matrix into exactly its rank many outer products. This gives the tensor-rank upper bound; the first flattening recovers the original matrix and gives the matching asymptotic lower bound.
\end{proof}

\renewcommand{\thetheorem}{A51}
\begin{lemma}[Flattening monotonicity under polynomial degeneration]
\label{aux:flattening-degeneration}
\lean{MatrixMultiplication.Tensor.DegenerationWitness.flatteningFirst,MatrixMultiplication.Tensor.DegenerationWitness.flatteningSecond,MatrixMultiplication.Tensor.DegenerationWitness.flatteningThird,MatrixMultiplication.Tensor.flatteningRank_mono_degeneration,MatrixMultiplication.Tensor.flatteningRank_mono,MatrixMultiplication.Tensor.flatteningRank_eq_of_equivalent}
\leanok
For the polynomial degeneration relation, $T\trianglerighteq S$ implies $R_i(S)\le R_i(T)$ for each of the three flattenings. Flattening rank is also monotone under restriction and invariant under tensor equivalence.
\end{lemma}

\noindent\textit{Auxiliary foundation. Source:} \texttt{material/main.tex:393}.
\textit{Status:} \texttt{formalized}. Main proof path.
\textit{Checked Lean declarations:} \leandecl{MatrixMultiplication.\allowbreak{}Tensor.\allowbreak{}DegenerationWitness.\allowbreak{}flatteningFirst}, \leandecl{MatrixMultiplication.\allowbreak{}Tensor.\allowbreak{}DegenerationWitness.\allowbreak{}flatteningSecond}, \leandecl{MatrixMultiplication.\allowbreak{}Tensor.\allowbreak{}DegenerationWitness.\allowbreak{}flatteningThird}, \leandecl{MatrixMultiplication.\allowbreak{}Tensor.\allowbreak{}flatteningRank\_\allowbreak{}mono\_\allowbreak{}degeneration}, \leandecl{MatrixMultiplication.\allowbreak{}Tensor.\allowbreak{}flatteningRank\_\allowbreak{}mono}, \leandecl{MatrixMultiplication.\allowbreak{}Tensor.\allowbreak{}flatteningRank\_\allowbreak{}eq\_\allowbreak{}of\_\allowbreak{}equivalent}.
\textit{Formalization note:} Checked over every field. The rational/polynomial degeneration equivalence transfers this result to the manuscript relation.

\begin{proof}
\proves{aux:flattening-degeneration}
\leanok
\uses{aux:matrix-tensor-rank,aux:polynomial-degeneration-definition,aux:asymptotic-rank-degeneration}
\textit{Checked proof.} Group the two complementary factor maps by their matrix Kronecker product, transporting the original polynomial witness to the corresponding matrix tensors without changing its order. Apply asymptotic-rank monotonicity and the equality of matrix-tensor asymptotic rank with matrix rank.
\end{proof}

\renewcommand{\thetheorem}{A52}
\begin{lemma}[Flattening ranks of repeated blocks]
\label{aux:flattening-repeated-blocks}
\lean{MatrixMultiplication.Tensor.flatteningRank_directSumCopies,MatrixMultiplication.Tensor.flatteningRank_directSumCopies_matrixMul_zero}
\leanok
For each flattening, $R_i(T^{\oplus a})=aR_i(T)$. In particular, the first flattening of $\langle b,b,b\rangle^{\oplus a}$ has rank $ab^2$ when $b\ge1$.
\end{lemma}

\noindent\textit{Auxiliary foundation. Source:} \texttt{material/main.tex:791}.
\textit{Status:} \texttt{formalized}. Main proof path.
\textit{Checked Lean declarations:} \leandecl{MatrixMultiplication.\allowbreak{}Tensor.\allowbreak{}flatteningRank\_\allowbreak{}directSumCopies}, \leandecl{MatrixMultiplication.\allowbreak{}Tensor.\allowbreak{}flatteningRank\_\allowbreak{}directSumCopies\_\allowbreak{}matrixMul\_\allowbreak{}zero}.
\textit{Formalization note:} Checked over every field, including zero copies. Only the displayed square-matrix specialization requires a positive matrix dimension.

\begin{proof}
\proves{aux:flattening-repeated-blocks}
\leanok
\uses{aux:direct-sum-copies,aux:flattening-operations,aux:diagonal-rank-normalization,aux:matrix-multiplication-flattening}
\textit{Checked proof.} Use the diagonal-product description of repeated blocks, flattening-rank multiplicativity, diagonal normalization, and the matrix-multiplication flattening formula.
\end{proof}

\renewcommand{\thetheorem}{A53}
\begin{lemma}[Nonzero tensors supply the unit]
\label{aux:nonzero-unit-restriction}
\lean{MatrixMultiplication.Tensor.exists_coeff_ne_zero,MatrixMultiplication.Tensor.restricts_diagonal_one_of_ne_zero,MatrixMultiplication.Tensor.kronecker_ne_zero,MatrixMultiplication.Tensor.kroneckerPower_ne_zero,MatrixMultiplication.Tensor.diagonal_ne_zero,MatrixMultiplication.Tensor.one_le_asymptoticRank_of_ne_zero}
\leanok
Every nonzero tensor over a field restricts to $\langle1\rangle$, hence has asymptotic rank at least one. Products and powers of nonzero tensors are nonzero, including the zeroth power.
\end{lemma}

\noindent\textit{Auxiliary foundation. Source:} \texttt{material/main.tex:682}.
\textit{Status:} \texttt{formalized}. Main proof path.
\textit{Checked Lean declarations:} \leandecl{MatrixMultiplication.\allowbreak{}Tensor.\allowbreak{}exists\_\allowbreak{}coeff\_\allowbreak{}ne\_\allowbreak{}zero}, \leandecl{MatrixMultiplication.\allowbreak{}Tensor.\allowbreak{}restricts\_\allowbreak{}diagonal\_\allowbreak{}one\_\allowbreak{}of\_\allowbreak{}ne\_\allowbreak{}zero}, \leandecl{MatrixMultiplication.\allowbreak{}Tensor.\allowbreak{}kronecker\_\allowbreak{}ne\_\allowbreak{}zero}, \leandecl{MatrixMultiplication.\allowbreak{}Tensor.\allowbreak{}kroneckerPower\_\allowbreak{}ne\_\allowbreak{}zero}, \leandecl{MatrixMultiplication.\allowbreak{}Tensor.\allowbreak{}diagonal\_\allowbreak{}ne\_\allowbreak{}zero}, \leandecl{MatrixMultiplication.\allowbreak{}Tensor.\allowbreak{}one\_\allowbreak{}le\_\allowbreak{}asymptoticRank\_\allowbreak{}of\_\allowbreak{}ne\_\allowbreak{}zero}.
\textit{Formalization note:} Checked over every field. No nonemptiness assumption on coordinate types is imposed; it follows from a nonzero coefficient when required.

\begin{proof}
\proves{aux:nonzero-unit-restriction}
\leanok
\uses{aux:coefficient-tensors,aux:diagonal-rank-normalization,aux:asymptotic-rank-degeneration}
\textit{Checked proof.} Select one nonzero coefficient and rescale by its inverse. Products use products of nonzero coefficients; powers use constant index tuples. Apply asymptotic-rank monotonicity for the lower bound.
\end{proof}

\renewcommand{\thetheorem}{A54}
\begin{definition}[Implemented fractional value and positive witnesses]
\label{aux:fractional-value-definition}
\lean{MatrixMultiplication.Tensor.ValueWitness,MatrixMultiplication.Tensor.ValueWitness.score,MatrixMultiplication.Tensor.ValueWitness.score_pos,MatrixMultiplication.Tensor.fractionalValue,MatrixMultiplication.Tensor.fractionalValue_zero,MatrixMultiplication.Tensor.ValueWitness.score_le_fractionalValue,MatrixMultiplication.Tensor.ValueWitness.tensor_ne_zero,MatrixMultiplication.Tensor.valueWitness_zero_isEmpty,MatrixMultiplication.Tensor.fractionalValue_eq_iSup}
\leanok
A witness consists of positive integers $a,b,c,d,N$ and $\langle d,d,d\rangle^{\oplus c}\otimes(T^{\otimes\mathfrak S_3})^{\otimes N}\trianglerighteq\langle b,b,b\rangle^{\oplus a}$. Its positive score is $(ab^{3\tau}/(cd^{3\tau}))^{1/(6N)}$. Define the value as the supremum of all witness scores in $[0,+\infty]$. Every witness forces $T\ne0$; therefore the zero tensor has value zero as the empty supremum.
\end{definition}

\noindent\textit{Auxiliary foundation. Source:} \texttt{material/main.tex:682}.
\textit{Status:} \texttt{formalized}. Main proof path.
\textit{Checked Lean declarations:} \leandecl{MatrixMultiplication.\allowbreak{}Tensor.\allowbreak{}ValueWitness}, \leandecl{MatrixMultiplication.\allowbreak{}Tensor.\allowbreak{}ValueWitness.\allowbreak{}score}, \leandecl{MatrixMultiplication.\allowbreak{}Tensor.\allowbreak{}ValueWitness.\allowbreak{}score\_\allowbreak{}pos}, \leandecl{MatrixMultiplication.\allowbreak{}Tensor.\allowbreak{}fractionalValue}, \leandecl{MatrixMultiplication.\allowbreak{}Tensor.\allowbreak{}fractionalValue\_\allowbreak{}zero}, \leandecl{MatrixMultiplication.\allowbreak{}Tensor.\allowbreak{}ValueWitness.\allowbreak{}score\_\allowbreak{}le\_\allowbreak{}fractionalValue}, \leandecl{MatrixMultiplication.\allowbreak{}Tensor.\allowbreak{}ValueWitness.\allowbreak{}tensor\_\allowbreak{}ne\_\allowbreak{}zero}, \leandecl{MatrixMultiplication.\allowbreak{}Tensor.\allowbreak{}valueWitness\_\allowbreak{}zero\_\allowbreak{}isEmpty}, \leandecl{MatrixMultiplication.\allowbreak{}Tensor.\allowbreak{}fractionalValue\_\allowbreak{}eq\_\allowbreak{}iSup}.
\textit{Formalization note:} The supremum formula holds definitionally for every tensor. The zero witness type is proved empty, giving value zero by the empty-supremum law. Scores and values extend to all real parameters; the paper uses positive parameters. The rational-value-equivalence node identifies the same supremum with rational-function witnesses.

\begin{proof}
\proves{aux:fractional-value-definition}
\leanok
\uses{def:symmetrization,aux:direct-sum-copies,aux:explicit-tensor-definitions,aux:polynomial-degeneration-definition}
\textit{Checked proof.} Store the positive sizes and degeneration as a witness structure. A zero source cannot degenerate to the nonzero target of a positive witness, so the witness type for the zero tensor is empty. Define value by the indexed supremum of the extended nonnegative real scores; the empty-supremum identity proves the zero case. Each witness bounds the supremum without a separate nonzero hypothesis.
\end{proof}

\renewcommand{\thetheorem}{A55}
\begin{lemma}[Bounding intermediate witness scores]
\label{aux:witness-endpoint-interpolation}
\lean{MatrixMultiplication.witness_ratio_le_of_endpoint_bounds,MatrixMultiplication.witness_root_le_of_bound,MatrixMultiplication.witness_root_le_of_endpoint_bounds}
\leanok
For positive real $a,b,c,d$, a natural $k\ge1$, and $R\ge0$, suppose $ab^2\le cd^2R^k$ and $ab^w\le cd^wR^k$. Then $(ab^t/(cd^t))^{1/k}\le R$ for every $2\le t\le w$.
\end{lemma}

\noindent\textit{Auxiliary foundation. Source:} \texttt{material/main.tex:791}.
\textit{Status:} \texttt{formalized}. Main proof path.
\textit{Checked Lean declarations:} \leandecl{MatrixMultiplication.\allowbreak{}witness\_\allowbreak{}ratio\_\allowbreak{}le\_\allowbreak{}of\_\allowbreak{}endpoint\_\allowbreak{}bounds}, \leandecl{MatrixMultiplication.\allowbreak{}witness\_\allowbreak{}root\_\allowbreak{}le\_\allowbreak{}of\_\allowbreak{}bound}, \leandecl{MatrixMultiplication.\allowbreak{}witness\_\allowbreak{}root\_\allowbreak{}le\_\allowbreak{}of\_\allowbreak{}endpoint\_\allowbreak{}bounds}.
\textit{Formalization note:} Checked without logarithms or division by w-2. The degenerate interval w=2 and zero right-hand rate are included; no value calculus is assumed.

\begin{proof}
\proves{aux:witness-endpoint-interpolation}
\leanok
\textit{Checked proof.} The ratio equals $(a/c)(b/d)^t$. If $b\le d$, the lower exponent controls the ratio; if $d\le b$, the upper exponent does. Apply monotonicity of the positive root to the appropriate bound.
\end{proof}

\renewcommand{\thetheorem}{A56}
\begin{lemma}[A score-one witness for nonzero tensors]
\label{aux:unit-value-witness}
\lean{MatrixMultiplication.Tensor.symmetrize_ne_zero,MatrixMultiplication.Tensor.symmetrize_restricts_diagonal_one,MatrixMultiplication.Tensor.unitValueWitness,MatrixMultiplication.Tensor.unitValueWitness_score,MatrixMultiplication.Tensor.valueWitness_nonempty,MatrixMultiplication.Tensor.one_le_fractionalValue}
\leanok
Every nonzero tensor has a fractional-value witness with $a=b=c=d=N=1$, whose score is one for every real $\tau$. Thus its witness set is nonempty and its polynomial-witness value is at least one.
\end{lemma}

\noindent\textit{Auxiliary foundation. Source:} \texttt{material/main.tex:682}.
\textit{Status:} \texttt{formalized}. Main proof path.
\textit{Checked Lean declarations:} \leandecl{MatrixMultiplication.\allowbreak{}Tensor.\allowbreak{}symmetrize\_\allowbreak{}ne\_\allowbreak{}zero}, \leandecl{MatrixMultiplication.\allowbreak{}Tensor.\allowbreak{}symmetrize\_\allowbreak{}restricts\_\allowbreak{}diagonal\_\allowbreak{}one}, \leandecl{MatrixMultiplication.\allowbreak{}Tensor.\allowbreak{}unitValueWitness}, \leandecl{MatrixMultiplication.\allowbreak{}Tensor.\allowbreak{}unitValueWitness\_\allowbreak{}score}, \leandecl{MatrixMultiplication.\allowbreak{}Tensor.\allowbreak{}valueWitness\_\allowbreak{}nonempty}, \leandecl{MatrixMultiplication.\allowbreak{}Tensor.\allowbreak{}one\_\allowbreak{}le\_\allowbreak{}fractionalValue}.
\textit{Formalization note:} Checked over every field, with no admissible-interval assumption. Nonemptiness is established before any manipulation of the supremum.

\begin{proof}
\proves{aux:unit-value-witness}
\leanok
\uses{aux:fractional-value-definition,aux:nonzero-unit-restriction,lem:sym-properties,aux:matrix-multiplication-capacity}
\textit{Checked proof.} The six-factor symmetrization is nonzero. Tensoring it with the scalar matrix-multiplication block remains nonzero, so it restricts to the unit diagonal and then to the target scalar block. This restriction is a degeneration and gives the all-one witness.
\end{proof}

\renewcommand{\thetheorem}{A57}
\begin{lemma}[Rational specialization and polynomial witnesses]
\label{aux:rational-degeneration-equivalence}
\lean{MatrixMultiplication.Tensor.liftRational,MatrixMultiplication.Tensor.RegularAtZeroValue,MatrixMultiplication.Tensor.RationalDegenerates,MatrixMultiplication.Tensor.Degenerates.rationalDegenerates,MatrixMultiplication.Tensor.RationalDegenerates.degenerates,MatrixMultiplication.Tensor.rationalDegenerates_iff_degenerates}
\leanok
For finite-dimensional tensors over every field, rational-function specialization is equivalent to polynomial degeneration witnesses. The rational factor maps may have poles at zero; their transformed tensor must be regular there and specialize to the target.
\end{lemma}

\noindent\textit{Auxiliary foundation. Source:} \texttt{material/main.tex:312}.
\textit{Status:} \texttt{formalized}. Main proof path.
\textit{Checked Lean declarations:} \leandecl{MatrixMultiplication.\allowbreak{}Tensor.\allowbreak{}liftRational}, \leandecl{MatrixMultiplication.\allowbreak{}Tensor.\allowbreak{}RegularAtZeroValue}, \leandecl{MatrixMultiplication.\allowbreak{}Tensor.\allowbreak{}RationalDegenerates}, \leandecl{MatrixMultiplication.\allowbreak{}Tensor.\allowbreak{}Degenerates.\allowbreak{}rationalDegenerates}, \leandecl{MatrixMultiplication.\allowbreak{}Tensor.\allowbreak{}RationalDegenerates.\allowbreak{}degenerates}, \leandecl{MatrixMultiplication.\allowbreak{}Tensor.\allowbreak{}rationalDegenerates\_\allowbreak{}iff\_\allowbreak{}degenerates}.
\textit{Formalization note:} Checked without assumptions on characteristic, field cardinality, or nonzero tensors. Regularity uses the reduced rational-function denominator, so removable poles are handled correctly.

\begin{proof}
\proves{aux:rational-degeneration-equivalence}
\leanok
\uses{aux:polynomial-degeneration-definition,aux:polynomial-degeneration-operations}
\textit{Checked proof.} In one direction divide a factor map by $X^d$, obtaining the regular output $S+XE$. Conversely clear the finite factor denominators and a separate common denominator of the regular output. Extract the common power of $X$ from the factor denominator product, normalize the remaining constant term, and use polynomial division by $X$ for the remainder.
\end{proof}

\renewcommand{\thetheorem}{A58}
\begin{lemma}[Witness transport along symmetrized degeneration]
\label{aux:symmetrized-value-monotonicity}
\lean{MatrixMultiplication.Tensor.Degenerates.rotate,MatrixMultiplication.Tensor.Degenerates.swapFirstTwo,MatrixMultiplication.Tensor.Degenerates.symmetrize,MatrixMultiplication.Tensor.ValueWitness.ofSymmetrizeDegenerates,MatrixMultiplication.Tensor.fractionalValue_mono_of_symmetrize_degenerates,MatrixMultiplication.Tensor.fractionalValue_mono,MatrixMultiplication.Tensor.fractionalValue_eq_of_equivalent,MatrixMultiplication.Tensor.fractionalValue_eq_of_symmetrize_equivalent,MatrixMultiplication.Tensor.fractionalValue_permute,MatrixMultiplication.Tensor.fractionalValue_permuteS3}
\leanok
If $T^{\otimes\mathfrak S_3}\trianglerighteq S^{\otimes\mathfrak S_3}$, every witness for $S$ is a witness for $T$ with the same parameters and score. Hence $V_\tau(S)\le V_\tau(T)$ for every real $\tau$. Value is invariant under all six factor permutations. A zero tensor has no positive witness, so the defining supremum also handles zero without a separate branch.
\end{lemma}

\noindent\textit{Auxiliary foundation. Source:} \texttt{material/main.tex:724}.
\textit{Status:} \texttt{formalized}. Main proof path.
\textit{Checked Lean declarations:} \leandecl{MatrixMultiplication.\allowbreak{}Tensor.\allowbreak{}Degenerates.\allowbreak{}rotate}, \leandecl{MatrixMultiplication.\allowbreak{}Tensor.\allowbreak{}Degenerates.\allowbreak{}swapFirstTwo}, \leandecl{MatrixMultiplication.\allowbreak{}Tensor.\allowbreak{}Degenerates.\allowbreak{}symmetrize}, \leandecl{MatrixMultiplication.\allowbreak{}Tensor.\allowbreak{}ValueWitness.\allowbreak{}ofSymmetrizeDegenerates}, \leandecl{MatrixMultiplication.\allowbreak{}Tensor.\allowbreak{}fractionalValue\_\allowbreak{}mono\_\allowbreak{}of\_\allowbreak{}symmetrize\_\allowbreak{}degenerates}, \leandecl{MatrixMultiplication.\allowbreak{}Tensor.\allowbreak{}fractionalValue\_\allowbreak{}mono}, \leandecl{MatrixMultiplication.\allowbreak{}Tensor.\allowbreak{}fractionalValue\_\allowbreak{}eq\_\allowbreak{}of\_\allowbreak{}equivalent}, \leandecl{MatrixMultiplication.\allowbreak{}Tensor.\allowbreak{}fractionalValue\_\allowbreak{}eq\_\allowbreak{}of\_\allowbreak{}symmetrize\_\allowbreak{}equivalent}, \leandecl{MatrixMultiplication.\allowbreak{}Tensor.\allowbreak{}fractionalValue\_\allowbreak{}permute}, \leandecl{MatrixMultiplication.\allowbreak{}Tensor.\allowbreak{}fractionalValue\_\allowbreak{}permuteS3}.
\textit{Formalization note:} Witness transport under degeneration of symmetrizations proves value monotonicity and permutation invariance. This supports the updated direct argument in the manuscript. The value may be infinite, and the parameter need not lie in the admissible interval.

\begin{proof}
\proves{aux:symmetrized-value-monotonicity}
\leanok
\uses{aux:fractional-value-definition,aux:polynomial-degeneration-operations,lem:sym-properties}
\textit{Checked proof.} Permute polynomial factor maps to transport degeneration through symmetrization. Compose a powered symmetrized degeneration with each witness and take the extended-real supremum. Apply the strengthened monotonicity in both directions to the symmetrization equivalences for factor permutations.
\end{proof}

\renewcommand{\thetheorem}{A59}
\begin{lemma}[Equality with the rational-function value supremum]
\label{aux:rational-value-equivalence}
\lean{MatrixMultiplication.Tensor.RationalValueWitness,MatrixMultiplication.Tensor.RationalValueWitness.toPolynomial,MatrixMultiplication.Tensor.ValueWitness.toRational,MatrixMultiplication.Tensor.valueWitnessEquivRational,MatrixMultiplication.Tensor.rationalFractionalValue,MatrixMultiplication.Tensor.fractionalValue_eq_rationalFractionalValue,MatrixMultiplication.Tensor.rationalFractionalValue_zero,MatrixMultiplication.Tensor.fractionalValue_mono_rational}
\leanok
The polynomial and rational-function value witnesses are equivalent with identical integer parameters and scores. Their extended-real suprema are equal for every tensor and every real parameter, including the zero tensor.
\end{lemma}

\noindent\textit{Auxiliary foundation. Source:} \texttt{material/main.tex:682}.
\textit{Status:} \texttt{formalized}. Main proof path.
\textit{Checked Lean declarations:} \leandecl{MatrixMultiplication.\allowbreak{}Tensor.\allowbreak{}RationalValueWitness}, \leandecl{MatrixMultiplication.\allowbreak{}Tensor.\allowbreak{}RationalValueWitness.\allowbreak{}toPolynomial}, \leandecl{MatrixMultiplication.\allowbreak{}Tensor.\allowbreak{}ValueWitness.\allowbreak{}toRational}, \leandecl{MatrixMultiplication.\allowbreak{}Tensor.\allowbreak{}valueWitnessEquivRational}, \leandecl{MatrixMultiplication.\allowbreak{}Tensor.\allowbreak{}rationalFractionalValue}, \leandecl{MatrixMultiplication.\allowbreak{}Tensor.\allowbreak{}fractionalValue\_\allowbreak{}eq\_\allowbreak{}rationalFractionalValue}, \leandecl{MatrixMultiplication.\allowbreak{}Tensor.\allowbreak{}rationalFractionalValue\_\allowbreak{}zero}, \leandecl{MatrixMultiplication.\allowbreak{}Tensor.\allowbreak{}fractionalValue\_\allowbreak{}mono\_\allowbreak{}rational}.
\textit{Formalization note:} This closes the correspondence obligation for Definition 4.3, rather than merely identifying the two degeneration relations.

\begin{proof}
\proves{aux:rational-value-equivalence}
\leanok
\uses{aux:rational-degeneration-equivalence,aux:fractional-value-definition,aux:symmetrized-value-monotonicity}
\textit{Checked proof.} Transfer the degeneration field of each witness using the rational/polynomial equivalence. The score expression is unchanged. Compare the two indexed suprema using these transports in both directions; the zero case follows from the absence of a positive witness.
\end{proof}

\renewcommand{\thetheorem}{A60}
\begin{lemma}[Witness products and exact power laws]
\label{aux:value-products-powers}
\lean{MatrixMultiplication.Tensor.Equivalent.kroneckerPower,MatrixMultiplication.Tensor.kroneckerPower_kronecker_equivalent,MatrixMultiplication.Tensor.kroneckerPower_mul_equivalent,MatrixMultiplication.Tensor.symmetrize_kroneckerPower_equivalent,MatrixMultiplication.Tensor.symmetrize_kroneckerPower_iterated_equivalent,MatrixMultiplication.Tensor.matrixMul_copies_kronecker_equivalent,MatrixMultiplication.Tensor.matrixMul_copies_kroneckerPower_equivalent,MatrixMultiplication.Tensor.ValueWitness.product,MatrixMultiplication.Tensor.ValueWitness.tensorPower,MatrixMultiplication.Tensor.ValueWitness.ofPower,MatrixMultiplication.Tensor.ValueWitness.log_score,MatrixMultiplication.Tensor.ValueWitness.product_score,MatrixMultiplication.Tensor.ValueWitness.tensorPower_score,MatrixMultiplication.Tensor.ValueWitness.ofPower_score_pow,MatrixMultiplication.Tensor.fractionalValue_kronecker_ge,MatrixMultiplication.Tensor.fractionalValue_pow_of_pos,MatrixMultiplication.Tensor.fractionalValue_pow}
\leanok
Cross-powering two witnesses gives a product witness whose score is the product of their scores. Transporting witnesses into and out of positive tensor powers gives $V_\tau(T\otimes S)\ge V_\tau(T)V_\tau(S)$ and $V_\tau(T^{\otimes n})=V_\tau(T)^n$ for $n\ge1$, for every real $\tau$. The power-zero identity holds in the admissible interval.
\end{lemma}

\noindent\textit{Auxiliary foundation. Source:} \texttt{material/main.tex:724}.
\textit{Status:} \texttt{formalized}. Main proof path.
\textit{Checked Lean declarations:} \leandecl{MatrixMultiplication.\allowbreak{}Tensor.\allowbreak{}Equivalent.\allowbreak{}kroneckerPower}, \leandecl{MatrixMultiplication.\allowbreak{}Tensor.\allowbreak{}kroneckerPower\_\allowbreak{}kronecker\_\allowbreak{}equivalent}, \leandecl{MatrixMultiplication.\allowbreak{}Tensor.\allowbreak{}kroneckerPower\_\allowbreak{}mul\_\allowbreak{}equivalent}, \leandecl{MatrixMultiplication.\allowbreak{}Tensor.\allowbreak{}symmetrize\_\allowbreak{}kroneckerPower\_\allowbreak{}equivalent}, \leandecl{MatrixMultiplication.\allowbreak{}Tensor.\allowbreak{}symmetrize\_\allowbreak{}kroneckerPower\_\allowbreak{}iterated\_\allowbreak{}equivalent}, \leandecl{MatrixMultiplication.\allowbreak{}Tensor.\allowbreak{}matrixMul\_\allowbreak{}copies\_\allowbreak{}kronecker\_\allowbreak{}equivalent}, \leandecl{MatrixMultiplication.\allowbreak{}Tensor.\allowbreak{}matrixMul\_\allowbreak{}copies\_\allowbreak{}kroneckerPower\_\allowbreak{}equivalent}, \leandecl{MatrixMultiplication.\allowbreak{}Tensor.\allowbreak{}ValueWitness.\allowbreak{}product}, \leandecl{MatrixMultiplication.\allowbreak{}Tensor.\allowbreak{}ValueWitness.\allowbreak{}tensorPower}, \leandecl{MatrixMultiplication.\allowbreak{}Tensor.\allowbreak{}ValueWitness.\allowbreak{}ofPower}, \leandecl{MatrixMultiplication.\allowbreak{}Tensor.\allowbreak{}ValueWitness.\allowbreak{}log\_\allowbreak{}score}, \leandecl{MatrixMultiplication.\allowbreak{}Tensor.\allowbreak{}ValueWitness.\allowbreak{}product\_\allowbreak{}score}, \leandecl{MatrixMultiplication.\allowbreak{}Tensor.\allowbreak{}ValueWitness.\allowbreak{}tensorPower\_\allowbreak{}score}, \leandecl{MatrixMultiplication.\allowbreak{}Tensor.\allowbreak{}ValueWitness.\allowbreak{}ofPower\_\allowbreak{}score\_\allowbreak{}pow}, \leandecl{MatrixMultiplication.\allowbreak{}Tensor.\allowbreak{}fractionalValue\_\allowbreak{}kronecker\_\allowbreak{}ge}, \leandecl{MatrixMultiplication.\allowbreak{}Tensor.\allowbreak{}fractionalValue\_\allowbreak{}pow\_\allowbreak{}of\_\allowbreak{}pos}, \leandecl{MatrixMultiplication.\allowbreak{}Tensor.\allowbreak{}fractionalValue\_\allowbreak{}pow}.
\textit{Formalization note:} Zero tensors and infinite values are included. Only the zero-power case requires the admissible parameter interval.

\begin{proof}
\proves{aux:value-products-powers}
\leanok
\uses{aux:fractional-value-definition,lem:sym-properties,aux:direct-sum-copies,aux:matrix-multiplication-calculus,aux:polynomial-degeneration-operations,aux:value-finiteness}
\textit{Checked proof.} Use explicit finite coordinate equivalences to rearrange products, iterated powers, and matrix blocks. For witnesses of orders $N,M$, take powers $M,N$ to obtain a common order $NM$. Prove the score identities via their logarithms, then distribute extended-real products and positive powers over suprema. Handle power zero by unit normalization.
\end{proof}

\renewcommand{\thetheorem}{A61}
\begin{lemma}[Absorbing matrix blocks into value witnesses]
\label{aux:value-block-scaling}
\lean{MatrixMultiplication.Tensor.valueSource_matrixBlock_equivalent,MatrixMultiplication.Tensor.ValueWitness.stripMatrixBlock,MatrixMultiplication.Tensor.ValueWitness.addMatrixBlock,MatrixMultiplication.Tensor.ValueWitness.score_mul_of_denominator_scale,MatrixMultiplication.Tensor.ValueWitness.score_of_numerator_scale,MatrixMultiplication.Tensor.fractionalValue_eq_mul_of_symmetrize_matrixBlock,MatrixMultiplication.Tensor.fractionalValue_diagonal_kronecker,MatrixMultiplication.Tensor.fractionalValue_directSumCopies,MatrixMultiplication.Tensor.fractionalValue_matrixMul_kronecker}
\leanok
If $S^{\otimes\mathfrak S_3}\cong\langle q,q,q\rangle^{\oplus r}\otimes T^{\otimes\mathfrak S_3}$ with $r,q\ge1$, then $V_\tau(S)=r^{1/6}q^{\tau/2}V_\tau(T)$. Consequently diagonal and matrix-multiplication factors obey the two exact scaling laws in Lemma 4.6.
\end{lemma}

\noindent\textit{Auxiliary foundation. Source:} \texttt{material/main.tex:724}.
\textit{Status:} \texttt{formalized}. Main proof path.
\textit{Checked Lean declarations:} \leandecl{MatrixMultiplication.\allowbreak{}Tensor.\allowbreak{}valueSource\_\allowbreak{}matrixBlock\_\allowbreak{}equivalent}, \leandecl{MatrixMultiplication.\allowbreak{}Tensor.\allowbreak{}ValueWitness.\allowbreak{}stripMatrixBlock}, \leandecl{MatrixMultiplication.\allowbreak{}Tensor.\allowbreak{}ValueWitness.\allowbreak{}addMatrixBlock}, \leandecl{MatrixMultiplication.\allowbreak{}Tensor.\allowbreak{}ValueWitness.\allowbreak{}score\_\allowbreak{}mul\_\allowbreak{}of\_\allowbreak{}denominator\_\allowbreak{}scale}, \leandecl{MatrixMultiplication.\allowbreak{}Tensor.\allowbreak{}ValueWitness.\allowbreak{}score\_\allowbreak{}of\_\allowbreak{}numerator\_\allowbreak{}scale}, \leandecl{MatrixMultiplication.\allowbreak{}Tensor.\allowbreak{}fractionalValue\_\allowbreak{}eq\_\allowbreak{}mul\_\allowbreak{}of\_\allowbreak{}symmetrize\_\allowbreak{}matrixBlock}, \leandecl{MatrixMultiplication.\allowbreak{}Tensor.\allowbreak{}fractionalValue\_\allowbreak{}diagonal\_\allowbreak{}kronecker}, \leandecl{MatrixMultiplication.\allowbreak{}Tensor.\allowbreak{}fractionalValue\_\allowbreak{}directSumCopies}, \leandecl{MatrixMultiplication.\allowbreak{}Tensor.\allowbreak{}fractionalValue\_\allowbreak{}matrixMul\_\allowbreak{}kronecker}.
\textit{Formalization note:} Checked for every real parameter, every tensor, and infinite values. Diagonal multiplicity may be zero; matrix dimensions are positive, as in the manuscript definition.

\begin{proof}
\proves{aux:value-block-scaling}
\leanok
\uses{aux:symmetrized-value-monotonicity,aux:value-products-powers,lem:sym-properties}
\textit{Checked proof.} For a witness of order $N$, stripping the block multiplies its denominator sizes by $r^N,q^N$. Adding the block multiplies its numerator sizes by the same factors. Their scores differ by $r^{1/6}q^{\tau/2}$. Take suprema in both directions, then apply the diagonal and matrix-multiplication symmetrization formulas.
\end{proof}

\renewcommand{\thetheorem}{A62}
\begin{lemma}[Uniform degeneration to polynomial multiplication]
\label{aux:polynomial-multiplication-degeneration}
\lean{MatrixMultiplication.Tensor.diagonal_restricts_polynomialMul_of_nodes,MatrixMultiplication.Tensor.diagonal_degenerates_polynomialMul}
\leanok
Over every field, $\langle a+b-1\rangle\trianglerighteq\mathrm{PM}_{a,b}$. More generally, distinct interpolation nodes give an explicit restriction from the diagonal over the field containing those nodes.
\end{lemma}

\noindent\textit{Auxiliary foundation. Source:} \texttt{material/main.tex:454}.
\textit{Status:} \texttt{formalized}. Main proof path.
\textit{Checked Lean declarations:} \leandecl{MatrixMultiplication.\allowbreak{}Tensor.\allowbreak{}diagonal\_\allowbreak{}restricts\_\allowbreak{}polynomialMul\_\allowbreak{}of\_\allowbreak{}nodes}, \leandecl{MatrixMultiplication.\allowbreak{}Tensor.\allowbreak{}diagonal\_\allowbreak{}degenerates\_\allowbreak{}polynomialMul}.
\textit{Formalization note:} The uniform degeneration part of Lemma 2.6 is checked, including zero input sizes. For positive input sizes, the exact border/asymptotic ranks and the ordinary-rank estimates are checked in their separate correspondence nodes.

\begin{proof}
\proves{aux:polynomial-multiplication-degeneration}
\leanok
\uses{aux:explicit-tensor-definitions,aux:rational-degeneration-equivalence}
\textit{Checked proof.} Evaluate at distinct nodes and extract coefficients from the Lagrange basis. Over the rational function field use $1,X,\ldots,X^{a+b-2}$, whose distinctness follows from polynomial degree. The output is exactly the constant polynomial-multiplication tensor, hence regular at zero. Apply the rational/polynomial degeneration equivalence.
\end{proof}

\renewcommand{\thetheorem}{A63}
\begin{lemma}[Lemma 2.6: exact border and asymptotic ranks]
\label{aux:polynomial-multiplication-ranks}
\manuscriptproof
\lean{MatrixMultiplication.Tensor.flatteningThird_rank_polynomialMul,MatrixMultiplication.Tensor.flatteningRank_polynomialMul_two,MatrixMultiplication.Tensor.add_sub_one_le_rank_polynomialMul,MatrixMultiplication.Tensor.borderRank_polynomialMul,MatrixMultiplication.Tensor.asymptoticRank_polynomialMul}
\leanok
\uses{def:polynomial-multiplication}
For every field $F$ and $a,b\ge1$, the third flattening of $\mathrm{PM}_{a,b}$ has rank $a+b-1$. Consequently, $R(\mathrm{PM}_{a,b})\ge a+b-1$ and $\underline R(\mathrm{PM}_{a,b})=\widetilde R(\mathrm{PM}_{a,b})=a+b-1$.
\end{lemma}

\noindent\textit{Auxiliary foundation. Source:} \texttt{material/main.tex:454}.
\textit{Status:} \texttt{formalized}. Additional manuscript coverage.
\textit{Proof supplied in manuscript.}
\textit{Checked Lean declarations:} \leandecl{MatrixMultiplication.\allowbreak{}Tensor.\allowbreak{}flatteningThird\_\allowbreak{}rank\_\allowbreak{}polynomialMul}, \leandecl{MatrixMultiplication.\allowbreak{}Tensor.\allowbreak{}flatteningRank\_\allowbreak{}polynomialMul\_\allowbreak{}two}, \leandecl{MatrixMultiplication.\allowbreak{}Tensor.\allowbreak{}add\_\allowbreak{}sub\_\allowbreak{}one\_\allowbreak{}le\_\allowbreak{}rank\_\allowbreak{}polynomialMul}, \leandecl{MatrixMultiplication.\allowbreak{}Tensor.\allowbreak{}borderRank\_\allowbreak{}polynomialMul}, \leandecl{MatrixMultiplication.\allowbreak{}Tensor.\allowbreak{}asymptoticRank\_\allowbreak{}polynomialMul}.
\textit{Formalization note:} This node isolates the checked field-uniform part of Lemma 2.6. Both input sizes must be positive; the formula does not extend to empty inputs. The separate ordinary-rank upper bounds are not claimed here.

\begin{proof}
\proves{aux:polynomial-multiplication-ranks}
\leanok
\uses{lem:flattening-lb,aux:border-rank-bound,aux:polynomial-multiplication-degeneration}
\textit{Checked proof.} For each output degree $k$, select the input pair $(\min(k,a-1),k-\min(k,a-1))$. These rows form an identity submatrix in the third flattening. Its rank bounds ordinary, border, and asymptotic rank from below. The diagonal degeneration gives border rank at most $a+b-1$, and $\widetilde R\le\underline R$ gives the asymptotic upper bound. This proves both identities without another finite-extension argument.
\end{proof}

\renewcommand{\thetheorem}{A64}
\begin{lemma}[Finite dependent tensor products and value bounds]
\label{aux:finite-tensor-products}
\lean{MatrixMultiplication.Tensor.tensorProductFamily,MatrixMultiplication.Tensor.tensorProductFamily_reindex,MatrixMultiplication.Tensor.tensorProductFamily_comp_equivalent,MatrixMultiplication.Tensor.tensorProductFamily_fin_succ_equivalent,MatrixMultiplication.Tensor.tensorProductFamily_fin_zero_equivalent,MatrixMultiplication.Tensor.fractionalValue_tensorProductFamily_ge}
\leanok
Finite tensor products may have dependent coordinate spaces. Reindexing their factors gives a coordinate equivalence. In the admissible interval, their fractional value is at least the product of the factor values, including the empty product.
\end{lemma}

\noindent\textit{Auxiliary foundation. Source:} \texttt{material/main.tex:626}.
\textit{Status:} \texttt{formalized}. Main proof path.
\textit{Checked Lean declarations:} \leandecl{MatrixMultiplication.\allowbreak{}Tensor.\allowbreak{}tensorProductFamily}, \leandecl{MatrixMultiplication.\allowbreak{}Tensor.\allowbreak{}tensorProductFamily\_\allowbreak{}reindex}, \leandecl{MatrixMultiplication.\allowbreak{}Tensor.\allowbreak{}tensorProductFamily\_\allowbreak{}comp\_\allowbreak{}equivalent}, \leandecl{MatrixMultiplication.\allowbreak{}Tensor.\allowbreak{}tensorProductFamily\_\allowbreak{}fin\_\allowbreak{}succ\_\allowbreak{}equivalent}, \leandecl{MatrixMultiplication.\allowbreak{}Tensor.\allowbreak{}tensorProductFamily\_\allowbreak{}fin\_\allowbreak{}zero\_\allowbreak{}equivalent}, \leandecl{MatrixMultiplication.\allowbreak{}Tensor.\allowbreak{}fractionalValue\_\allowbreak{}tensorProductFamily\_\allowbreak{}ge}.
\textit{Formalization note:} The tensor identities hold over commutative semirings; the value inequality holds over every field.

\begin{proof}
\proves{aux:finite-tensor-products}
\leanok
\uses{aux:value-products-powers,aux:symmetrized-value-monotonicity,aux:value-finiteness}
\textit{Checked proof.} Group coordinates into dependent function spaces. Prove invariance by reindexing the finite coefficient product. Split a finite product into its first factor and the remaining factors; use product supermultiplicativity inductively and unit normalization for the empty product.
\end{proof}

\renewcommand{\thetheorem}{A65}
\begin{lemma}[Monomial weights on coordinate pullbacks]
\label{aux:monomial-degeneration}
\lean{MatrixMultiplication.Tensor.degenerates_of_monomial_pullback}
\leanok
Assign nonnegative integer weights to the three factor maps of a coordinate pullback. If every supported coefficient has total weight at least $d$, the terms of weight exactly $d$ form a degeneration target.
\end{lemma}

\noindent\textit{Auxiliary foundation. Source:} \texttt{material/main.tex:565}.
\textit{Status:} \texttt{formalized}. Main proof path.
\textit{Checked Lean declarations:} \leandecl{MatrixMultiplication.\allowbreak{}Tensor.\allowbreak{}degenerates\_\allowbreak{}of\_\allowbreak{}monomial\_\allowbreak{}pullback}.
\textit{Formalization note:} The helper allows repeated selected coordinates and does not impose characteristic or field-size assumptions.

\begin{proof}
\proves{aux:monomial-degeneration}
\leanok
\uses{aux:polynomial-degeneration-definition}
\textit{Checked proof.} Use polynomial factor matrices with one selected coordinate per row and a monomial scalar. Divide the transformed tensor by the common power $X^d$ coefficientwise, separate its constant term, and record the remaining terms in the polynomial remainder.
\end{proof}

\renewcommand{\thetheorem}{A66}
\begin{lemma}[Independent maps in ordinary direct sums]
\label{aux:ordinary-sum-block-transport}
\lean{MatrixMultiplication.Tensor.directSumFamily,MatrixMultiplication.Tensor.map_directSumFamily,MatrixMultiplication.Tensor.Restricts.directSumFamily,MatrixMultiplication.Tensor.Equivalent.directSumFamily,MatrixMultiplication.Tensor.directSumFamily_const_equivalent_copies,MatrixMultiplication.Tensor.directSumFamily_equivalent_copies}
\leanok
Ordinary direct sums of dependent families respect independent restrictions and equivalences in every block. A direct sum of $Q$ blocks equivalent to a common tensor is equivalent to $Q$ copies of that tensor.
\end{lemma}

\noindent\textit{Auxiliary foundation. Source:} \texttt{material/main.tex:626}.
\textit{Status:} \texttt{formalized}. Main proof path.
\textit{Checked Lean declarations:} \leandecl{MatrixMultiplication.\allowbreak{}Tensor.\allowbreak{}directSumFamily}, \leandecl{MatrixMultiplication.\allowbreak{}Tensor.\allowbreak{}map\_\allowbreak{}directSumFamily}, \leandecl{MatrixMultiplication.\allowbreak{}Tensor.\allowbreak{}Restricts.\allowbreak{}directSumFamily}, \leandecl{MatrixMultiplication.\allowbreak{}Tensor.\allowbreak{}Equivalent.\allowbreak{}directSumFamily}, \leandecl{MatrixMultiplication.\allowbreak{}Tensor.\allowbreak{}directSumFamily\_\allowbreak{}const\_\allowbreak{}equivalent\_\allowbreak{}copies}, \leandecl{MatrixMultiplication.\allowbreak{}Tensor.\allowbreak{}directSumFamily\_\allowbreak{}equivalent\_\allowbreak{}copies}.
\textit{Formalization note:} These independent block maps are used only after separation. No analogous unconstrained congruence is asserted for partially direct sums.

\begin{proof}
\proves{aux:ordinary-sum-block-transport}
\leanok
\uses{aux:coefficient-tensors,aux:direct-sum-copies}
\textit{Checked proof.} Place each block factor map on the diagonal of the corresponding family matrix. Off-diagonal block tags give zero coefficients. Reindex a constant family by the finite cardinality to identify it with repeated copies.
\end{proof}

\renewcommand{\thetheorem}{A67}
\begin{lemma}[Words with prescribed counts]
\label{aux:word-type-cardinality}
\lean{MatrixMultiplication.wordCounts,MatrixMultiplication.WordType,MatrixMultiplication.wordCounts_sum,MatrixMultiplication.wordCounts_comp_perm,MatrixMultiplication.perm_of_wordCounts_eq,MatrixMultiplication.wordType_nonempty,MatrixMultiplication.wordType_nonempty_iff,MatrixMultiplication.card_wordType}
\leanok
Words of length $N$ with symbol counts $c_h$ exist exactly when $\sum_h c_h=N$. When this equality holds, their number is $\binom{N}{(c_h)_h}$, and any two such words differ by a permutation of their positions.
\end{lemma}

\noindent\textit{Auxiliary foundation. Source:} \texttt{material/main.tex:626}.
\textit{Status:} \texttt{formalized}. Main proof path.
\textit{Checked Lean declarations:} \leandecl{MatrixMultiplication.\allowbreak{}wordCounts}, \leandecl{MatrixMultiplication.\allowbreak{}WordType}, \leandecl{MatrixMultiplication.\allowbreak{}wordCounts\_\allowbreak{}sum}, \leandecl{MatrixMultiplication.\allowbreak{}wordCounts\_\allowbreak{}comp\_\allowbreak{}perm}, \leandecl{MatrixMultiplication.\allowbreak{}perm\_\allowbreak{}of\_\allowbreak{}wordCounts\_\allowbreak{}eq}, \leandecl{MatrixMultiplication.\allowbreak{}wordType\_\allowbreak{}nonempty}, \leandecl{MatrixMultiplication.\allowbreak{}wordType\_\allowbreak{}nonempty\_\allowbreak{}iff}, \leandecl{MatrixMultiplication.\allowbreak{}card\_\allowbreak{}wordType}.
\textit{Formalization note:} Zero counts, length zero, and empty alphabets are included.

\begin{proof}
\proves{aux:word-type-cardinality}
\leanok
\textit{Checked proof.} Construct a permutation from bijections of the symbol fibers. Apply orbit--stabilizer to the action of position permutations: the stabilizer has size $\prod_h c_h!$, while the full permutation group has size $N!$.
\end{proof}

\renewcommand{\thetheorem}{A68}
\begin{lemma}[Grouping tensor products by symbol counts]
\label{aux:word-product-grouping}
\lean{MatrixMultiplication.Tensor.wordCoordinateEquiv,MatrixMultiplication.Tensor.reindex_wordProduct,MatrixMultiplication.Tensor.wordProduct_equivalent_grouped,MatrixMultiplication.Tensor.wordCoordinateEquivOfCounts,MatrixMultiplication.Tensor.reindex_wordProduct_of_counts,MatrixMultiplication.Tensor.wordProduct_equivalent_grouped_of_counts}
\leanok
The tensor product along a word with counts $c_h$ is coordinate-isomorphic to $\bigotimes_h T_h^{\otimes c_h}$. All three coordinate types may depend on the symbol.
\end{lemma}

\noindent\textit{Auxiliary foundation. Source:} \texttt{material/main.tex:626}.
\textit{Status:} \texttt{formalized}. Main proof path.
\textit{Checked Lean declarations:} \leandecl{MatrixMultiplication.\allowbreak{}Tensor.\allowbreak{}wordCoordinateEquiv}, \leandecl{MatrixMultiplication.\allowbreak{}Tensor.\allowbreak{}reindex\_\allowbreak{}wordProduct}, \leandecl{MatrixMultiplication.\allowbreak{}Tensor.\allowbreak{}wordProduct\_\allowbreak{}equivalent\_\allowbreak{}grouped}, \leandecl{MatrixMultiplication.\allowbreak{}Tensor.\allowbreak{}wordCoordinateEquivOfCounts}, \leandecl{MatrixMultiplication.\allowbreak{}Tensor.\allowbreak{}reindex\_\allowbreak{}wordProduct\_\allowbreak{}of\_\allowbreak{}counts}, \leandecl{MatrixMultiplication.\allowbreak{}Tensor.\allowbreak{}wordProduct\_\allowbreak{}equivalent\_\allowbreak{}grouped\_\allowbreak{}of\_\allowbreak{}counts}.
\textit{Formalization note:} The explicit coordinate identifications include zero counts and empty products. They are applied separately to ordinary-direct-sum blocks after separation.

\begin{proof}
\proves{aux:word-product-grouping}
\leanok
\uses{aux:word-type-cardinality,aux:finite-tensor-products}
\textit{Checked proof.} Split positions into fibers of the word and enumerate each fiber by its count. The resulting dependent coordinate equivalences group the coefficient product into the corresponding tensor powers.
\end{proof}

\renewcommand{\thetheorem}{A69}
\begin{lemma}[Selecting words from powers of partial sums]
\label{aux:partial-sum-word-selection}
\lean{MatrixMultiplication.Tensor.restricts_pullback,MatrixMultiplication.Tensor.kroneckerPower_partialSumFamily_restricts}
\leanok
A power of a partial direct sum restricts to the partial sum of tensor products indexed by any selected family of distinct block words, with its first factor still shared.
\end{lemma}

\noindent\textit{Auxiliary foundation. Source:} \texttt{material/main.tex:626}.
\textit{Status:} \texttt{formalized}. Main proof path.
\textit{Checked Lean declarations:} \leandecl{MatrixMultiplication.\allowbreak{}Tensor.\allowbreak{}restricts\_\allowbreak{}pullback}, \leandecl{MatrixMultiplication.\allowbreak{}Tensor.\allowbreak{}kroneckerPower\_\allowbreak{}partialSumFamily\_\allowbreak{}restricts}.
\textit{Formalization note:} This proves the needed selection directly, without assuming that independent first-factor equivalences preserve partial sums.

\begin{proof}
\proves{aux:partial-sum-word-selection}
\leanok
\uses{def:partial-direct-sum,aux:finite-tensor-products}
\textit{Checked proof.} Pull back the second and third coordinate functions along the selected words and leave the first coordinate unchanged. Distinct words differ at some position, making the cross-block coefficient vanish.
\end{proof}

\renewcommand{\thetheorem}{A70}
\begin{lemma}[From all type bounds to a sum bound]
\label{aux:multinomial-sum-bound}
\lean{MatrixMultiplication.le_of_pow_le_polynomial_mul_pow,MatrixMultiplication.card_piAntidiag_univ_le,MatrixMultiplication.sum_le_of_multinomial_bounds}
\leanok
For nonnegative $a_h,B$ and $C>0$, suppose every count vector of total $N\ge1$ satisfies $\binom{N}{(c_h)_h}\prod_h a_h^{c_h}\le C B^N$. Then $\sum_h a_h\le B$.
\end{lemma}

\noindent\textit{Auxiliary foundation. Source:} \texttt{material/main.tex:868}.
\textit{Status:} \texttt{formalized}. Main proof path.
\textit{Checked Lean declarations:} \leandecl{MatrixMultiplication.\allowbreak{}le\_\allowbreak{}of\_\allowbreak{}pow\_\allowbreak{}le\_\allowbreak{}polynomial\_\allowbreak{}mul\_\allowbreak{}pow}, \leandecl{MatrixMultiplication.\allowbreak{}card\_\allowbreak{}piAntidiag\_\allowbreak{}univ\_\allowbreak{}le}, \leandecl{MatrixMultiplication.\allowbreak{}sum\_\allowbreak{}le\_\allowbreak{}of\_\allowbreak{}multinomial\_\allowbreak{}bounds}.
\textit{Formalization note:} Zero weights and empty families are allowed. This replaces the entropy asymptotic and rational approximation on the main proof path.

\begin{proof}
\proves{aux:multinomial-sum-bound}
\leanok
\uses{aux:polynomial-overhead-limit}
\textit{Checked proof.} Expand $(\sum_h a_h)^N$ by the multinomial theorem and sum the individual bounds. The count vectors inject into $[0,N]^H$, so there are at most $(N+1)^{|H|}$. Take $N$-th roots; the polynomial factor and the fixed positive constant both tend to one.
\end{proof}

\renewcommand{\thetheorem}{A71}
\begin{lemma}[Weighted AM--GM form of the entropy bound]
\label{aux:entropy-variational-upper-bound}
\lean{MatrixMultiplication.two_rpow_entropy_eq_prod,MatrixMultiplication.entropy_variational_le,MatrixMultiplication.entropy_variational_rpow_le}
\leanok
For a probability vector $q$ and nonnegative $a_h$, $2^{H(q)}\prod_h a_h^{q_h}\le\sum_h a_h$. For $\tau>0$, this gives $2^{\tau H(q)}\prod_h v_h^{q_h}\le(\sum_h v_h^{1/\tau})^\tau$.
\end{lemma}

\noindent\textit{Auxiliary foundation. Source:} \texttt{material/main.tex:868}.
\textit{Status:} \texttt{formalized}. Additional manuscript coverage.
\textit{Checked Lean declarations:} \leandecl{MatrixMultiplication.\allowbreak{}two\_\allowbreak{}rpow\_\allowbreak{}entropy\_\allowbreak{}eq\_\allowbreak{}prod}, \leandecl{MatrixMultiplication.\allowbreak{}entropy\_\allowbreak{}variational\_\allowbreak{}le}, \leandecl{MatrixMultiplication.\allowbreak{}entropy\_\allowbreak{}variational\_\allowbreak{}rpow\_\allowbreak{}le}.
\textit{Formalization note:} This is the checked upper-bound half of the entropy variational identity. The attained maximum and equality are checked separately in the manuscript lemma node.

\begin{proof}
\proves{aux:entropy-variational-upper-bound}
\leanok
\uses{aux:entropy}
\textit{Checked proof.} Use weighted AM--GM on the positive support and the product formula for exponential entropy. Handle a positive probability at a zero base directly; zero probabilities use the convention $0^0=1$.
\end{proof}

\renewcommand{\thetheorem}{A72}
\begin{lemma}[Coordinate models for polynomial multiplication]
\label{aux:polynomial-coordinate-models}
\lean{MatrixMultiplication.Tensor.swapLastTwo,MatrixMultiplication.Tensor.polynomialMulCopies,MatrixMultiplication.Tensor.polynomialMulCopies_equivalent,MatrixMultiplication.Tensor.polynomialMul_swapFirstTwo,MatrixMultiplication.Tensor.polynomialMul_one_equivalent_matrixMul}
\leanok
The repeated coefficient model $\mathrm{PM}^{[r]}_{a,b}$ has the literal shared first factor $F^a$ and equals $\mathrm{PM}_{a,b}\otimes\langle1,1,r\rangle$ up to canonical singleton-coordinate identifications. Swapping the inputs exchanges $a,b$, and $\mathrm{PM}_{1,b}$ is equivalent to $\langle1,1,b\rangle$.
\end{lemma}

\noindent\textit{Auxiliary foundation. Source:} \texttt{material/main.tex:446}.
\textit{Status:} \texttt{formalized}. Main proof path.
\textit{Checked Lean declarations:} \leandecl{MatrixMultiplication.\allowbreak{}Tensor.\allowbreak{}swapLastTwo}, \leandecl{MatrixMultiplication.\allowbreak{}Tensor.\allowbreak{}polynomialMulCopies}, \leandecl{MatrixMultiplication.\allowbreak{}Tensor.\allowbreak{}polynomialMulCopies\_\allowbreak{}equivalent}, \leandecl{MatrixMultiplication.\allowbreak{}Tensor.\allowbreak{}polynomialMul\_\allowbreak{}swapFirstTwo}, \leandecl{MatrixMultiplication.\allowbreak{}Tensor.\allowbreak{}polynomialMul\_\allowbreak{}one\_\allowbreak{}equivalent\_\allowbreak{}matrixMul}.

\begin{proof}
\proves{aux:polynomial-coordinate-models}
\leanok
\uses{def:polynomial-multiplication,aux:factor-permutations,aux:coefficient-tensors}
\textit{Checked proof.} Give explicit coordinate bijections. Repeated copies retain the first input index; only the canonical one-dimensional matrix coordinate is inserted. The last-two-factor swap is also defined explicitly.
\end{proof}

\renewcommand{\thetheorem}{A73}
\begin{lemma}[Degeneration from three split-map identities]
\label{aux:split-coefficient-degeneration}
\lean{MatrixMultiplication.Tensor.degenerates_partialSum_of_split_maps}
\leanok
If $\operatorname{map}(1,B_0,C_0)T=T_0$, $\operatorname{map}(1,B_0,C_1)T=0$, and $\operatorname{map}(1,B_1,C_1)T=T_1$, then $T\trianglerighteq T_0\diamondplus T_1$. The fourth block is unrestricted. This holds over every commutative semiring.
\end{lemma}

\noindent\textit{Auxiliary foundation. Source:} \texttt{material/main.tex:885}.
\textit{Status:} \texttt{formalized}. Main proof path.
\textit{Checked Lean declarations:} \leandecl{MatrixMultiplication.\allowbreak{}Tensor.\allowbreak{}degenerates\_\allowbreak{}partialSum\_\allowbreak{}of\_\allowbreak{}split\_\allowbreak{}maps}.
\textit{Formalization note:} A checked coordinate criterion sufficient for polynomial multiplication. It does not claim the full abstract short-exact-sequence lemma.

\begin{proof}
\proves{aux:split-coefficient-degeneration}
\leanok
\uses{aux:monomial-degeneration,aux:degeneration-calculus,def:partial-direct-sum}
\textit{Checked proof.} Combine the two second-factor maps and two third-factor maps into sum-indexed maps. Assign second-factor weights $(0,1)$ and third-factor weights $(1,0)$. At order one the two diagonal blocks survive, the zero block causes no lower term, and the fourth block has higher order.
\end{proof}

\renewcommand{\thetheorem}{A74}
\begin{lemma}[Explicit maps on bivariate coefficient grids]
\label{aux:bivariate-grid-maps}
\lean{MatrixMultiplication.Tensor.gridEdgeMap,MatrixMultiplication.Tensor.sum_gridEdgeMap_mul,MatrixMultiplication.Tensor.gridSection,MatrixMultiplication.Tensor.gridSection_degree,MatrixMultiplication.Tensor.gridSectionMap,MatrixMultiplication.Tensor.sum_gridSectionMap_mul,MatrixMultiplication.Tensor.gridDiagonalCoefficient,MatrixMultiplication.Tensor.gridTailCoefficient,MatrixMultiplication.Tensor.gridDiagonalCoefficient_edge,MatrixMultiplication.Tensor.gridTailCoefficient_edge}
\leanok
On the rectangular coefficient grid, multiplication by $x-y$ is an adjacent-edge difference; $x^s\mapsto x^{\min(s,b-1)}y^{s-\min(s,b-1)}$ is a section of diagonal evaluation. Diagonal summation kills edge differences and diagonal-tail extraction recovers their coefficients.
\end{lemma}

\noindent\textit{Auxiliary foundation. Source:} \texttt{material/main.tex:918}.
\textit{Status:} \texttt{formalized}. Main proof path.
\textit{Checked Lean declarations:} \leandecl{MatrixMultiplication.\allowbreak{}Tensor.\allowbreak{}gridEdgeMap}, \leandecl{MatrixMultiplication.\allowbreak{}Tensor.\allowbreak{}sum\_\allowbreak{}gridEdgeMap\_\allowbreak{}mul}, \leandecl{MatrixMultiplication.\allowbreak{}Tensor.\allowbreak{}gridSection}, \leandecl{MatrixMultiplication.\allowbreak{}Tensor.\allowbreak{}gridSection\_\allowbreak{}degree}, \leandecl{MatrixMultiplication.\allowbreak{}Tensor.\allowbreak{}gridSectionMap}, \leandecl{MatrixMultiplication.\allowbreak{}Tensor.\allowbreak{}sum\_\allowbreak{}gridSectionMap\_\allowbreak{}mul}, \leandecl{MatrixMultiplication.\allowbreak{}Tensor.\allowbreak{}gridDiagonalCoefficient}, \leandecl{MatrixMultiplication.\allowbreak{}Tensor.\allowbreak{}gridTailCoefficient}, \leandecl{MatrixMultiplication.\allowbreak{}Tensor.\allowbreak{}gridDiagonalCoefficient\_\allowbreak{}edge}, \leandecl{MatrixMultiplication.\allowbreak{}Tensor.\allowbreak{}gridTailCoefficient\_\allowbreak{}edge}.

\begin{proof}
\proves{aux:bivariate-grid-maps}
\leanok
\uses{aux:coefficient-tensors}
\textit{Checked proof.} Check the two endpoint indices and their degree bounds, then prove the diagonal and tail indicator identities. The cancellation uses only ring arithmetic, so it also holds in characteristic two.
\end{proof}

\renewcommand{\thetheorem}{A75}
\begin{lemma}[Real-root value identities for binary partial sums]
\label{aux:binary-partial-value-support}
\lean{MatrixMultiplication.Tensor.fractionalValue_partialSum_toReal_add_le,MatrixMultiplication.Tensor.polynomialMulCopies_value_root,MatrixMultiplication.Tensor.fractionalValue_swapLastTwo,MatrixMultiplication.Tensor.polynomialMul_value_swap}
\leanok
In the admissible interval, $V_\tau(T\diamondplus S)^{1/\tau}\ge V_\tau(T)^{1/\tau}+V_\tau(S)^{1/\tau}$, and $V_\tau(\mathrm{PM}^{[r]}_{a,b})^{1/\tau}=rV_\tau(\mathrm{PM}_{a,b})^{1/\tau}$ for $r\ge1$. Exchanging the last two factors preserves value, as does exchanging the two polynomial input sizes.
\end{lemma}

\noindent\textit{Auxiliary foundation. Source:} \texttt{material/main.tex:936}.
\textit{Status:} \texttt{formalized}. Main proof path.
\textit{Checked Lean declarations:} \leandecl{MatrixMultiplication.\allowbreak{}Tensor.\allowbreak{}fractionalValue\_\allowbreak{}partialSum\_\allowbreak{}toReal\_\allowbreak{}add\_\allowbreak{}le}, \leandecl{MatrixMultiplication.\allowbreak{}Tensor.\allowbreak{}polynomialMulCopies\_\allowbreak{}value\_\allowbreak{}root}, \leandecl{MatrixMultiplication.\allowbreak{}Tensor.\allowbreak{}fractionalValue\_\allowbreak{}swapLastTwo}, \leandecl{MatrixMultiplication.\allowbreak{}Tensor.\allowbreak{}polynomialMul\_\allowbreak{}value\_\allowbreak{}swap}.

\begin{proof}
\proves{aux:binary-partial-value-support}
\leanok
\uses{lem:value-diamondsum,lem:value-properties,aux:value-finiteness,aux:polynomial-coordinate-models}
\textit{Checked proof.} Transport the finite-family inequality through the two-element family equivalence while preserving the shared first coordinate. Apply rectangular matrix scaling and permutation invariance; admissible-interval finiteness justifies passing to real values.
\end{proof}

\renewcommand{\thetheorem}{A76}
\begin{lemma}[Uniform upper bound for polynomial value]
\label{aux:polynomial-value-upper-bound}
\lean{MatrixMultiplication.Tensor.polynomialMulValue_rpow_le,MatrixMultiplication.Tensor.polynomialMulValue_diagonal_rpow_le}
\leanok
For admissible $\tau$, $P(a,b)^\tau\le a+b-1$ for positive $a,b$. In particular, $P(a,a)^\tau\le2a-1$ for every $a\ge1$.
\end{lemma}

\noindent\textit{Auxiliary foundation. Source:} \texttt{material/main.tex:1182}.
\textit{Status:} \texttt{formalized}. Main proof path.
\textit{Checked Lean declarations:} \leandecl{MatrixMultiplication.\allowbreak{}Tensor.\allowbreak{}polynomialMulValue\_\allowbreak{}rpow\_\allowbreak{}le}, \leandecl{MatrixMultiplication.\allowbreak{}Tensor.\allowbreak{}polynomialMulValue\_\allowbreak{}diagonal\_\allowbreak{}rpow\_\allowbreak{}le}.

\begin{proof}
\proves{aux:polynomial-value-upper-bound}
\leanok
\uses{aux:polynomial-value-hypotheses,aux:value-finiteness,aux:polynomial-multiplication-degeneration,aux:asymptotic-rank-degeneration,aux:diagonal-rank-normalization}
\textit{Checked proof.} Cancel the positive reciprocal exponent in $P$. Bound fractional value by asymptotic rank, and use the diagonal-to-polynomial-multiplication degeneration. Only the upper bound is needed, not the exact rank formula.
\end{proof}

\renewcommand{\thetheorem}{A77}
\begin{lemma}[Subrank maxima and asymptotic subrank supremum]
\label{aux:subrank-invariants}
\lean{MatrixMultiplication.Tensor.subrank,MatrixMultiplication.Tensor.borderSubrank,MatrixMultiplication.Tensor.asymptoticSubrank,MatrixMultiplication.Tensor.subrankValues_bddAbove,MatrixMultiplication.Tensor.restricts_diagonal_of_subrank,MatrixMultiplication.Tensor.le_subrank_of_restricts_diagonal,MatrixMultiplication.Tensor.subrank_mono,MatrixMultiplication.Tensor.subrank_le_rank,MatrixMultiplication.Tensor.borderSubrankValues_bddAbove,MatrixMultiplication.Tensor.degenerates_diagonal_of_borderSubrank,MatrixMultiplication.Tensor.le_borderSubrank_of_degenerates_diagonal,MatrixMultiplication.Tensor.borderSubrank_mono,MatrixMultiplication.Tensor.subrank_le_borderSubrank,MatrixMultiplication.Tensor.borderSubrank_le_borderRank,MatrixMultiplication.Tensor.subrankRootValues_nonempty,MatrixMultiplication.Tensor.subrankRootValues_bddAbove,MatrixMultiplication.Tensor.root_le_asymptoticSubrank,MatrixMultiplication.Tensor.asymptoticSubrank_nonneg,MatrixMultiplication.Tensor.subrank_le_asymptoticSubrank,MatrixMultiplication.Tensor.asymptoticSubrank_le_rank,MatrixMultiplication.Tensor.subrank_zero,MatrixMultiplication.Tensor.borderSubrank_zero,MatrixMultiplication.Tensor.asymptoticSubrank_zero,MatrixMultiplication.Tensor.subrank_diagonal,MatrixMultiplication.Tensor.borderSubrank_diagonal,MatrixMultiplication.Tensor.asymptoticSubrank_diagonal}
\leanok
The maxima defining $Q(T)$ and $\underline Q(T)$ are attained, with $Q(T)\le\underline Q(T)\le\underline R(T)$. Restriction is monotone for subrank and degeneration is monotone for border subrank. The set $\{Q(T^{\otimes n})^{1/n}:n\ge1\}$ is nonempty and bounded by $R(T)$, and its supremum defines $\widetilde Q(T)$. All three invariants are zero on the zero tensor and equal $r$ on $\langle r\rangle$.
\end{lemma}

\noindent\textit{Auxiliary foundation. Source:} \texttt{material/main.tex:324}.
\textit{Status:} \texttt{formalized}. Additional manuscript coverage.
\textit{Checked Lean declarations:} \leandecl{MatrixMultiplication.\allowbreak{}Tensor.\allowbreak{}subrank}, \leandecl{MatrixMultiplication.\allowbreak{}Tensor.\allowbreak{}borderSubrank}, \leandecl{MatrixMultiplication.\allowbreak{}Tensor.\allowbreak{}asymptoticSubrank}, \leandecl{MatrixMultiplication.\allowbreak{}Tensor.\allowbreak{}subrankValues\_\allowbreak{}bddAbove}, \leandecl{MatrixMultiplication.\allowbreak{}Tensor.\allowbreak{}restricts\_\allowbreak{}diagonal\_\allowbreak{}of\_\allowbreak{}subrank}, \leandecl{MatrixMultiplication.\allowbreak{}Tensor.\allowbreak{}le\_\allowbreak{}subrank\_\allowbreak{}of\_\allowbreak{}restricts\_\allowbreak{}diagonal}, \leandecl{MatrixMultiplication.\allowbreak{}Tensor.\allowbreak{}subrank\_\allowbreak{}mono}, \leandecl{MatrixMultiplication.\allowbreak{}Tensor.\allowbreak{}subrank\_\allowbreak{}le\_\allowbreak{}rank}, \leandecl{MatrixMultiplication.\allowbreak{}Tensor.\allowbreak{}borderSubrankValues\_\allowbreak{}bddAbove}, \leandecl{MatrixMultiplication.\allowbreak{}Tensor.\allowbreak{}degenerates\_\allowbreak{}diagonal\_\allowbreak{}of\_\allowbreak{}borderSubrank}, \leandecl{MatrixMultiplication.\allowbreak{}Tensor.\allowbreak{}le\_\allowbreak{}borderSubrank\_\allowbreak{}of\_\allowbreak{}degenerates\_\allowbreak{}diagonal}, \leandecl{MatrixMultiplication.\allowbreak{}Tensor.\allowbreak{}borderSubrank\_\allowbreak{}mono}, \leandecl{MatrixMultiplication.\allowbreak{}Tensor.\allowbreak{}subrank\_\allowbreak{}le\_\allowbreak{}borderSubrank}, \leandecl{MatrixMultiplication.\allowbreak{}Tensor.\allowbreak{}borderSubrank\_\allowbreak{}le\_\allowbreak{}borderRank}, \leandecl{MatrixMultiplication.\allowbreak{}Tensor.\allowbreak{}subrankRootValues\_\allowbreak{}nonempty}, \leandecl{MatrixMultiplication.\allowbreak{}Tensor.\allowbreak{}subrankRootValues\_\allowbreak{}bddAbove}, \leandecl{MatrixMultiplication.\allowbreak{}Tensor.\allowbreak{}root\_\allowbreak{}le\_\allowbreak{}asymptoticSubrank}, \leandecl{MatrixMultiplication.\allowbreak{}Tensor.\allowbreak{}asymptoticSubrank\_\allowbreak{}nonneg}, \leandecl{MatrixMultiplication.\allowbreak{}Tensor.\allowbreak{}subrank\_\allowbreak{}le\_\allowbreak{}asymptoticSubrank}, \leandecl{MatrixMultiplication.\allowbreak{}Tensor.\allowbreak{}asymptoticSubrank\_\allowbreak{}le\_\allowbreak{}rank}, \leandecl{MatrixMultiplication.\allowbreak{}Tensor.\allowbreak{}subrank\_\allowbreak{}zero}, \leandecl{MatrixMultiplication.\allowbreak{}Tensor.\allowbreak{}borderSubrank\_\allowbreak{}zero}, \leandecl{MatrixMultiplication.\allowbreak{}Tensor.\allowbreak{}asymptoticSubrank\_\allowbreak{}zero}, \leandecl{MatrixMultiplication.\allowbreak{}Tensor.\allowbreak{}subrank\_\allowbreak{}diagonal}, \leandecl{MatrixMultiplication.\allowbreak{}Tensor.\allowbreak{}borderSubrank\_\allowbreak{}diagonal}, \leandecl{MatrixMultiplication.\allowbreak{}Tensor.\allowbreak{}asymptoticSubrank\_\allowbreak{}diagonal}.
\textit{Formalization note:} Definitions agree with Definition 2.1 over every field, including zero tensors and empty coordinate sets. Maxima are attained, and the real supremum is bounded; no convention for an unbounded supremum is used.

\begin{proof}
\proves{aux:subrank-invariants}
\leanok
\uses{aux:ordinary-rank-definition,aux:border-rank-bound,aux:diagonal-rank-normalization,aux:power-rank-submultiplicativity}
\textit{Checked proof.} Zero is an admissible target. Rank bounds every diagonal restriction target, and border rank bounds every diagonal degeneration target, giving finite attained maxima. Submultiplicativity of rank uniformly bounds the subrank roots. Normalization follows from the unit-power lower bound and the diagonal rank upper bound.
\end{proof}

\renewcommand{\thetheorem}{A78}
\begin{lemma}[Finiteness forces unit normalization and the upper endpoint]
\label{aux:finite-unit-value}
\lean{MatrixMultiplication.Tensor.fractionalValue_unit_eq_one_of_ne_top,MatrixMultiplication.Tensor.fractionalValue_diagonal_of_unit_ne_top,MatrixMultiplication.Tensor.fractionalValue_matrixMul_of_unit_ne_top,MatrixMultiplication.Tensor.fractionalValue_le_rank_of_unit_ne_top,MatrixMultiplication.Tensor.fractionalValue_ne_top_of_unit_ne_top,MatrixMultiplication.Tensor.fractionalValue_le_asymptoticRank_of_unit_ne_top,MatrixMultiplication.Tensor.le_omega_div_three_of_unit_ne_top,MatrixMultiplication.Tensor.fractionalValue_unit_eq_top_of_omega_div_three_lt}
\leanok
If $V_\tau(\langle1\rangle)<\infty$, then $V_\tau(\langle1\rangle)=1$, $V_\tau(\langle r\rangle)=r$, $V_\tau(\langle\ell,m,n\rangle)=(\ell mn)^\tau$, and $V_\tau(T)\le\widetilde R(T)$ for every tensor. Consequently $\tau\le\omega/3$.
\end{lemma}

\noindent\textit{Auxiliary foundation. Source:} \texttt{material/main.tex:697}.
\textit{Status:} \texttt{formalized}. Additional manuscript coverage.
\textit{Checked Lean declarations:} \leandecl{MatrixMultiplication.\allowbreak{}Tensor.\allowbreak{}fractionalValue\_\allowbreak{}unit\_\allowbreak{}eq\_\allowbreak{}one\_\allowbreak{}of\_\allowbreak{}ne\_\allowbreak{}top}, \leandecl{MatrixMultiplication.\allowbreak{}Tensor.\allowbreak{}fractionalValue\_\allowbreak{}diagonal\_\allowbreak{}of\_\allowbreak{}unit\_\allowbreak{}ne\_\allowbreak{}top}, \leandecl{MatrixMultiplication.\allowbreak{}Tensor.\allowbreak{}fractionalValue\_\allowbreak{}matrixMul\_\allowbreak{}of\_\allowbreak{}unit\_\allowbreak{}ne\_\allowbreak{}top}, \leandecl{MatrixMultiplication.\allowbreak{}Tensor.\allowbreak{}fractionalValue\_\allowbreak{}le\_\allowbreak{}rank\_\allowbreak{}of\_\allowbreak{}unit\_\allowbreak{}ne\_\allowbreak{}top}, \leandecl{MatrixMultiplication.\allowbreak{}Tensor.\allowbreak{}fractionalValue\_\allowbreak{}ne\_\allowbreak{}top\_\allowbreak{}of\_\allowbreak{}unit\_\allowbreak{}ne\_\allowbreak{}top}, \leandecl{MatrixMultiplication.\allowbreak{}Tensor.\allowbreak{}fractionalValue\_\allowbreak{}le\_\allowbreak{}asymptoticRank\_\allowbreak{}of\_\allowbreak{}unit\_\allowbreak{}ne\_\allowbreak{}top}, \leandecl{MatrixMultiplication.\allowbreak{}Tensor.\allowbreak{}le\_\allowbreak{}omega\_\allowbreak{}div\_\allowbreak{}three\_\allowbreak{}of\_\allowbreak{}unit\_\allowbreak{}ne\_\allowbreak{}top}, \leandecl{MatrixMultiplication.\allowbreak{}Tensor.\allowbreak{}fractionalValue\_\allowbreak{}unit\_\allowbreak{}eq\_\allowbreak{}top\_\allowbreak{}of\_\allowbreak{}omega\_\allowbreak{}div\_\allowbreak{}three\_\allowbreak{}lt}.
\textit{Formalization note:} The real parameter is arbitrary; no admissible-interval hypothesis or circular appeal to normalization is used. Finiteness is established before converting values to real numbers.

\begin{proof}
\proves{aux:finite-unit-value}
\leanok
\uses{aux:value-products-powers,aux:value-block-scaling,aux:unit-value-witness,aux:ordinary-rank-definition,aux:asymptotic-rank-definition,eq:omega-log-n}
\textit{Checked proof.} Supermultiplicativity gives $V_\tau(\langle1\rangle)^2\le V_\tau(\langle1\rangle)$; the score-one witness gives the reverse lower bound by one. Under finiteness these imply unit normalization. Apply the scaling laws, rank restriction sources, positive-power multiplicativity, and the infimum of rank roots. At the size-two matrix tensor the bound becomes $8^\tau\le2^\omega$.
\end{proof}

\renewcommand{\thetheorem}{A79}
\begin{lemma}[An explicit quadratic diagonal degeneration]
\label{aux:quadratic-matrix-subrank}
\lean{MatrixMultiplication.Tensor.matrixMul_degenerates_diagonal_quadratic,MatrixMultiplication.Tensor.sq_le_borderSubrank_matrixMul_double}
\leanok
For every $q\ge0$, $\langle2q,2q,2q\rangle\trianglerighteq\langle q^2\rangle$. Over fields, $q^2\le\underline Q(\langle2q,2q,2q\rangle)$.
\end{lemma}

\noindent\textit{Auxiliary foundation. Source:} \texttt{material/main.tex:697}.
\textit{Status:} \texttt{formalized}. Additional manuscript coverage.
\textit{Checked Lean declarations:} \leandecl{MatrixMultiplication.\allowbreak{}Tensor.\allowbreak{}matrixMul\_\allowbreak{}degenerates\_\allowbreak{}diagonal\_\allowbreak{}quadratic}, \leandecl{MatrixMultiplication.\allowbreak{}Tensor.\allowbreak{}sq\_\allowbreak{}le\_\allowbreak{}borderSubrank\_\allowbreak{}matrixMul\_\allowbreak{}double}.
\textit{Formalization note:} This elementary quadratic bound is sufficient for lower-parameter divergence and is distinct from the sharper cited bound. The degeneration works over every commutative semiring and includes q=0.

\begin{proof}
\proves{aux:quadratic-matrix-subrank}
\leanok
\uses{aux:monomial-degeneration,aux:coefficient-tensors,aux:subrank-invariants}
\textit{Checked proof.} Index the target diagonal by pairs $(i,j)$ with $0\le i,j<q$, selecting source edges $(i,j)$, $(j,i+j)$, and $(i,i+j)$. Use the nonnegative natural weights $(i+j)^2$, $q^2+u^2-v^2$, and $4q^2-2w(w+z)$ on the three target factors. On the support, their sum is $5q^2$ plus a square, with equality exactly when all target pairs coincide. Apply the monomial degeneration criterion and reindex the $q^2$ pairs.
\end{proof}

\renewcommand{\thetheorem}{A80}
\begin{theorem}[Prove2Me exponent identification and nine-quarters corollary]
\label{aux:prove2me-statement}
\lean{MME.MMSpace,MME.MMTensor,MME.tensorRank,MME.matMulExp,MME.mmCoefficients_MMTensor,MME.mm_decomposition_iff,MME.tensorRank_MMTensor_eq_rank,MatrixMultiplication.omega_le_log_rank_matrixMul,MatrixMultiplication.log_rank_matrixMul_pow_two_tendsto,MatrixMultiplication.iInf_log_rank_matrixMul_eq_omega,MME.matMulExp_eq_omega,mme_omega_le_225}
\leanok
\uses{aux:matrix-multiplication-exponent}
For every field $K$, the downloaded Prove2Me exponent satisfies $\operatorname{MME.matMulExp}(K)=\omega_K\le 9/4=2.25$. Here $\operatorname{MME.matMulExp}$ is the infimum of $\log R(\mathrm{MM}_n)/\log n$ for $n\ge2$, with value $3$ at sizes zero and one, using pure-tensor decompositions in a tensor product of vector spaces.
\end{theorem}

\noindent\textit{Auxiliary foundation. Source:} \texttt{material/prove2me/mme\_omega\_le\_225.lean.txt:5}.
\textit{Status:} \texttt{formalized}. Additional manuscript coverage.
\textit{Checked Lean declarations:} \leandecl{MME.\allowbreak{}MMSpace}, \leandecl{MME.\allowbreak{}MMTensor}, \leandecl{MME.\allowbreak{}tensorRank}, \leandecl{MME.\allowbreak{}matMulExp}, \leandecl{MME.\allowbreak{}mmCoefficients\_\allowbreak{}MMTensor}, \leandecl{MME.\allowbreak{}mm\_\allowbreak{}decomposition\_\allowbreak{}iff}, \leandecl{MME.\allowbreak{}tensorRank\_\allowbreak{}MMTensor\_\allowbreak{}eq\_\allowbreak{}rank}, \leandecl{MatrixMultiplication.\allowbreak{}omega\_\allowbreak{}le\_\allowbreak{}log\_\allowbreak{}rank\_\allowbreak{}matrixMul}, \leandecl{MatrixMultiplication.\allowbreak{}log\_\allowbreak{}rank\_\allowbreak{}matrixMul\_\allowbreak{}pow\_\allowbreak{}two\_\allowbreak{}tendsto}, \leandecl{MatrixMultiplication.\allowbreak{}iInf\_\allowbreak{}log\_\allowbreak{}rank\_\allowbreak{}matrixMul\_\allowbreak{}eq\_\allowbreak{}omega}, \leandecl{MME.\allowbreak{}matMulExp\_\allowbreak{}eq\_\allowbreak{}omega}, \leandecl{mme\_\allowbreak{}omega\_\allowbreak{}le\_\allowbreak{}225}.
\textit{Formalization note:} External statement downloaded from Prove2Me on 2026-11-02; source URLs and hashes are preserved in material/prove2me/provenance.json. Its definition bodies are unchanged. The tensor bridge accounts for the reversed third-factor indices and includes zero dimensions. Equality of exponents is proved independently of the nine-quarters bound; the exact public target mme\_omega\_le\_225 then follows from our main theorem. No external solution is imported.

\begin{proof}
\proves{aux:prove2me-statement}
\leanok
\uses{aux:tensor-calculus,aux:matrix-multiplication-calculus,aux:asymptotic-rank-limit,aux:omega-bounds,eq:omega-log-n,thm:omega-le-nine-fourths}
\textit{Checked proof.} Take coordinates in the tensor-product standard basis, reversing the third pair of indices. Pure-tensor decompositions and diagonal restrictions have the same lengths, so the two ordinary ranks agree. Asymptotic rank gives $\omega_K\le\log R(\mathrm{MM}_n)/\log n$ for $n\ge2$. Along $n=2^k$, these ratios converge to $\omega_K$, proving equality of the infima; $\omega_K\le3$ handles the two small sizes. Rewrite the downloaded target using this equality and apply the manuscript main theorem.
\end{proof}

\renewcommand{\thetheorem}{A81}
\begin{lemma}[Rational restrictions and border rank]
\label{aux:rational-specialization-rank}
\manuscriptproof
\lean{MatrixMultiplication.Tensor.degenerates_of_restricts_liftRational,MatrixMultiplication.Tensor.borderRank_le_rank_liftRational}
\leanok
A restriction of scalar-extended tensors over $F(\varepsilon)$ gives a degeneration over $F$. Consequently $\underline R_F(T)\le R_{F(\varepsilon)}(T)$.
\end{lemma}

\noindent\textit{Auxiliary foundation. Source:} \texttt{material/main.tex:477}.
\textit{Status:} \texttt{formalized}. Additional manuscript coverage.
\textit{Proof supplied in manuscript.}
\textit{Checked Lean declarations:} \leandecl{MatrixMultiplication.\allowbreak{}Tensor.\allowbreak{}degenerates\_\allowbreak{}of\_\allowbreak{}restricts\_\allowbreak{}liftRational}, \leandecl{MatrixMultiplication.\allowbreak{}Tensor.\allowbreak{}borderRank\_\allowbreak{}le\_\allowbreak{}rank\_\allowbreak{}liftRational}.

\begin{proof}
\proves{aux:rational-specialization-rank}
\leanok
\uses{aux:rational-degeneration-equivalence,aux:border-rank-bound}
\textit{Checked proof.} The factor-map output is the constant rational tensor, hence regular at zero. Apply rational-to-polynomial degeneration equivalence to a minimal diagonal restriction witness.
\end{proof}

\renewcommand{\thetheorem}{A82}
\begin{lemma}[Asymptotic Gamma coefficient]
\label{aux:gamma-asymptotic}
\manuscriptproof
\lean{MatrixMultiplication.gamma_ratio_one_third_asymptotic,MatrixMultiplication.gammaGrowthBound_div_rpow_tendsto}
\leanok
$\Gamma(a+1/3)/(\Gamma(a)a^{1/3})\to1$ and $G(a)/a^{4/3}\to3/(2\Gamma(4/3))$.
\end{lemma}

\noindent\textit{Auxiliary foundation. Source:} \texttt{material/main.tex:1176}.
\textit{Status:} \texttt{formalized}. Additional manuscript coverage.
\textit{Proof supplied in manuscript.}
\textit{Checked Lean declarations:} \leandecl{MatrixMultiplication.\allowbreak{}gamma\_\allowbreak{}ratio\_\allowbreak{}one\_\allowbreak{}third\_\allowbreak{}asymptotic}, \leandecl{MatrixMultiplication.\allowbreak{}gammaGrowthBound\_\allowbreak{}div\_\allowbreak{}rpow\_\allowbreak{}tendsto}.
\textit{Formalization note:} The manuscript supplies a Stirling derivation; the formal proof uses Euler’s limit formula.

\begin{proof}
\proves{aux:gamma-asymptotic}
\leanok
\uses{lem:4/3}
\textit{Checked proof.} Euler’s Gamma limit gives the Gamma quotient limit directly. Multiplication by the elementary prefactor yields the coefficient.
\end{proof}

\renewcommand{\thetheorem}{A83}
\begin{lemma}[Certified Gamma coefficient enclosure]
\label{aux:gamma-coefficient-enclosure}
\lean{MatrixMultiplication.gamma_growth_coefficient_bounds,MatrixMultiplication.gammaThirdProduct,MatrixMultiplication.gammaThird_bracket}
\leanok
$1.67976<3/(2\Gamma(4/3))<1.67977$.
\end{lemma}

\noindent\textit{Auxiliary foundation. Source:} \texttt{material/main.tex:1179}.
\textit{Status:} \texttt{formalized}. Additional manuscript coverage.
\textit{Checked Lean declarations:} \leandecl{MatrixMultiplication.\allowbreak{}gamma\_\allowbreak{}growth\_\allowbreak{}coefficient\_\allowbreak{}bounds}, \leandecl{MatrixMultiplication.\allowbreak{}gammaThirdProduct}, \leandecl{MatrixMultiplication.\allowbreak{}gammaThird\_\allowbreak{}bracket}.
\textit{Formalization note:} The decimal assertion is stated without a manuscript derivation. The library supplies exact rational certificates.

\begin{proof}
\proves{aux:gamma-coefficient-enclosure}
\leanok
\textit{Checked proof.} Bound $\Gamma(4/3)^{-3}$ between monotone rational product sequences and evaluate at $n=320$. Compare the cubes of the decimal endpoints using exact arithmetic.
\end{proof}

\renewcommand{\thetheorem}{A84}
\begin{remark}[Arithmetic-operation interpretation of the exponent]
\label{aux:arithmetic-exponent-interpretation}
The introduction describes the exponent through asymptotic arithmetic-operation counts. The formal theorem uses the tensor-rank exponent defined in the preliminaries.
\end{remark}

\noindent\textit{Auxiliary foundation. Source:} \texttt{material/main.tex:166}.
\textit{Status:} \texttt{expository}. Additional manuscript coverage.
No independent Lean declaration is assigned to this expository note.
\textit{Formalization note:} The equality of the tensor-rank exponent with an arithmetic-operation exponent is conventional background, not formalized here. ExponentInfimum and the Prove2Me bridge compare tensor-rank definitions only; the project contains no computational model.

\section{Manuscript correspondence}

\renewcommand{\thetheorem}{2.1}
\begin{equationresult}[Factor-map definition of restriction]
\label{eq:restriction}
\lean{MatrixMultiplication.Tensor.Restricts,MatrixMultiplication.Tensor.map}
\leanok
$T\ge S$ means that $(g_1\otimes g_2\otimes g_3)T=S$ for three linear factor maps.
\end{equationresult}

\noindent\textit{Paper equation 2.1. Source:} \texttt{material/main.tex:299}.
\textit{Status:} \texttt{formalized}. Additional manuscript coverage.
\textit{Checked Lean declarations:} \leandecl{MatrixMultiplication.\allowbreak{}Tensor.\allowbreak{}Restricts}, \leandecl{MatrixMultiplication.\allowbreak{}Tensor.\allowbreak{}map}.

\begin{proof}
\proves{eq:restriction}
\leanok
\uses{aux:tensor-calculus}
\textit{Checked proof.} Use the coefficient-array factor-map definition.
\end{proof}

\renewcommand{\thetheorem}{2.2}
\begin{equationresult}[Polynomial-witness form of degeneration]
\label{eq:degen-def}
\lean{MatrixMultiplication.Tensor.DegenerationWitness,MatrixMultiplication.Tensor.rationalDegenerates_iff_degenerates}
\leanok
Rational specialization is equivalent to polynomial factor maps satisfying $(g_1(\varepsilon)\otimes g_2(\varepsilon)\otimes g_3(\varepsilon))T=\varepsilon^dS+\varepsilon^{d+1}U(\varepsilon)$ for $d\ge0$.
\end{equationresult}

\noindent\textit{Paper equation 2.2. Source:} \texttt{material/main.tex:312}.
\textit{Status:} \texttt{formalized}. Additional manuscript coverage.
\textit{Checked Lean declarations:} \leandecl{MatrixMultiplication.\allowbreak{}Tensor.\allowbreak{}DegenerationWitness}, \leandecl{MatrixMultiplication.\allowbreak{}Tensor.\allowbreak{}rationalDegenerates\_\allowbreak{}iff\_\allowbreak{}degenerates}.

\begin{proof}
\proves{eq:degen-def}
\leanok
\uses{aux:rational-degeneration-equivalence}
\textit{Checked proof.} Clear rational denominators, factor their lowest power of the parameter, and normalize the remaining nonzero constant coefficient. This also treats a zero target. Conversely divide out the polynomial witness order.
\end{proof}

\renewcommand{\thetheorem}{2.1}
\begin{definition}[Diagonal tensors and rank invariants]
\label{def:rank-invariants}
\lean{MatrixMultiplication.Tensor.diagonal,MatrixMultiplication.Tensor.rank,MatrixMultiplication.Tensor.borderRank,MatrixMultiplication.Tensor.asymptoticRank,MatrixMultiplication.Tensor.subrank,MatrixMultiplication.Tensor.borderSubrank,MatrixMultiplication.Tensor.asymptoticSubrank}
\leanok
\uses{aux:coefficient-tensors,aux:explicit-tensor-definitions,aux:degeneration-calculus,aux:ordinary-rank-definition,aux:asymptotic-rank-definition,aux:border-rank-bound,aux:subrank-invariants}
Define $\langle r\rangle=\sum_{i=1}^r e_i\otimes e_i\otimes e_i$. Rank and border rank are the least $r$ such that $\langle r\rangle$ restricts or degenerates to $T$. Subrank and border subrank reverse these relations. Define $\widetilde R(T)=\inf_{n\ge1}R(T^{\otimes n})^{1/n}$ and $\widetilde Q(T)=\sup_{n\ge1}Q(T^{\otimes n})^{1/n}$.
\end{definition}

\noindent\textit{Paper definition 2.1. Source:} \texttt{material/main.tex:324}.
\textit{Status:} \texttt{formalized}. Additional manuscript coverage.
\textit{Checked Lean declarations:} \leandecl{MatrixMultiplication.\allowbreak{}Tensor.\allowbreak{}diagonal}, \leandecl{MatrixMultiplication.\allowbreak{}Tensor.\allowbreak{}rank}, \leandecl{MatrixMultiplication.\allowbreak{}Tensor.\allowbreak{}borderRank}, \leandecl{MatrixMultiplication.\allowbreak{}Tensor.\allowbreak{}asymptoticRank}, \leandecl{MatrixMultiplication.\allowbreak{}Tensor.\allowbreak{}subrank}, \leandecl{MatrixMultiplication.\allowbreak{}Tensor.\allowbreak{}borderSubrank}, \leandecl{MatrixMultiplication.\allowbreak{}Tensor.\allowbreak{}asymptoticSubrank}.
\textit{Formalization note:} All seven definitions are checked. Subrank and border subrank are attained finite maxima; the positive-power subrank roots form a nonempty bounded set, so their real supremum agrees with the manuscript definition. Polynomial and rational degeneration are equivalent over every field.

\renewcommand{\thetheorem}{2.3}
\begin{equationresult}[Power multiplicativity of asymptotic rank]
\label{eq:monomial}
\manuscriptproof
\lean{MatrixMultiplication.Tensor.asymptoticRank_pow}
\leanok
For $N\ge1$, $\widetilde R(T^{\otimes N})=\widetilde R(T)^N$.
\end{equationresult}

\noindent\textit{Paper equation 2.3. Source:} \texttt{material/main.tex:354}.
\textit{Status:} \texttt{formalized}. Additional manuscript coverage.
\textit{Proof supplied in manuscript.}
\textit{Checked Lean declarations:} \leandecl{MatrixMultiplication.\allowbreak{}Tensor.\allowbreak{}asymptoticRank\_\allowbreak{}pow}.

\begin{proof}
\proves{eq:monomial}
\leanok
\uses{aux:asymptotic-rank-powers-products}
\textit{Checked proof.} Compare the rank-root limit along multiples of the exponent.
\end{proof}

\renewcommand{\thetheorem}{2.2}
\begin{lemma}[Asymptotic rank under degeneration]
\label{lem:asymp-border-ineq}
\manuscriptproof
\lean{MatrixMultiplication.Tensor.asymptoticRank_mono_degeneration,MatrixMultiplication.Tensor.asymptoticRank_le_borderRank}
\leanok
If $S\trianglelefteq T$, then $\widetilde R(S)\le\widetilde R(T)$. In particular, $\widetilde R(T)\le\underline R(T)$.
\end{lemma}

\noindent\textit{Paper lemma 2.2. Source:} \texttt{material/main.tex:361}.
\textit{Status:} \texttt{formalized}. Main proof path.
\textit{Proof supplied in manuscript.}
\textit{Checked Lean declarations:} \leandecl{MatrixMultiplication.\allowbreak{}Tensor.\allowbreak{}asymptoticRank\_\allowbreak{}mono\_\allowbreak{}degeneration}, \leandecl{MatrixMultiplication.\allowbreak{}Tensor.\allowbreak{}asymptoticRank\_\allowbreak{}le\_\allowbreak{}borderRank}.
\textit{Formalization note:} The exact binomial coefficient bound is checked separately. The existing limit proof retains its coarser quadratic overhead, which is sufficient to prove the same conclusion over every field.

\begin{proof}
\proves{lem:asymp-border-ineq}
\leanok
\uses{aux:asymptotic-rank-degeneration,aux:border-rank-bound}
\textit{Checked proof.} Extract degree $dN$ coefficients from the tensor-power polynomial witness. There are $\binom{dN+2}{2}$ restrictions. Bound their ranks, take $N$-th roots, and remove the polynomial overhead. Apply the result to a minimal border-rank witness.
\end{proof}

\renewcommand{\thetheorem}{2.4}
\begin{equationresult}[The leading-coefficient witness used in rank comparison]
\label{eq:a-degen}
\manuscriptproof
\lean{MatrixMultiplication.Tensor.DegenerationWitness,MatrixMultiplication.Tensor.Degenerates}
\leanok
A degeneration admits polynomial maps with $(g_1\otimes g_2\otimes g_3)T=\varepsilon^dS+\varepsilon^{d+1}U(\varepsilon)$.
\end{equationresult}

\noindent\textit{Paper equation 2.4. Source:} \texttt{material/main.tex:367}.
\textit{Status:} \texttt{formalized}. Additional manuscript coverage.
\textit{Proof supplied in manuscript.}
\textit{Checked Lean declarations:} \leandecl{MatrixMultiplication.\allowbreak{}Tensor.\allowbreak{}DegenerationWitness}, \leandecl{MatrixMultiplication.\allowbreak{}Tensor.\allowbreak{}Degenerates}.

\begin{proof}
\proves{eq:a-degen}
\leanok
\uses{aux:polynomial-degeneration-definition}
\textit{Checked proof.} Choose a polynomial degeneration witness.
\end{proof}

\renewcommand{\thetheorem}{2.3}
\begin{lemma}[Flattening lower bounds]
\label{lem:flattening-lb}
\manuscriptproof
\lean{MatrixMultiplication.Tensor.flatteningRank_le_asymptoticRank}
\leanok
For every tensor $T$ and each factor $i$, $R^{(i)}(T)\le\widetilde R(T)$.
\end{lemma}

\noindent\textit{Paper lemma 2.3. Source:} \texttt{material/main.tex:393}.
\textit{Status:} \texttt{formalized}. Main proof path.
\textit{Proof supplied in manuscript.}
\textit{Checked Lean declarations:} \leandecl{MatrixMultiplication.\allowbreak{}Tensor.\allowbreak{}flatteningRank\_\allowbreak{}le\_\allowbreak{}asymptoticRank}.
\textit{Formalization note:} The library also retains the stronger border-rank lower bound as auxiliary coverage.

\begin{proof}
\proves{lem:flattening-lb}
\leanok
\uses{aux:flattening-operations,aux:flattening-powers,aux:asymptotic-rank-definition}
\textit{Checked proof.} Tensor-power multiplicativity and the ordinary-rank bound give $R^{(i)}(T)^N\le R(T^{\otimes N})$. Take positive-power roots and then their infimum.
\end{proof}

\renewcommand{\thetheorem}{2.4}
\begin{definition}[Matrix multiplication and its exponent]
\label{def:matrix-multiplication}
\lean{MatrixMultiplication.Tensor.matrixMul,MatrixMultiplication.omega,MatrixMultiplication.Tensor.matrixMul_n_one_one_equivalent,MatrixMultiplication.Tensor.matrixMul_one_n_one_equivalent,MatrixMultiplication.Tensor.matrixMul_one_one_n_equivalent}
\leanok
\uses{aux:explicit-tensor-definitions,aux:asymptotic-rank-definition,aux:matrix-multiplication-exponent}
For positive integers $\ell,m,n$, define $\langle\ell,m,n\rangle=\sum_{i,j,k}e_{ij}\otimes e_{jk}\otimes e_{ik}$. For every field $F$, set $\omega_F=\log_2\widetilde R_F(\langle2,2,2\rangle)$. The special shapes have the coordinate presentations $\langle n,1,1\rangle\cong\sum_i e_i\otimes e_1\otimes e_i$, $\langle1,n,1\rangle\cong\sum_i e_i\otimes e_i\otimes e_1$, and $\langle1,1,n\rangle\cong\sum_i e_1\otimes e_i\otimes e_i$.
\end{definition}

\noindent\textit{Paper definition 2.4. Source:} \texttt{material/main.tex:404}.
\textit{Status:} \texttt{formalized}. Additional manuscript coverage.
\textit{Checked Lean declarations:} \leandecl{MatrixMultiplication.\allowbreak{}Tensor.\allowbreak{}matrixMul}, \leandecl{MatrixMultiplication.\allowbreak{}omega}, \leandecl{MatrixMultiplication.\allowbreak{}Tensor.\allowbreak{}matrixMul\_\allowbreak{}n\_\allowbreak{}one\_\allowbreak{}one\_\allowbreak{}equivalent}, \leandecl{MatrixMultiplication.\allowbreak{}Tensor.\allowbreak{}matrixMul\_\allowbreak{}one\_\allowbreak{}n\_\allowbreak{}one\_\allowbreak{}equivalent}, \leandecl{MatrixMultiplication.\allowbreak{}Tensor.\allowbreak{}matrixMul\_\allowbreak{}one\_\allowbreak{}one\_\allowbreak{}n\_\allowbreak{}equivalent}.
\textit{Formalization note:} The coordinate formula, exponent definition, and all three special-shape equivalences are checked. The coordinate equivalences also hold at size zero and over any commutative semiring.

\begin{proof}
\proves{def:matrix-multiplication}
\leanok
\uses{aux:coefficient-tensors}
\textit{Checked proof.} Use the checked coordinate formula and exponent definition. Construct the three special-shape equivalences by removing singleton coordinates in the factor index types and applying coordinate reindexing.
\end{proof}

\renewcommand{\thetheorem}{2.5}
\begin{equationresult}[Matrix-size comparison by floor and ceiling powers]
\label{eq:n-matmul-to-2}
\manuscriptproof
\lean{MatrixMultiplication.Tensor.matrixMul_kroneckerPower_log_sandwich}
\leanok
For $m\ge1$, $n\ge2$, $\ell=\lfloor m\log_2n\rfloor$, and $k=\lceil m\log_2n\rceil$, $\langle2,2,2\rangle^{\otimes\ell}\le\langle n,n,n\rangle^{\otimes m}\cong\langle n^m,n^m,n^m\rangle\le\langle2,2,2\rangle^{\otimes k}$.
\end{equationresult}

\noindent\textit{Paper equation 2.5. Source:} \texttt{material/main.tex:433}.
\textit{Status:} \texttt{formalized}. Additional manuscript coverage.
\textit{Proof supplied in manuscript.}
\textit{Checked Lean declarations:} \leandecl{MatrixMultiplication.\allowbreak{}Tensor.\allowbreak{}matrixMul\_\allowbreak{}kroneckerPower\_\allowbreak{}log\_\allowbreak{}sandwich}.

\begin{proof}
\proves{eq:n-matmul-to-2}
\leanok
\uses{aux:matrix-multiplication-calculus}
\textit{Checked proof.} Bound the integer matrix sizes using floor and ceiling logarithms; apply dimension restrictions and tensor-power equivalences.
\end{proof}

\renewcommand{\thetheorem}{2.6}
\begin{equationresult}[Matrix size independence of the exponent]
\label{eq:omega-log-n}
\manuscriptproof
\lean{MatrixMultiplication.Tensor.asymptoticRank_matrixMul_square}
\leanok
\uses{aux:matrix-multiplication-exponent}
For every field $F$ and every integer $n\ge1$, $\widetilde R_F(\langle n,n,n\rangle)=n^{\omega_F}$.
\end{equationresult}

\noindent\textit{Paper equation 2.6. Source:} \texttt{material/main.tex:438}.
\textit{Status:} \texttt{formalized}. Main proof path.
\textit{Proof supplied in manuscript.}
\textit{Checked Lean declarations:} \leandecl{MatrixMultiplication.\allowbreak{}Tensor.\allowbreak{}asymptoticRank\_\allowbreak{}matrixMul\_\allowbreak{}square}.
\textit{Formalization note:} Checked in the power formulation over every field for every natural n with 1 <= n, including scalar multiplication. The manuscript leaves positivity implicit in Definition 2.4; Lean retains the explicit hypothesis. This proof needs neither numerical bounds on omega nor the rational/polynomial degeneration comparison.

\begin{proof}
\proves{eq:omega-log-n}
\leanok
\uses{aux:matrix-multiplication-calculus,aux:asymptotic-rank-powers-products,aux:monotone-power-functions}
\textit{Checked proof.} Apply the monotone-power-function theorem to $f(n)=\widetilde R(\langle n,n,n\rangle)$. Dimension restrictions give monotonicity, scalar multiplication gives $f(1)=1$, and the tensor-power coordinate equivalence together with asymptotic-rank power multiplicativity gives $f(n^k)=f(n)^k$.
\end{proof}

\renewcommand{\thetheorem}{2.5}
\begin{definition}[Polynomial multiplication tensors]
\label{def:polynomial-multiplication}
\lean{MatrixMultiplication.Tensor.polynomialMul,MatrixMultiplication.Tensor.coefficientPolynomial,MatrixMultiplication.Tensor.polynomialMul_bilinear}
\leanok
\uses{aux:explicit-tensor-definitions}
For positive integers $a,b$, define $\mathrm{PM}_{a,b}=\sum_{i=0}^{a-1}\sum_{j=0}^{b-1}e_i\otimes e_j\otimes e_{i+j}$. This is the structure tensor of multiplication of polynomials of degree less than $a$ and $b$.
\end{definition}

\noindent\textit{Paper definition 2.5. Source:} \texttt{material/main.tex:446}.
\textit{Status:} \texttt{formalized}. Main proof path.
\textit{Checked Lean declarations:} \leandecl{MatrixMultiplication.\allowbreak{}Tensor.\allowbreak{}polynomialMul}, \leandecl{MatrixMultiplication.\allowbreak{}Tensor.\allowbreak{}coefficientPolynomial}, \leandecl{MatrixMultiplication.\allowbreak{}Tensor.\allowbreak{}polynomialMul\_\allowbreak{}bilinear}.
\textit{Formalization note:} The coefficient-vector expansion is checked against ordinary polynomial multiplication over any commutative semiring.

\begin{proof}
\proves{def:polynomial-multiplication}
\leanok
\textit{Checked proof.} Expand the two finite sums of monomials, multiply them, and extract the coefficient at each output index.
\end{proof}

\renewcommand{\thetheorem}{2.6}
\begin{lemma}[Ranks of polynomial multiplication]
\label{lem:poly-mult-from-sum}
\manuscriptproof
\lean{MatrixMultiplication.Tensor.rank_polynomialMul_of_card,MatrixMultiplication.Tensor.rank_polynomialMul_finiteField_le,MatrixMultiplication.Tensor.diagonal_degenerates_polynomialMul,MatrixMultiplication.Tensor.borderRank_polynomialMul,MatrixMultiplication.Tensor.asymptoticRank_polynomialMul,MatrixMultiplication.Tensor.rank_polynomialMul_of_nodes_atInfinity,MatrixMultiplication.Tensor.rank_polynomialMul_finiteField_le_clog,MatrixMultiplication.Tensor.rank_polynomialMul_smallFiniteField_le}
\leanok
For positive $a,b$, a field with at least $a+b-2$ elements satisfies $R(\mathrm{PM}_{a,b})=a+b-1$. Over a finite field of size $q<a+b-2$, $R(\mathrm{PM}_{a,b})\le(a+b-1)\lceil\log_q(a+b-2)\rceil^2$. Over every field, $\underline R(\mathrm{PM}_{a,b})=\widetilde R(\mathrm{PM}_{a,b})=a+b-1$.
\end{lemma}

\noindent\textit{Paper lemma 2.6. Source:} \texttt{material/main.tex:454}.
\textit{Status:} \texttt{formalized}. Additional manuscript coverage.
\textit{Proof supplied in manuscript.}
\textit{Checked Lean declarations:} \leandecl{MatrixMultiplication.\allowbreak{}Tensor.\allowbreak{}rank\_\allowbreak{}polynomialMul\_\allowbreak{}of\_\allowbreak{}card}, \leandecl{MatrixMultiplication.\allowbreak{}Tensor.\allowbreak{}rank\_\allowbreak{}polynomialMul\_\allowbreak{}finiteField\_\allowbreak{}le}, \leandecl{MatrixMultiplication.\allowbreak{}Tensor.\allowbreak{}diagonal\_\allowbreak{}degenerates\_\allowbreak{}polynomialMul}, \leandecl{MatrixMultiplication.\allowbreak{}Tensor.\allowbreak{}borderRank\_\allowbreak{}polynomialMul}, \leandecl{MatrixMultiplication.\allowbreak{}Tensor.\allowbreak{}asymptoticRank\_\allowbreak{}polynomialMul}, \leandecl{MatrixMultiplication.\allowbreak{}Tensor.\allowbreak{}rank\_\allowbreak{}polynomialMul\_\allowbreak{}of\_\allowbreak{}nodes\_\allowbreak{}atInfinity}, \leandecl{MatrixMultiplication.\allowbreak{}Tensor.\allowbreak{}rank\_\allowbreak{}polynomialMul\_\allowbreak{}finiteField\_\allowbreak{}le\_\allowbreak{}clog}, \leandecl{MatrixMultiplication.\allowbreak{}Tensor.\allowbreak{}rank\_\allowbreak{}polynomialMul\_\allowbreak{}smallFiniteField\_\allowbreak{}le}.
\textit{Formalization note:} The small-field hypothesis makes the ceiling logarithm greater than one, so the maximum in the stronger uniform library bound can be removed. Interpolation includes the scalar case; the exact border and asymptotic ranks hold over every field.

\begin{proof}
\proves{lem:poly-mult-from-sum}
\leanok
\uses{def:polynomial-multiplication,eq:extensionscalars,lem:flattening-lb,aux:asymptotic-rank-calculus,aux:polynomial-multiplication-degeneration,aux:polynomial-multiplication-ranks,aux:rational-specialization-rank}
\textit{Checked proof.} Use finite evaluations and a leading-coefficient summand whose output polynomial is $X^{a+b-2}$ minus its finite-node interpolant. This gives $a+b-1$ pure summands at the sharp field-size threshold; the third flattening gives the matching lower bound. For finite fields choose the checked degree-$k$ extension with cardinality $q^k$, where $k=\lceil\log_q(a+b-2)\rceil>1$, and apply scalar descent. The existing uniform diagonal degeneration and flattening bound give the exact border and asymptotic ranks.
\end{proof}

\renewcommand{\thetheorem}{2.7}
\begin{equationresult}[Rank descent through a finite extension]
\label{eq:extensionscalars}
\lean{MatrixMultiplication.Tensor.rank_le_extension_degree_sq_mul,MatrixMultiplication.Tensor.scalarExtension,MatrixMultiplication.Tensor.scalarExtension_polynomialMul}
\leanok
If $E/F$ has finite degree $k$ and $T$ is defined over $F$, then $R_F(T)\le k^2R_E(T_E)$.
\end{equationresult}

\noindent\textit{Paper equation 2.7. Source:} \texttt{material/main.tex:468}.
\textit{Status:} \texttt{formalized}. Additional manuscript coverage.
\textit{Checked Lean declarations:} \leandecl{MatrixMultiplication.\allowbreak{}Tensor.\allowbreak{}rank\_\allowbreak{}le\_\allowbreak{}extension\_\allowbreak{}degree\_\allowbreak{}sq\_\allowbreak{}mul}, \leandecl{MatrixMultiplication.\allowbreak{}Tensor.\allowbreak{}scalarExtension}, \leandecl{MatrixMultiplication.\allowbreak{}Tensor.\allowbreak{}scalarExtension\_\allowbreak{}polynomialMul}.
\textit{Formalization note:} The finite-degree rank descent is checked for arbitrary field extensions and finite tensor coordinate types, including zero tensors. The proof uses an arbitrary basis of the extension and a linear retraction of the algebra map.

\begin{proof}
\proves{eq:extensionscalars}
\leanok
\uses{def:rank-invariants,aux:tensor-calculus}
\textit{Checked proof.} Choose any basis of the extension over the base field and a linear left inverse of the algebra map. Expand the first two factors of every extension-field pure summand in this basis and retract the remaining coefficient to the base field. Each summand gives at most $k^2$ pure summands. A basis containing one is unnecessary.
\end{proof}

\renewcommand{\thetheorem}{2.7}
\begin{lemma}[Equal-block asymptotic sum identity]
\label{cor:asymptotic-rank-lower-bound}
\manuscriptproof
\lean{MatrixMultiplication.Tensor.asymptoticRank_matrixMul_directPower}
\leanok
\uses{aux:direct-sum-copies,aux:matrix-multiplication-exponent}
For positive integers $a,b$, $\widetilde R(\langle b,b,b\rangle^{\oplus a})=ab^\omega$.
\end{lemma}

\noindent\textit{Paper lemma 2.7. Source:} \texttt{material/main.tex:492}.
\textit{Status:} \texttt{formalized}. Main proof path.
\textit{Proof supplied in manuscript.}
\textit{Checked Lean declarations:} \leandecl{MatrixMultiplication.\allowbreak{}Tensor.\allowbreak{}asymptoticRank\_\allowbreak{}matrixMul\_\allowbreak{}directPower}.
\textit{Formalization note:} Checked for every field K and all positive natural a,b, including a=1 and b=1. Repeated copies use explicit block-labelled coordinates. The revised manuscript explicitly supplies the integer rounding argument. The existing cor: label is preserved although the active environment is a lemma.

\begin{proof}
\proves{cor:asymptotic-rank-lower-bound}
\leanok
\uses{eq:omega-log-n,aux:matrix-multiplication-capacity,aux:integer-capacity-growth,aux:asymptotic-rank-limit,aux:omega-lower-bound,aux:flattening-powers}
\textit{Checked proof.} The repeated-copy upper bound and the square-size identity give $\widetilde R(\langle b,b,b\rangle^{\oplus a})\le ab^\omega$. For the lower bound, treat $a=1$ by selecting the sole block. For $a\ge2$, set $r_n=R(\langle2,2,2\rangle^{\otimes n})$, $q=\widetilde R(\langle2,2,2\rangle)>1$, and $p=\widetilde R(\langle b,b,b\rangle^{\oplus a})/b^\omega\ge1$. The matrix capacity comparison gives $q^n\le p^m$ whenever $r_n\le a^m$. Apply the integer-capacity growth lemma and the root-limit characterization to get $a\le p$.
\end{proof}

\renewcommand{\thetheorem}{2.8}
\begin{lemma}[Entropy variational identity]
\label{lem:optimize-sum}
\lean{MatrixMultiplication.entropy_variational_sum,MatrixMultiplication.entropy_variational_normalized,MatrixMultiplication.entropy_variational_attained,MatrixMultiplication.entropy_variational_sSup}
\leanok
For a nonempty finite set $I$ and $a_i\ge0$, $\max_q 2^{H(q)}\prod_i a_i^{q_i}=\sum_i a_i$, where $q$ ranges over probability distributions and $0^0=1$.
\end{lemma}

\noindent\textit{Paper lemma 2.8. Source:} \texttt{material/main.tex:522}.
\textit{Status:} \texttt{formalized}. Additional manuscript coverage.
\textit{Checked Lean declarations:} \leandecl{MatrixMultiplication.\allowbreak{}entropy\_\allowbreak{}variational\_\allowbreak{}sum}, \leandecl{MatrixMultiplication.\allowbreak{}entropy\_\allowbreak{}variational\_\allowbreak{}normalized}, \leandecl{MatrixMultiplication.\allowbreak{}entropy\_\allowbreak{}variational\_\allowbreak{}attained}, \leandecl{MatrixMultiplication.\allowbreak{}entropy\_\allowbreak{}variational\_\allowbreak{}sSup}.
\textit{Formalization note:} The full maximum is proved as an IsGreatest statement, including attainment. Normalized coordinates give a maximizer for positive total mass; a point mass handles zero total mass. Zero coordinates respect the convention 0\textasciicircum{}0=1. This result remains independent of the main theorem's direct multinomial proof.

\begin{proof}
\proves{lem:optimize-sum}
\leanok
\uses{aux:entropy}
\textit{Checked proof.} The checked weighted AM--GM inequality gives the upper bound. If the total mass is positive, set $q_i=a_i/\sum_j a_j$ and use the entropy product identity to obtain equality, treating zero coordinates explicitly. If the total mass is zero, choose a point mass; the upper bound and nonnegativity give equality. Record attainment and the upper bound together as a maximum.
\end{proof}

\renewcommand{\thetheorem}{3.1}
\begin{definition}[Partially direct sums]
\label{def:partial-direct-sum}
\lean{MatrixMultiplication.Tensor.partialSumFamily,MatrixMultiplication.Tensor.partialSum}
\leanok
\uses{aux:coefficient-tensors}
For $T_h\in U\otimes V_h\otimes W_h$, define $\mathop{\diamondplus}_hT_h$ using the canonical inclusions into $U\otimes(\bigoplus_hV_h)\otimes(\bigoplus_hW_h)$.
\end{definition}

\noindent\textit{Paper definition 3.1. Source:} \texttt{material/main.tex:534}.
\textit{Status:} \texttt{formalized}. Main proof path.
\textit{Checked Lean declarations:} \leandecl{MatrixMultiplication.\allowbreak{}Tensor.\allowbreak{}partialSumFamily}, \leandecl{MatrixMultiplication.\allowbreak{}Tensor.\allowbreak{}partialSum}.
\textit{Formalization note:} Implemented for every finite family with a common first coordinate space and dependent second/third coordinate spaces; the empty family is included.

\begin{proof}
\proves{def:partial-direct-sum}
\leanok
\textit{Checked proof.} Use tagged basis indices in the second and third factors. The coefficient is the corresponding block coefficient when the tags match and zero otherwise.
\end{proof}

\renewcommand{\thetheorem}{3.2}
\begin{lemma}[Partial-sum calculus]
\label{diamondplus-basic}
\manuscriptproof
\lean{MatrixMultiplication.Tensor.partialSum_restricts_summand,MatrixMultiplication.Tensor.partialSum_restricts_left,MatrixMultiplication.Tensor.partialSum_restricts_right,MatrixMultiplication.Tensor.partialSum_comm,MatrixMultiplication.Tensor.partialSum_assoc,MatrixMultiplication.Tensor.kronecker_partialSum,MatrixMultiplication.Tensor.partialSum_kronecker,MatrixMultiplication.Tensor.partialSumCopies,MatrixMultiplication.Tensor.partialPower_equiv_matrixMul,MatrixMultiplication.Tensor.partialPower_unit_equiv_matrixMul}
\leanok
A partially direct sum restricts to each summand. Partial sums are commutative and associative up to the canonical equivalences, and Kronecker products distribute over them on either side. Moreover, $T^{\diamondplus r}\cong T\otimes\langle1,1,r\rangle$.
\end{lemma}

\noindent\textit{Paper lemma 3.2. Source:} \texttt{material/main.tex:540}.
\textit{Status:} \texttt{formalized}. Additional manuscript coverage.
\textit{Proof supplied in manuscript.}
\textit{Checked Lean declarations:} \leandecl{MatrixMultiplication.\allowbreak{}Tensor.\allowbreak{}partialSum\_\allowbreak{}restricts\_\allowbreak{}summand}, \leandecl{MatrixMultiplication.\allowbreak{}Tensor.\allowbreak{}partialSum\_\allowbreak{}restricts\_\allowbreak{}left}, \leandecl{MatrixMultiplication.\allowbreak{}Tensor.\allowbreak{}partialSum\_\allowbreak{}restricts\_\allowbreak{}right}, \leandecl{MatrixMultiplication.\allowbreak{}Tensor.\allowbreak{}partialSum\_\allowbreak{}comm}, \leandecl{MatrixMultiplication.\allowbreak{}Tensor.\allowbreak{}partialSum\_\allowbreak{}assoc}, \leandecl{MatrixMultiplication.\allowbreak{}Tensor.\allowbreak{}kronecker\_\allowbreak{}partialSum}, \leandecl{MatrixMultiplication.\allowbreak{}Tensor.\allowbreak{}partialSum\_\allowbreak{}kronecker}, \leandecl{MatrixMultiplication.\allowbreak{}Tensor.\allowbreak{}partialSumCopies}, \leandecl{MatrixMultiplication.\allowbreak{}Tensor.\allowbreak{}partialPower\_\allowbreak{}equiv\_\allowbreak{}matrixMul}, \leandecl{MatrixMultiplication.\allowbreak{}Tensor.\allowbreak{}partialPower\_\allowbreak{}unit\_\allowbreak{}equiv\_\allowbreak{}matrixMul}.
\textit{Formalization note:} All six items are checked over commutative semirings, including zero repeated copies. The common first-factor map is fixed in the partial-sum reindexings.

\begin{proof}
\proves{diamondplus-basic}
\leanok
\uses{def:partial-direct-sum,aux:explicit-tensor-definitions,aux:restriction-operations}
\textit{Checked proof.} Construct the coordinate inclusions, projections, and reindexings. Every identity must keep the common first-factor map consistent.
\end{proof}

\renewcommand{\thetheorem}{3.3}
\begin{remark}[The common-factor compatibility caveat]
\label{rem:partial-sum-equivalence}
\lean{MatrixMultiplication.Tensor.map_partialSum,MatrixMultiplication.Tensor.partialSum_restricts_fixedFirst,MatrixMultiplication.Tensor.partialSum_congr_fixedFirst,MatrixMultiplication.Tensor.partialSumCounterexampleBlock,MatrixMultiplication.Tensor.partialSumCounterexample_source_rank,MatrixMultiplication.Tensor.partialSumCounterexample_target_rank,MatrixMultiplication.Tensor.partialSum_equivalence_counterexample,MatrixMultiplication.Tensor.partialSum_restriction_counterexample}
\leanok
Separate equivalences $S_h\cong T_h$ do not in general induce an equivalence of their partially direct sums: they can use incompatible maps on the common factor. The analogous implication for restrictions also fails.
\end{remark}

\noindent\textit{Paper remark 3.3. Source:} \texttt{material/main.tex:555}.
\textit{Status:} \texttt{formalized}. Additional manuscript coverage.
\textit{Checked Lean declarations:} \leandecl{MatrixMultiplication.\allowbreak{}Tensor.\allowbreak{}map\_\allowbreak{}partialSum}, \leandecl{MatrixMultiplication.\allowbreak{}Tensor.\allowbreak{}partialSum\_\allowbreak{}restricts\_\allowbreak{}fixedFirst}, \leandecl{MatrixMultiplication.\allowbreak{}Tensor.\allowbreak{}partialSum\_\allowbreak{}congr\_\allowbreak{}fixedFirst}, \leandecl{MatrixMultiplication.\allowbreak{}Tensor.\allowbreak{}partialSumCounterexampleBlock}, \leandecl{MatrixMultiplication.\allowbreak{}Tensor.\allowbreak{}partialSumCounterexample\_\allowbreak{}source\_\allowbreak{}rank}, \leandecl{MatrixMultiplication.\allowbreak{}Tensor.\allowbreak{}partialSumCounterexample\_\allowbreak{}target\_\allowbreak{}rank}, \leandecl{MatrixMultiplication.\allowbreak{}Tensor.\allowbreak{}partialSum\_\allowbreak{}equivalence\_\allowbreak{}counterexample}, \leandecl{MatrixMultiplication.\allowbreak{}Tensor.\allowbreak{}partialSum\_\allowbreak{}restriction\_\allowbreak{}counterexample}.
\textit{Formalization note:} Both negative assertions are checked over every field. The same first-factor map is sufficient to assemble restrictions; common maps in both directions suffice for equivalences. The explicit counterexample uses scalar blocks supported at one shared basis vector versus two distinct basis vectors.

\begin{proof}
\proves{rem:partial-sum-equivalence}
\leanok
\uses{def:partial-direct-sum,aux:flattening-calculus}
\textit{Checked proof.} Use a common first-factor matrix and block-diagonal matrices on the other factors to assemble compatible maps. For the counterexample, each block is a scalar tensor supported at $e_0$ or $e_1$. The partial sum with both blocks at $e_0$ has first-flattening rank $1$, while the sum with blocks at $e_0,e_1$ has rank $2$. Flattening monotonicity excludes a restriction, and hence an equivalence, although the individual blocks are equivalent.
\end{proof}

\renewcommand{\thetheorem}{3.4}
\begin{lemma}[Separation by polynomial multiplication]
\label{lem:diamond-plus-deg}
\manuscriptproof
\lean{MatrixMultiplication.Tensor.partialSumFamily_kronecker_polynomialMul_degenerates,MatrixMultiplication.Tensor.partialSumFamily_separation}
\leanok
Let $T=\mathop{\diamondplus}_{h=1}^mT_h$ with $m,r\ge1$. Then $T^{\oplus(m+r-1)}\trianglerighteq(\bigoplus_{h=1}^mT_h)\otimes\langle1,1,r\rangle$.
\end{lemma}

\noindent\textit{Paper lemma 3.4. Source:} \texttt{material/main.tex:565}.
\textit{Status:} \texttt{formalized}. Main proof path.
\textit{Proof supplied in manuscript.}
\textit{Checked Lean declarations:} \leandecl{MatrixMultiplication.\allowbreak{}Tensor.\allowbreak{}partialSumFamily\_\allowbreak{}kronecker\_\allowbreak{}polynomialMul\_\allowbreak{}degenerates}, \leandecl{MatrixMultiplication.\allowbreak{}Tensor.\allowbreak{}partialSumFamily\_\allowbreak{}separation}.
\textit{Formalization note:} The full separation statement is checked over every field. The implementation also handles empty families and zero auxiliary size; the intermediate weighted degeneration holds over commutative semirings.

\begin{proof}
\proves{lem:diamond-plus-deg}
\leanok
\uses{def:partial-direct-sum,aux:polynomial-multiplication-degeneration,aux:monomial-degeneration,aux:polynomial-degeneration-operations}
\textit{Checked proof.} Degenerate the diagonal auxiliary tensor to polynomial multiplication. With zero-based indices use weights $i^2$, $h^2+2hj$, and $D-2hk$, choosing $D$ large enough to keep every weight nonnegative. Supported terms have total weight $D+(i-h)^2$, so the leading terms separate the blocks. Coordinate pullbacks identify the target with an ordinary direct sum tensored with $\langle1,1,r\rangle$.
\end{proof}

\renewcommand{\thetheorem}{3.5}
\begin{lemma}[Separation of a multinomial type class]
\label{lem:decoupling-asymptotic}
\manuscriptproof
\lean{MatrixMultiplication.Tensor.partialSum_type_separation}
\leanok
For nonnegative counts $N_h$ with sum $N$ and $Q=\binom{N}{N_1,\ldots,N_m}$, $(T^{\otimes N})^{\oplus(2Q-1)}\trianglerighteq(\bigotimes_hT_h^{\otimes N_h})^{\oplus Q}\otimes\langle1,1,Q\rangle$.
\end{lemma}

\noindent\textit{Paper lemma 3.5. Source:} \texttt{material/main.tex:626}.
\textit{Status:} \texttt{formalized}. Main proof path.
\textit{Proof supplied in manuscript.}
\textit{Checked Lean declarations:} \leandecl{MatrixMultiplication.\allowbreak{}Tensor.\allowbreak{}partialSum\_\allowbreak{}type\_\allowbreak{}separation}.
\textit{Formalization note:} Checked with the exact multinomial coefficient and exact factor 2Q-1. Zero counts, N=0, and empty families are included. Independent first-factor equivalences are used only after ordinary separation.

\begin{proof}
\proves{lem:decoupling-asymptotic}
\leanok
\uses{lem:diamond-plus-deg,aux:word-type-cardinality,aux:partial-sum-word-selection,aux:word-product-grouping,aux:ordinary-sum-block-transport}
\textit{Checked proof.} Expand the tensor power as a partial sum indexed by words, restrict to the specified type class, and apply separation with $m=r=Q$. Only after obtaining an ordinary direct sum identify the equivalent blocks by permuting tensor-power positions. Include the zero-count and $N=0$ unit conventions.
\end{proof}

\renewcommand{\thetheorem}{4.1}
\begin{definition}[Kronecker symmetrization]
\label{def:symmetrization}
\lean{MatrixMultiplication.Tensor.symmetrize}
\leanok
\uses{aux:coefficient-tensors,aux:factor-permutations}
Define $T^{\otimes\mathfrak S_3}=\bigotimes_{\pi\in\mathfrak S_3}\pi\cdot T$, using all six permutations of the three tensor factors.
\end{definition}

\noindent\textit{Paper definition 4.1. Source:} \texttt{material/main.tex:662}.
\textit{Status:} \texttt{formalized}. Main proof path.
\textit{Checked Lean declarations:} \leandecl{MatrixMultiplication.\allowbreak{}Tensor.\allowbreak{}symmetrize}.
\textit{Formalization note:} Checked as the product of the three even and three odd factor permutations, each occurring exactly once. The explicit six orders are identified bijectively with S3.

\renewcommand{\thetheorem}{4.2}
\begin{lemma}[Symmetrization identities]
\label{lem:sym-properties}
\manuscriptproof
\lean{MatrixMultiplication.Tensor.symmetrize_permuteS3,MatrixMultiplication.Tensor.symmetrize_kronecker_equivalent,MatrixMultiplication.Tensor.symmetrize_directSumCopies_equivalent,MatrixMultiplication.Tensor.symmetrize_diagonal_equivalent,MatrixMultiplication.Tensor.symmetrize_matrixMul}
\leanok
\uses{def:symmetrization}
Symmetrization is invariant under permuting tensor factors and commutes with Kronecker product up to equivalence. It takes $T^{\oplus r}$ to $(T^{\otimes\mathfrak S_3})^{\oplus r^6}$ and takes $\langle\ell,m,n\rangle$ to $\langle(\ell mn)^2,(\ell mn)^2,(\ell mn)^2\rangle$.
\end{lemma}

\noindent\textit{Paper lemma 4.2. Source:} \texttt{material/main.tex:667}.
\textit{Status:} \texttt{formalized}. Main proof path.
\textit{Proof supplied in manuscript.}
\textit{Checked Lean declarations:} \leandecl{MatrixMultiplication.\allowbreak{}Tensor.\allowbreak{}symmetrize\_\allowbreak{}permuteS3}, \leandecl{MatrixMultiplication.\allowbreak{}Tensor.\allowbreak{}symmetrize\_\allowbreak{}kronecker\_\allowbreak{}equivalent}, \leandecl{MatrixMultiplication.\allowbreak{}Tensor.\allowbreak{}symmetrize\_\allowbreak{}directSumCopies\_\allowbreak{}equivalent}, \leandecl{MatrixMultiplication.\allowbreak{}Tensor.\allowbreak{}symmetrize\_\allowbreak{}diagonal\_\allowbreak{}equivalent}, \leandecl{MatrixMultiplication.\allowbreak{}Tensor.\allowbreak{}symmetrize\_\allowbreak{}matrixMul}.
\textit{Formalization note:} All four properties are checked over commutative semirings, including zero dimensions and zero copies. Permutation invariance quantifies over the actual symmetric group; diagonal and matrix identities use explicit coordinate equivalences.

\begin{proof}
\proves{lem:sym-properties}
\leanok
\uses{aux:factor-permutations,aux:direct-sum-copies,aux:matrix-multiplication-calculus}
\textit{Checked proof.} Construct finite index equivalences for the six factors. Combine the coordinate formulas for direct sums and matrix multiplication.
\end{proof}

\renewcommand{\thetheorem}{4.3}
\begin{definition}[Fractional symmetrized value]
\label{def:fractional-value}
\lean{MatrixMultiplication.Tensor.ValueWitness,MatrixMultiplication.Tensor.fractionalValue,MatrixMultiplication.Tensor.valueWitnessEquivRational,MatrixMultiplication.Tensor.fractionalValue_eq_rationalFractionalValue}
\leanok
\uses{def:symmetrization,aux:matrix-multiplication-exponent,aux:direct-sum-copies}
For $\tau>0$, let $V_\tau(T)$ be the supremum in $[0,+\infty]$ of $\bigl(ab^{3\tau}/(cd^{3\tau})\bigr)^{1/(6N)}$ over positive integers $a,b,c,d,N$ satisfying $\langle d,d,d\rangle^{\oplus c}\otimes(T^{\otimes\mathfrak S_3})^{\otimes N}\trianglerighteq\langle b,b,b\rangle^{\oplus a}$. For $T=0$ the witness set is empty, so $V_\tau(0)=0$.
\end{definition}

\noindent\textit{Paper definition 4.3. Source:} \texttt{material/main.tex:682}.
\textit{Status:} \texttt{formalized}. Main proof path.
\textit{Checked Lean declarations:} \leandecl{MatrixMultiplication.\allowbreak{}Tensor.\allowbreak{}ValueWitness}, \leandecl{MatrixMultiplication.\allowbreak{}Tensor.\allowbreak{}fractionalValue}, \leandecl{MatrixMultiplication.\allowbreak{}Tensor.\allowbreak{}valueWitnessEquivRational}, \leandecl{MatrixMultiplication.\allowbreak{}Tensor.\allowbreak{}fractionalValue\_\allowbreak{}eq\_\allowbreak{}rationalFractionalValue}.
\textit{Formalization note:} The value is defined directly as the witness supremum in the extended nonnegative reals, and agrees exactly with rational-function witnesses. The zero tensor has no witnesses, so its value follows from the empty supremum without a separate branch. Lean extends the formula to all real parameters; the paper uses positive parameters. The revised definition explicitly uses the extended supremum and derives the zero value from it.

\begin{proof}
\proves{def:fractional-value}
\leanok
\uses{aux:fractional-value-definition,aux:rational-value-equivalence}
\textit{Checked proof.} Transfer polynomial and rational-function witnesses in both directions without changing their parameters or scores, then identify the two indexed suprema.
\end{proof}

\renewcommand{\thetheorem}{4.4}
\begin{remark}[The exponent convention]
\label{rem:value-exponent-convention}
\uses{def:fractional-value,cor:value-diag-matrix}
The power $3\tau$ in the definition agrees with the convention that matrix multiplication has value $(\ell mn)^\tau$.
\end{remark}

\noindent\textit{Paper remark 4.4. Source:} \texttt{material/main.tex:691}.
\textit{Status:} \texttt{expository}. Additional manuscript coverage.
No independent Lean declaration is assigned to this expository note.
\textit{Formalization note:} Expository convention; no separate Lean theorem is required.

\renewcommand{\thetheorem}{4.5}
\begin{lemma}[Value outside the admissible interval]
\label{lem:value-outside-range}
\manuscriptproof
\lean{MatrixMultiplication.Tensor.fractionalValue_eq_top_of_outside_range,MatrixMultiplication.Tensor.fractionalValue_parameter_bounds_of_ne_top,MatrixMultiplication.Tensor.two_thirds_le_of_unit_ne_top,MatrixMultiplication.Tensor.fractionalValue_unit_eq_top_of_lt_two_thirds,MatrixMultiplication.Tensor.fractionalValue_ne_top_iff,MatrixMultiplication.two_le_of_nat_sq_le_double_rpow}
\leanok
For $T\ne0$ and $\tau>0$ outside $[2/3,\omega/3]$, $V_\tau(T)=+\infty$.
\end{lemma}

\noindent\textit{Paper lemma 4.5. Source:} \texttt{material/main.tex:697}.
\textit{Status:} \texttt{formalized}. Additional manuscript coverage.
\textit{Proof supplied in manuscript.}
\textit{Checked Lean declarations:} \leandecl{MatrixMultiplication.\allowbreak{}Tensor.\allowbreak{}fractionalValue\_\allowbreak{}eq\_\allowbreak{}top\_\allowbreak{}of\_\allowbreak{}outside\_\allowbreak{}range}, \leandecl{MatrixMultiplication.\allowbreak{}Tensor.\allowbreak{}fractionalValue\_\allowbreak{}parameter\_\allowbreak{}bounds\_\allowbreak{}of\_\allowbreak{}ne\_\allowbreak{}top}, \leandecl{MatrixMultiplication.\allowbreak{}Tensor.\allowbreak{}two\_\allowbreak{}thirds\_\allowbreak{}le\_\allowbreak{}of\_\allowbreak{}unit\_\allowbreak{}ne\_\allowbreak{}top}, \leandecl{MatrixMultiplication.\allowbreak{}Tensor.\allowbreak{}fractionalValue\_\allowbreak{}unit\_\allowbreak{}eq\_\allowbreak{}top\_\allowbreak{}of\_\allowbreak{}lt\_\allowbreak{}two\_\allowbreak{}thirds}, \leandecl{MatrixMultiplication.\allowbreak{}Tensor.\allowbreak{}fractionalValue\_\allowbreak{}ne\_\allowbreak{}top\_\allowbreak{}iff}, \leandecl{MatrixMultiplication.\allowbreak{}two\_\allowbreak{}le\_\allowbreak{}of\_\allowbreak{}nat\_\allowbreak{}sq\_\allowbreak{}le\_\allowbreak{}double\_\allowbreak{}rpow}.
\textit{Formalization note:} The full manuscript lemma is checked, with its positive-parameter hypothesis explicit. A stronger exact finiteness criterion holds for the library definition at every real parameter. The proof uses finite-unit normalization and a direct q² diagonal degeneration from matrix multiplication of size2q; it uses no sharper bound; the cited ceil(3n²/4) estimate is now also checked separately.

\begin{proof}
\proves{lem:value-outside-range}
\leanok
\uses{def:fractional-value,aux:finite-unit-value,aux:quadratic-matrix-subrank,aux:nonzero-unit-restriction,aux:value-finiteness}
\textit{Checked proof.} If a nonzero tensor has finite value, monotonicity makes the unit value finite. The finite-unit lemma then forces $3\tau\le\omega$. The explicit degeneration $\langle2q,2q,2q\rangle\trianglerighteq\langle q^2\rangle$ and normalization give $q^2\le(2q)^{3\tau}$ for every positive integer $q$, hence $2\le3\tau$. Therefore all parameters outside the interval have infinite value. Combine this with the independently checked finiteness within the interval for an exact criterion.
\end{proof}

\renewcommand{\thetheorem}{4.6}
\begin{lemma}[Fractional value calculus]
\label{lem:value-properties}
\manuscriptproof
\lean{MatrixMultiplication.Tensor.fractionalValue_mono_rational,MatrixMultiplication.Tensor.fractionalValue_mono_of_symmetrize_degenerates,MatrixMultiplication.Tensor.fractionalValue_permuteS3,MatrixMultiplication.Tensor.fractionalValue_pow_of_pos,MatrixMultiplication.Tensor.fractionalValue_kronecker_ge,MatrixMultiplication.Tensor.fractionalValue_diagonal_kronecker,MatrixMultiplication.Tensor.fractionalValue_matrixMul_kronecker}
\leanok
For every $\tau>0$, value in $[0,+\infty]$ is monotone under degeneration and invariant under factor permutation. For $n\ge1$, $V_\tau(T^{\otimes n})=V_\tau(T)^n$. Moreover $V_\tau(T\otimes S)\ge V_\tau(T)V_\tau(S)$, $V_\tau(\langle r\rangle\otimes T)=rV_\tau(T)$, and $V_\tau(\langle\ell,m,n\rangle\otimes T)=(\ell mn)^\tau V_\tau(T)$ for positive matrix dimensions, with $0\cdot\infty=0$.
\end{lemma}

\noindent\textit{Paper lemma 4.6. Source:} \texttt{material/main.tex:724}.
\textit{Status:} \texttt{formalized}. Main proof path.
\textit{Proof supplied in manuscript.}
\textit{Checked Lean declarations:} \leandecl{MatrixMultiplication.\allowbreak{}Tensor.\allowbreak{}fractionalValue\_\allowbreak{}mono\_\allowbreak{}rational}, \leandecl{MatrixMultiplication.\allowbreak{}Tensor.\allowbreak{}fractionalValue\_\allowbreak{}mono\_\allowbreak{}of\_\allowbreak{}symmetrize\_\allowbreak{}degenerates}, \leandecl{MatrixMultiplication.\allowbreak{}Tensor.\allowbreak{}fractionalValue\_\allowbreak{}permuteS3}, \leandecl{MatrixMultiplication.\allowbreak{}Tensor.\allowbreak{}fractionalValue\_\allowbreak{}pow\_\allowbreak{}of\_\allowbreak{}pos}, \leandecl{MatrixMultiplication.\allowbreak{}Tensor.\allowbreak{}fractionalValue\_\allowbreak{}kronecker\_\allowbreak{}ge}, \leandecl{MatrixMultiplication.\allowbreak{}Tensor.\allowbreak{}fractionalValue\_\allowbreak{}diagonal\_\allowbreak{}kronecker}, \leandecl{MatrixMultiplication.\allowbreak{}Tensor.\allowbreak{}fractionalValue\_\allowbreak{}matrixMul\_\allowbreak{}kronecker}.
\textit{Formalization note:} All six properties are checked. Permutation invariance follows directly from the defining witness set and equivalence of the symmetrizations, as in the updated manuscript. Positive-power, product, and scaling laws hold for all real parameters; the zero-power case uses unit normalization in the admissible interval. Matrix dimensions are positive as in Definition 2.4.

\begin{proof}
\proves{lem:value-properties}
\leanok
\uses{def:fractional-value,aux:symmetrized-value-monotonicity,aux:value-products-powers,aux:value-block-scaling}
\textit{Checked proof.} Transport witnesses under degeneration of the symmetrizations. Cross-power two witnesses to prove product supermultiplicativity, and transfer witnesses both into and out of tensor powers to prove equality. Absorb a diagonal or matrix block into denominator sizes for one inequality and numerator sizes for the reverse. Use extended-real supremum identities directly; restrict the power law to positive powers.
\end{proof}

\renewcommand{\thetheorem}{4.7}
\begin{lemma}[Normalization at the unit tensor]
\label{lem:value-unit}
\manuscriptproof
\lean{MatrixMultiplication.Tensor.fractionalValue_unit}
\leanok
For $2/3\le\tau\le\omega/3$, $V_\tau(\langle1\rangle)=1$.
\end{lemma}

\noindent\textit{Paper lemma 4.7. Source:} \texttt{material/main.tex:791}.
\textit{Status:} \texttt{formalized}. Main proof path.
\textit{Proof supplied in manuscript.}
\textit{Checked Lean declarations:} \leandecl{MatrixMultiplication.\allowbreak{}Tensor.\allowbreak{}fractionalValue\_\allowbreak{}unit}.
\textit{Formalization note:} Checked in aux:value-finiteness, with exact equality of the polynomial and rational-function definitions supplied by def:fractional-value.

\begin{proof}
\proves{lem:value-unit}
\leanok
\uses{aux:value-finiteness,def:fractional-value}
\textit{Checked proof.} The identity degeneration gives the lower bound. Apply the witness bound to the unit, whose asymptotic rank is one. Endpoint monotonicity bounds each witness score directly.
\end{proof}

\renewcommand{\thetheorem}{4.8}
\begin{corollary}[Diagonal and matrix multiplication normalization]
\label{cor:value-diag-matrix}
\manuscriptproof
\lean{MatrixMultiplication.Tensor.fractionalValue_diagonal,MatrixMultiplication.Tensor.fractionalValue_matrixMul}
\leanok
In the admissible interval, $V_\tau(\langle r\rangle)=r$ and $V_\tau(\langle\ell,m,n\rangle)=(\ell mn)^\tau$.
\end{corollary}

\noindent\textit{Paper corollary 4.8. Source:} \texttt{material/main.tex:809}.
\textit{Status:} \texttt{formalized}. Main proof path.
\textit{Proof supplied in manuscript.}
\textit{Checked Lean declarations:} \leandecl{MatrixMultiplication.\allowbreak{}Tensor.\allowbreak{}fractionalValue\_\allowbreak{}diagonal}, \leandecl{MatrixMultiplication.\allowbreak{}Tensor.\allowbreak{}fractionalValue\_\allowbreak{}matrixMul}.
\textit{Formalization note:} Checked over every field in the admissible parameter interval. Matrix dimensions are positive; diagonal normalization also includes size zero.

\begin{proof}
\proves{cor:value-diag-matrix}
\leanok
\uses{lem:value-unit,lem:value-properties}
\textit{Checked proof.} Apply the two exact scaling laws with the unit tensor.
\end{proof}

\renewcommand{\thetheorem}{4.9}
\begin{corollary}[Value between asymptotic subrank and rank]
\label{lem:rank-value}
\manuscriptproof
\lean{MatrixMultiplication.Tensor.asymptoticSubrank_le_fractionalValue,MatrixMultiplication.Tensor.fractionalValue_le_asymptoticRank,MatrixMultiplication.Tensor.fractionalValue_rank_sandwich,MatrixMultiplication.Tensor.fractionalValue_toReal_rank_sandwich}
\leanok
\uses{def:rank-invariants}
For $2/3\le\tau\le\omega/3$, $\widetilde Q(T)\le V_\tau(T)\le\widetilde R(T)$.
\end{corollary}

\noindent\textit{Paper corollary 4.9. Source:} \texttt{material/main.tex:816}.
\textit{Status:} \texttt{formalized}. Additional manuscript coverage.
\textit{Proof supplied in manuscript.}
\textit{Checked Lean declarations:} \leandecl{MatrixMultiplication.\allowbreak{}Tensor.\allowbreak{}asymptoticSubrank\_\allowbreak{}le\_\allowbreak{}fractionalValue}, \leandecl{MatrixMultiplication.\allowbreak{}Tensor.\allowbreak{}fractionalValue\_\allowbreak{}le\_\allowbreak{}asymptoticRank}, \leandecl{MatrixMultiplication.\allowbreak{}Tensor.\allowbreak{}fractionalValue\_\allowbreak{}rank\_\allowbreak{}sandwich}, \leandecl{MatrixMultiplication.\allowbreak{}Tensor.\allowbreak{}fractionalValue\_\allowbreak{}toReal\_\allowbreak{}rank\_\allowbreak{}sandwich}.
\textit{Formalization note:} Both inequalities are checked over every field, including the zero tensor. The value is proved finite before passing to real numbers. The lower bound takes the supremum of the individual positive-power root bounds directly, without invoking convergence.

\begin{proof}
\proves{lem:rank-value}
\leanok
\uses{lem:value-properties,cor:value-diag-matrix,aux:asymptotic-rank-calculus,aux:value-finiteness}
\textit{Checked proof.} Use a maximal diagonal restriction, monotonicity, and diagonal normalization to get $Q(T^{\otimes N})\le V_\tau(T)^N$ for each $N\ge1$. Take roots and then the supremum from the definition of asymptotic subrank. Combine this with the independently checked upper bound from witness finiteness. No subrank limit theorem is needed.
\end{proof}

\renewcommand{\thetheorem}{4.10}
\begin{lemma}[Entropy bound for partially direct sums]
\label{lem:diamondsum-value-prob}
\manuscriptproof
\lean{MatrixMultiplication.Tensor.fractionalValue_partialSum_entropy}
\leanok
For an admissible $\tau$, $T=\mathop{\diamondplus}_{h=1}^mT_h$, and a probability distribution $q$ on the blocks, $V_\tau(T)\ge2^{\tau H(q)}\prod_hV_\tau(T_h)^{q_h}$.
\end{lemma}

\noindent\textit{Paper lemma 4.10. Source:} \texttt{material/main.tex:837}.
\textit{Status:} \texttt{formalized}. Additional manuscript coverage.
\textit{Proof supplied in manuscript.}
\textit{Checked Lean declarations:} \leandecl{MatrixMultiplication.\allowbreak{}Tensor.\allowbreak{}fractionalValue\_\allowbreak{}partialSum\_\allowbreak{}entropy}.
\textit{Formalization note:} Checked for every real probability vector, including zero probabilities and zero block values. The proof uses the already checked superadditivity inequality and weighted AM-GM; no entropy limit or rational approximation is assumed.

\begin{proof}
\proves{lem:diamondsum-value-prob}
\leanok
\uses{lem:value-diamondsum,aux:entropy-variational-upper-bound,aux:value-finiteness}
\textit{Checked proof.} Apply the weighted AM--GM entropy bound to the finite block values, obtaining an upper bound by $(\sum_h V_\tau(T_h)^{1/\tau})^\tau$. Superadditivity and positivity of $\tau$ bound this by $V_\tau(T)$.
\end{proof}

\renewcommand{\thetheorem}{4.11}
\begin{remark}[Comparison with the laser-method estimate]
\label{rem:laser-comparison}
\lean{MatrixMultiplication.Tensor.strassenLeft,MatrixMultiplication.Tensor.strassenRight,MatrixMultiplication.Tensor.strassenTensor,MatrixMultiplication.Tensor.diagonal_degenerates_strassenTensor,MatrixMultiplication.Tensor.flatteningSecond_rank_strassenTensor,MatrixMultiplication.Tensor.borderRank_strassenTensor,MatrixMultiplication.Tensor.asymptoticRank_strassenTensor,MatrixMultiplication.Tensor.fractionalValue_strassenLeft,MatrixMultiplication.Tensor.fractionalValue_strassenRight,MatrixMultiplication.Tensor.rpow_le_fractionalValue_strassenTensor,MatrixMultiplication.Tensor.strassen_value_inequality,MatrixMultiplication.strassen_exponent_inequality,MatrixMultiplication.omega_le_strassen_bound,MatrixMultiplication.strassen_bound_lt_decimal,MatrixMultiplication.omega_lt_strassen_decimal}
\leanok
The entropy coefficient $\tau\ge2/3$ improves the coefficient $2/3$ from the cited laser-method estimate. For $n\ge1$, Strassen's tensor $S_n=\langle n,1,1\rangle\diamondplus\langle1,n,1\rangle$ has border rank and asymptotic rank $n+1$. For every admissible $\tau$, $(2n)^\tau\le V_\tau(S_n)\le n+1$. Taking $\tau=\omega/3$ and $n=3$ yields $\omega\le3\log4/\log6<2.32112$.
\end{remark}

\noindent\textit{Paper remark 4.11. Source:} \texttt{material/main.tex:859}.
\textit{Status:} \texttt{formalized}. Additional manuscript coverage.
\textit{Checked Lean declarations:} \leandecl{MatrixMultiplication.\allowbreak{}Tensor.\allowbreak{}strassenLeft}, \leandecl{MatrixMultiplication.\allowbreak{}Tensor.\allowbreak{}strassenRight}, \leandecl{MatrixMultiplication.\allowbreak{}Tensor.\allowbreak{}strassenTensor}, \leandecl{MatrixMultiplication.\allowbreak{}Tensor.\allowbreak{}diagonal\_\allowbreak{}degenerates\_\allowbreak{}strassenTensor}, \leandecl{MatrixMultiplication.\allowbreak{}Tensor.\allowbreak{}flatteningSecond\_\allowbreak{}rank\_\allowbreak{}strassenTensor}, \leandecl{MatrixMultiplication.\allowbreak{}Tensor.\allowbreak{}borderRank\_\allowbreak{}strassenTensor}, \leandecl{MatrixMultiplication.\allowbreak{}Tensor.\allowbreak{}asymptoticRank\_\allowbreak{}strassenTensor}, \leandecl{MatrixMultiplication.\allowbreak{}Tensor.\allowbreak{}fractionalValue\_\allowbreak{}strassenLeft}, \leandecl{MatrixMultiplication.\allowbreak{}Tensor.\allowbreak{}fractionalValue\_\allowbreak{}strassenRight}, \leandecl{MatrixMultiplication.\allowbreak{}Tensor.\allowbreak{}rpow\_\allowbreak{}le\_\allowbreak{}fractionalValue\_\allowbreak{}strassenTensor}, \leandecl{MatrixMultiplication.\allowbreak{}Tensor.\allowbreak{}strassen\_\allowbreak{}value\_\allowbreak{}inequality}, \leandecl{MatrixMultiplication.\allowbreak{}strassen\_\allowbreak{}exponent\_\allowbreak{}inequality}, \leandecl{MatrixMultiplication.\allowbreak{}omega\_\allowbreak{}le\_\allowbreak{}strassen\_\allowbreak{}bound}, \leandecl{MatrixMultiplication.\allowbreak{}strassen\_\allowbreak{}bound\_\allowbreak{}lt\_\allowbreak{}decimal}, \leandecl{MatrixMultiplication.\allowbreak{}omega\_\allowbreak{}lt\_\allowbreak{}strassen\_\allowbreak{}decimal}.
\textit{Formalization note:} The Strassen tensor construction, exact border and asymptotic ranks, value inequality, and strict decimal upper bound are checked over every field. The proof uses binary partial-sum superadditivity, which gives the same bound as the equal-weight entropy inequality here. It does not use the main nine-quarters theorem. The historical attribution and comparison with other methods remain expository.

\begin{proof}
\proves{rem:laser-comparison}
\leanok
\uses{def:matrix-multiplication,lem:value-diamondsum,lem:rank-value,cor:value-diag-matrix,lem:flattening-lb}
\textit{Checked proof.} The coefficient of $\varepsilon$ in $\sum_{i=1}^n x_i(y_0+\varepsilon y_i)(z_0+\varepsilon z_i)-(\sum_{i=1}^n x_i)y_0z_0$ is $S_n$. This gives a degeneration from $\langle n+1\rangle$ over every field. An identity submatrix of the second flattening gives the matching lower bound. Each scalar-multiplication block has value $n^\tau$; binary superadditivity gives $(2n)^\tau$. Specialize to $n=3$, take logarithms, and certify the decimal comparison by exact rational bounds for $\log4$ and $\log6$.
\end{proof}

\renewcommand{\thetheorem}{4.12}
\begin{lemma}[Superadditivity for partially direct sums]
\label{lem:value-diamondsum}
\manuscriptproof
\lean{MatrixMultiplication.Tensor.fractionalValue_partialSum_superadditive,MatrixMultiplication.Tensor.fractionalValue_partialSum_toReal_superadditive,MatrixMultiplication.Tensor.fractionalValue_partialSum_toReal_add_le}
\leanok
For $2/3\le\tau\le\omega/3$ and $T=\mathop{\diamondplus}_hT_h$, $V_\tau(T)^{1/\tau}\ge\sum_hV_\tau(T_h)^{1/\tau}$.
\end{lemma}

\noindent\textit{Paper lemma 4.12. Source:} \texttt{material/main.tex:868}.
\textit{Status:} \texttt{formalized}. Main proof path.
\textit{Proof supplied in manuscript.}
\textit{Checked Lean declarations:} \leandecl{MatrixMultiplication.\allowbreak{}Tensor.\allowbreak{}fractionalValue\_\allowbreak{}partialSum\_\allowbreak{}superadditive}, \leandecl{MatrixMultiplication.\allowbreak{}Tensor.\allowbreak{}fractionalValue\_\allowbreak{}partialSum\_\allowbreak{}toReal\_\allowbreak{}superadditive}, \leandecl{MatrixMultiplication.\allowbreak{}Tensor.\allowbreak{}fractionalValue\_\allowbreak{}partialSum\_\allowbreak{}toReal\_\allowbreak{}add\_\allowbreak{}le}.
\textit{Formalization note:} Checked over every field, including empty families and zero-valued blocks. The Lean proof reverses the manuscript order: superadditivity is proved first by summing finite type bounds, and the entropy bound follows afterward. The alternative proof is documented here rather than in a manuscript annotation.

\begin{proof}
\proves{lem:value-diamondsum}
\leanok
\uses{lem:decoupling-asymptotic,lem:value-properties,aux:value-finiteness,aux:finite-tensor-products,aux:multinomial-sum-bound}
\textit{Checked proof.} Write $a_h=V_\tau(T_h)^{1/\tau}$ and $B=V_\tau(T)^{1/\tau}$. Type separation and exact value scaling give $Q^\tau\prod_hV_\tau(T_h)^{c_h}\le2V_\tau(T)^N$. Raise to $1/\tau$ and apply the checked multinomial-sum bound with $C=2^{1/\tau}$. Admissible finiteness justifies real arithmetic; transfer the result back to extended nonnegative reals.
\end{proof}

\renewcommand{\thetheorem}{5.1}
\begin{lemma}[Degeneration from compatible bilinear maps]
\label{lem:degeneration-sequences}
\manuscriptproof
\lean{MatrixMultiplication.Tensor.degenerates_partialSum_of_compatible_maps,MatrixMultiplication.Tensor.structureTensor,MatrixMultiplication.Tensor.map_structureTensor,MatrixMultiplication.Tensor.degenerates_partialSum_of_exact}
\leanok
Let $\iota_V:V_0\to V$, $\pi_V:V\to V_1$, $\iota_W:W_0\to W$, and $\pi_W:W\to W_1$ be linear maps. Assume $\pi_V$ is surjective, $\iota_W$ is injective, and $\pi_W\iota_W=0$. If bilinear maps satisfy $\varphi(u,\iota_Vv_0)=\iota_W\varphi_0(u,v_0)$ and $\varphi_1(u,\pi_Vv)=\pi_W\varphi(u,v)$, their structure tensors satisfy $T\trianglerighteq T_0\diamondplus T_1$.
\end{lemma}

\noindent\textit{Paper lemma 5.1. Source:} \texttt{material/main.tex:885}.
\textit{Status:} \texttt{formalized}. Additional manuscript coverage.
\textit{Proof supplied in manuscript.}
\textit{Checked Lean declarations:} \leandecl{MatrixMultiplication.\allowbreak{}Tensor.\allowbreak{}degenerates\_\allowbreak{}partialSum\_\allowbreak{}of\_\allowbreak{}compatible\_\allowbreak{}maps}, \leandecl{MatrixMultiplication.\allowbreak{}Tensor.\allowbreak{}structureTensor}, \leandecl{MatrixMultiplication.\allowbreak{}Tensor.\allowbreak{}map\_\allowbreak{}structureTensor}, \leandecl{MatrixMultiplication.\allowbreak{}Tensor.\allowbreak{}degenerates\_\allowbreak{}partialSum\_\allowbreak{}of\_\allowbreak{}exact}.
\textit{Formalization note:} The revised weaker hypotheses exactly match the compatible-map theorem. Arbitrary finite bases are allowed, with the shared first-factor basis preserved. The short-exact-sequence theorem is retained as a special case.

\begin{proof}
\proves{lem:degeneration-sequences}
\leanok
\uses{aux:degeneration-calculus,def:partial-direct-sum,aux:split-coefficient-degeneration}
\textit{Checked proof.} Choose a linear left inverse of the output inclusion and a linear right inverse of the input quotient. Compatibility gives the two prescribed diagonal blocks; the zero composite on the output side makes the remaining required block zero. The structure-tensor action lemma translates input precomposition and output postcomposition into coefficient maps, and the checked split-map criterion gives the degeneration.
\end{proof}

\renewcommand{\thetheorem}{5.2}
\begin{remark}[The induced quotient bilinear map]
\label{rem:quotient-degeneration}
\lean{MatrixMultiplication.Tensor.quotientBilinearMap,MatrixMultiplication.Tensor.quotientBilinearMap_mkQ,MatrixMultiplication.Tensor.quotientBilinearMap_unique,MatrixMultiplication.Tensor.existsUnique_quotientBilinearMap,MatrixMultiplication.Tensor.degenerates_partialSum_quotient}
\leanok
Under the compatible-map hypotheses, $\bar\varphi:U\times(V/\operatorname{im}\iota_V)\to W/\operatorname{im}\iota_W$, $\bar\varphi(u,[v])=[\varphi(u,v)]$, exists uniquely. If $\iota_W$ is injective, $T_\varphi\trianglerighteq T_{\varphi_0}\diamondplus T_{\bar\varphi}$.
\end{remark}

\noindent\textit{Paper remark 5.2. Source:} \texttt{material/main.tex:914}.
\textit{Status:} \texttt{formalized}. Additional manuscript coverage.
\textit{Checked Lean declarations:} \leandecl{MatrixMultiplication.\allowbreak{}Tensor.\allowbreak{}quotientBilinearMap}, \leandecl{MatrixMultiplication.\allowbreak{}Tensor.\allowbreak{}quotientBilinearMap\_\allowbreak{}mkQ}, \leandecl{MatrixMultiplication.\allowbreak{}Tensor.\allowbreak{}quotientBilinearMap\_\allowbreak{}unique}, \leandecl{MatrixMultiplication.\allowbreak{}Tensor.\allowbreak{}existsUnique\_\allowbreak{}quotientBilinearMap}, \leandecl{MatrixMultiplication.\allowbreak{}Tensor.\allowbreak{}degenerates\_\allowbreak{}partialSum\_\allowbreak{}quotient}.
\textit{Formalization note:} The manuscript states the quotient construction and its consequence without a separate proof. The formalization supplies well-definedness, the representative formula, uniqueness, and the degeneration.

\begin{proof}
\proves{rem:quotient-degeneration}
\leanok
\uses{lem:degeneration-sequences}
\textit{Checked proof.} Compatibility sends the input image into the output image, so quotient lifting constructs a bilinear map. Surjectivity of the quotient projection gives uniqueness. Apply the compatible-map degeneration theorem to the quotient projections.
\end{proof}

\renewcommand{\thetheorem}{5.3}
\begin{lemma}[Polynomial multiplication concavity degeneration]
\label{lem:pm-concavity-degeneration}
\manuscriptproof
\lean{MatrixMultiplication.Tensor.concavity_edge_tail,MatrixMultiplication.Tensor.concavity_edge_diagonal,MatrixMultiplication.Tensor.concavity_section_diagonal,MatrixMultiplication.Tensor.polynomialMul_concavity_degeneration}
\leanok
For $a,p\ge1$ and $b\ge2$, $\mathrm{PM}_{a,b}\otimes\langle1,1,p+1\rangle\trianglerighteq(\mathrm{PM}_{a,b-1}\otimes\langle1,1,p\rangle)\diamondplus\mathrm{PM}_{a,b+p}$.
\end{lemma}

\noindent\textit{Paper lemma 5.3. Source:} \texttt{material/main.tex:918}.
\textit{Status:} \texttt{formalized}. Main proof path.
\textit{Proof supplied in manuscript.}
\textit{Checked Lean declarations:} \leandecl{MatrixMultiplication.\allowbreak{}Tensor.\allowbreak{}concavity\_\allowbreak{}edge\_\allowbreak{}tail}, \leandecl{MatrixMultiplication.\allowbreak{}Tensor.\allowbreak{}concavity\_\allowbreak{}edge\_\allowbreak{}diagonal}, \leandecl{MatrixMultiplication.\allowbreak{}Tensor.\allowbreak{}concavity\_\allowbreak{}section\_\allowbreak{}diagonal}, \leandecl{MatrixMultiplication.\allowbreak{}Tensor.\allowbreak{}polynomialMul\_\allowbreak{}concavity\_\allowbreak{}degeneration}.
\textit{Formalization note:} The checked coefficient theorem is stronger: a and p may be zero and b need only be at least one. The repeated model is canonically equivalent to the displayed matrix-product tensor, with a literal common first factor. The general exact-sequence lemma remains separate.

\begin{proof}
\proves{lem:pm-concavity-degeneration}
\leanok
\uses{aux:polynomial-coordinate-models,aux:split-coefficient-degeneration,aux:bivariate-grid-maps}
\textit{Checked proof.} Use multiplication by $x-y$ and the boundary monomial section as input maps. Use diagonal-tail extraction and evaluation at $y=x$ as output maps. Check the three split-map identities coefficientwise and apply the order-one degeneration criterion. The revised abstract lemma uses weaker compatibility hypotheses; the concrete coefficient proof remains a direct implementation of the same block identities.
\end{proof}

\renewcommand{\thetheorem}{5.4}
\begin{corollary}[Concavity of polynomial multiplication value]
\label{cor:pm-concavity}
\manuscriptproof
\lean{MatrixMultiplication.Tensor.polynomialMul_value_concavity}
\leanok
For admissible $\tau$, let $P(a,b)=V_\tau(\mathrm{PM}_{a,b})^{1/\tau}$. If $a,p\ge1$ and $b\ge2$, $(p+1)P(a,b)\ge pP(a,b-1)+P(a,b+p)$. In particular, $P(a,\cdot)$ is discretely concave.
\end{corollary}

\noindent\textit{Paper corollary 5.4. Source:} \texttt{material/main.tex:936}.
\textit{Status:} \texttt{formalized}. Main proof path.
\textit{Proof supplied in manuscript.}
\textit{Checked Lean declarations:} \leandecl{MatrixMultiplication.\allowbreak{}Tensor.\allowbreak{}polynomialMul\_\allowbreak{}value\_\allowbreak{}concavity}.
\textit{Formalization note:} Checked for all natural a,p and b at least one; zero multiplicity is handled as equality. No degeneration is left as an assumption.

\begin{proof}
\proves{cor:pm-concavity}
\leanok
\uses{lem:pm-concavity-degeneration,aux:binary-partial-value-support,aux:value-finiteness}
\textit{Checked proof.} Apply monotonicity and partial-sum superadditivity to the preceding degeneration, then use matrix multiplication scaling to pull out $p$ and $p+1$.
\end{proof}

\renewcommand{\thetheorem}{5.5}
\begin{lemma}[Three-block polynomial multiplication degeneration]
\label{lem:pm-tripling-degeneration}
\manuscriptproof
\lean{MatrixMultiplication.Tensor.polynomialMul_threeBlock_degeneration}
\leanok
For positive $a,b,c,d$, $\mathrm{PM}_{a,a+b+c+d-1}\trianglerighteq\mathrm{PM}_{a,b}\diamondplus(23)\cdot\mathrm{PM}_{a,c}\diamondplus\mathrm{PM}_{a,d}$.
\end{lemma}

\noindent\textit{Paper lemma 5.5. Source:} \texttt{material/main.tex:954}.
\textit{Status:} \texttt{formalized}. Main proof path.
\textit{Proof supplied in manuscript.}
\textit{Checked Lean declarations:} \leandecl{MatrixMultiplication.\allowbreak{}Tensor.\allowbreak{}polynomialMul\_\allowbreak{}threeBlock\_\allowbreak{}degeneration}.
\textit{Formalization note:} The coefficient degeneration is proved over every commutative semiring. It retains the identical first-factor map in all three blocks.

\begin{proof}
\proves{lem:pm-tripling-degeneration}
\leanok
\uses{def:polynomial-multiplication,aux:polynomial-coordinate-models,aux:monomial-degeneration,def:partial-direct-sum}
\textit{Checked proof.} Partition the second and third index sets into the three intervals in the manuscript. Two off-diagonal blocks receive a positive parameter weight and vanish. Reindex the surviving middle block in reverse order to identify the transposed polynomial multiplication tensor. All reindexings act identically on the shared first factor.
\end{proof}

\renewcommand{\thetheorem}{5.6}
\begin{corollary}[Tripling inequality]
\label{cor:pm-tripling}
\manuscriptproof
\lean{MatrixMultiplication.Tensor.polynomialMul_value_threeBlock,MatrixMultiplication.Tensor.polynomialMul_value_tripling}
\leanok
For admissible $\tau$ and positive $a$, $P(a,4a-1)\ge3P(a,a)$.
\end{corollary}

\noindent\textit{Paper corollary 5.6. Source:} \texttt{material/main.tex:1114}.
\textit{Status:} \texttt{formalized}. Main proof path.
\textit{Proof supplied in manuscript.}
\textit{Checked Lean declarations:} \leandecl{MatrixMultiplication.\allowbreak{}Tensor.\allowbreak{}polynomialMul\_\allowbreak{}value\_\allowbreak{}threeBlock}, \leandecl{MatrixMultiplication.\allowbreak{}Tensor.\allowbreak{}polynomialMul\_\allowbreak{}value\_\allowbreak{}tripling}.
\textit{Formalization note:} The general three-block value inequality is checked before specializing all four sizes to a.

\begin{proof}
\proves{cor:pm-tripling}
\leanok
\uses{lem:pm-tripling-degeneration,aux:binary-partial-value-support,aux:value-finiteness}
\textit{Checked proof.} Set $a=b=c=d$ in the preceding degeneration, use permutation invariance of the value, and apply partial-sum superadditivity.
\end{proof}

\renewcommand{\thetheorem}{6.1}
\begin{lemma}[Quantitative Gamma-function growth]
\label{lem:4/3}
\manuscriptproof
\lean{MatrixMultiplication.Tensor.polynomialMul_value_gamma_lower_bound,MatrixMultiplication.gammaGrowthBound,MatrixMultiplication.GrowthHypotheses.diagonal_gamma_lower_bound,MatrixMultiplication.rpow_le_gammaGrowthBound}
\leanok
For admissible $\tau$ and $a\ge1$, $P_\tau(a,a)\ge G(a)\ge a^{4/3}$, where $G(a)=\frac{3a-1}{2\Gamma(4/3)}\frac{\Gamma(a+1/3)}{\Gamma(a)}$.
\end{lemma}

\noindent\textit{Paper lemma 6.1. Source:} \texttt{material/main.tex:1126}.
\textit{Status:} \texttt{formalized}. Additional manuscript coverage.
\textit{Proof supplied in manuscript.}
\textit{Checked Lean declarations:} \leandecl{MatrixMultiplication.\allowbreak{}Tensor.\allowbreak{}polynomialMul\_\allowbreak{}value\_\allowbreak{}gamma\_\allowbreak{}lower\_\allowbreak{}bound}, \leandecl{MatrixMultiplication.\allowbreak{}gammaGrowthBound}, \leandecl{MatrixMultiplication.\allowbreak{}GrowthHypotheses.\allowbreak{}diagonal\_\allowbreak{}gamma\_\allowbreak{}lower\_\allowbreak{}bound}, \leandecl{MatrixMultiplication.\allowbreak{}rpow\_\allowbreak{}le\_\allowbreak{}gammaGrowthBound}.
\textit{Formalization note:} Both inequalities hold uniformly for every positive integer. The Gamma sharpening is independent of the elementary main-theorem route. The asymptotic coefficient is recorded separately under Remark 6.2.

\begin{proof}
\proves{lem:4/3}
\leanok
\uses{eq:PM-recurrence}
\textit{Checked proof.} The Gamma expression equals one at $a=1$ and its successive ratio is the recurrence multiplier. Induction proves the first inequality. The same multiplier dominates $((a+1)/a)^{4/3}$, proving the second inequality by induction.
\end{proof}

\renewcommand{\thetheorem}{6.1}
\begin{equationresult}[Uniform Gamma and elementary lower bounds]
\label{eq:lem4/3}
\manuscriptproof
\lean{MatrixMultiplication.Tensor.polynomialMul_value_gamma_bounds}
\leanok
For admissible $\tau$ and $a\ge1$, $P_\tau(a,a)\ge G(a)\ge a^{4/3}$.
\end{equationresult}

\noindent\textit{Paper equation 6.1. Source:} \texttt{material/main.tex:1128}.
\textit{Status:} \texttt{formalized}. Additional manuscript coverage.
\textit{Proof supplied in manuscript.}
\textit{Checked Lean declarations:} \leandecl{MatrixMultiplication.\allowbreak{}Tensor.\allowbreak{}polynomialMul\_\allowbreak{}value\_\allowbreak{}gamma\_\allowbreak{}bounds}.

\begin{proof}
\proves{eq:lem4/3}
\leanok
\uses{lem:4/3}
\textit{Checked proof.} Combine the exact Gamma lower bound with its elementary lower bound.
\end{proof}

\renewcommand{\thetheorem}{6.2}
\begin{equationresult}[Diagonal recurrence]
\label{eq:PM-recurrence}
\manuscriptproof
\lean{MatrixMultiplication.Tensor.polynomialMul_value_diagonal_recurrence}
\leanok
For the manuscript function $P(a,b)=V_\tau(\mathrm{PM}_{a,b})^{1/\tau}$ with admissible $\tau$, $P(a+1,a+1)\ge\frac{(3a+1)(3a+2)}{3a(3a-1)}P(a,a)$ for $a\ge1$.
\end{equationresult}

\noindent\textit{Paper equation 6.2. Source:} \texttt{material/main.tex:1153}.
\textit{Status:} \texttt{formalized}. Additional manuscript coverage.
\textit{Proof supplied in manuscript.}
\textit{Checked Lean declarations:} \leandecl{MatrixMultiplication.\allowbreak{}Tensor.\allowbreak{}polynomialMul\_\allowbreak{}value\_\allowbreak{}diagonal\_\allowbreak{}recurrence}.
\textit{Formalization note:} The generic numerical recurrence is now instantiated with the checked polynomial-value hypotheses.

\begin{proof}
\proves{eq:PM-recurrence}
\leanok
\uses{aux:polynomial-value-hypotheses,aux:numerical-recurrence}
\textit{Checked proof.} Apply concavity first with $(b,p)=(a+1,3a-2)$ and then with $(b,p)=(a+1,3a+2)$. Use symmetry and tripling, collect the repeated $P(a+1,a+1)$ term, and divide only by strictly positive denominators.
\end{proof}

\renewcommand{\thetheorem}{6.3}
\begin{equationresult}[Finite-product expression for the Gamma lower bound]
\label{eq:PM-Gamma}
\manuscriptproof
\lean{MatrixMultiplication.gammaGrowthBound_eq_prod,MatrixMultiplication.gammaGrowthBound_prod_eq_gamma,MatrixMultiplication.Tensor.polynomialMul_value_prod_lower_bound}
\leanok
$P_\tau(a,a)\ge\frac{3a-1}{2}\prod_{i=1}^{a-1}\frac{3i+1}{3i}=\frac{3a-1}{2}\prod_{i=1}^{a-1}\frac{i+1/3}{i}=\frac{3a-1}{2}\frac{\Gamma(a+1/3)}{\Gamma(4/3)\Gamma(a)}=G(a)$ for $a\ge1$. The parameter satisfies $2/3\le\tau\le\omega/3$.
\end{equationresult}

\noindent\textit{Paper equation 6.3. Source:} \texttt{material/main.tex:1161}.
\textit{Status:} \texttt{formalized}. Additional manuscript coverage.
\textit{Proof supplied in manuscript.}
\textit{Checked Lean declarations:} \leandecl{MatrixMultiplication.\allowbreak{}gammaGrowthBound\_\allowbreak{}eq\_\allowbreak{}prod}, \leandecl{MatrixMultiplication.\allowbreak{}gammaGrowthBound\_\allowbreak{}prod\_\allowbreak{}eq\_\allowbreak{}gamma}, \leandecl{MatrixMultiplication.\allowbreak{}Tensor.\allowbreak{}polynomialMul\_\allowbreak{}value\_\allowbreak{}prod\_\allowbreak{}lower\_\allowbreak{}bound}.

\begin{proof}
\proves{eq:PM-Gamma}
\leanok
\uses{lem:4/3}
\textit{Checked proof.} Iterate the Gamma functional equation and simplify the products termwise. Apply the checked exact lower bound.
\end{proof}

\renewcommand{\thetheorem}{6.2}
\begin{remark}[Gamma asymptotics and numerical coefficient]
\label{rem:gamma-coefficient}
\lean{MatrixMultiplication.gamma_ratio_one_third_asymptotic,MatrixMultiplication.gammaGrowthBound_div_rpow_tendsto,MatrixMultiplication.gamma_growth_coefficient_bounds}
\leanok
The Gamma expression satisfies $G(a)=\frac{3}{2\Gamma(4/3)}a^{4/3}(1+o(1))$, hence the same asymptotic lower bound holds for $P_\tau(a,a)$. Its coefficient lies strictly between $1.67976$ and $1.67977$.
\end{remark}

\noindent\textit{Paper remark 6.2. Source:} \texttt{material/main.tex:1175}.
\textit{Status:} \texttt{formalized}. Additional manuscript coverage.
\textit{Checked Lean declarations:} \leandecl{MatrixMultiplication.\allowbreak{}gamma\_\allowbreak{}ratio\_\allowbreak{}one\_\allowbreak{}third\_\allowbreak{}asymptotic}, \leandecl{MatrixMultiplication.\allowbreak{}gammaGrowthBound\_\allowbreak{}div\_\allowbreak{}rpow\_\allowbreak{}tendsto}, \leandecl{MatrixMultiplication.\allowbreak{}gamma\_\allowbreak{}growth\_\allowbreak{}coefficient\_\allowbreak{}bounds}.
\textit{Formalization note:} The manuscript derives the asymptotic using Stirling but states the decimal enclosure without a derivation. Accordingly this combined node has a thin border; the separate asymptotic node records the supplied derivation. Lean uses Euler’s Gamma limit and exact rational arithmetic.

\begin{proof}
\proves{rem:gamma-coefficient}
\leanok
\uses{lem:4/3,aux:gamma-asymptotic,aux:gamma-coefficient-enclosure}
\textit{Checked proof.} Combine the exact lower bound with the Gamma-ratio limit and the independently certified rational enclosure.
\end{proof}

\renewcommand{\thetheorem}{6.3}
\begin{theorem}[Main result: Matrix multiplication exponent over every field]
\label{thm:omega-le-nine-fourths}
\manuscriptproof
\mainresult{ω ≤ 9/4}
\lean{MatrixMultiplication.omega_le_nine_quarters}
\leanok
For every field $F$, $\omega_F\le\frac94$. The only algebraic hypothesis is that $F$ is a field.
\end{theorem}

\noindent\textit{Paper theorem 6.3. Source:} \texttt{material/main.tex:1182}.
\textit{Status:} \texttt{formalized}. Main proof path.
\textit{Proof supplied in manuscript.}
\textit{Checked Lean declarations:} \leandecl{MatrixMultiplication.\allowbreak{}omega\_\allowbreak{}le\_\allowbreak{}nine\_\allowbreak{}quarters}.
\textit{Formalization note:} Checked as a closed theorem for an arbitrary type K with only a Field K instance. There are no characteristic, cardinality, value-functional, or unproved growth assumptions.

\begin{proof}
\proves{thm:omega-le-nine-fourths}
\leanok
\uses{aux:omega-bounds,aux:polynomial-value-hypotheses,aux:polynomial-value-elementary-growth,aux:polynomial-value-upper-bound,aux:numerical-endpoint}
\textit{Checked proof.} Take $\tau=\omega_F/3$, which is admissible because $\omega_F\ge2$. The concrete polynomial value satisfies all numerical growth hypotheses. Its elementary growth bound and the uniform upper bound give $a^{4\omega_F/9}\le V_{\omega_F/3}(\mathrm{PM}_{a,a})\le2a-1$ for every $a\ge1$. The checked numerical endpoint theorem yields $\omega_F\le9/4$. The exact rank formula and the Gamma-function sharpening are not required.
\end{proof}

\section{Introduction cross-references}

\renewcommand{\thetheorem}{1.1}
\begin{equationresult}[Introduction: matrix multiplication scaling]
\label{eq:matmul-normalization}
\manuscriptproof
\lean{MatrixMultiplication.Tensor.fractionalValue_kronecker_matrixMul}
\leanok
For $2\le3\tau\le\omega$ and positive $\ell,m,n$, $V_\tau(T\otimes\langle\ell,m,n\rangle)=(\ell mn)^\tau V_\tau(T)$.
\end{equationresult}

\noindent\textit{Paper equation 1.1. Source:} \texttt{material/main.tex:193}.
\textit{Status:} \texttt{formalized}. Additional manuscript coverage.
\textit{Proof supplied in manuscript.}
\textit{Checked Lean declarations:} \leandecl{MatrixMultiplication.\allowbreak{}Tensor.\allowbreak{}fractionalValue\_\allowbreak{}kronecker\_\allowbreak{}matrixMul}.
\textit{Formalization note:} The library theorem holds for every real parameter, including infinite values.

\begin{proof}
\proves{eq:matmul-normalization}
\leanok
\uses{lem:value-properties}
\textit{Checked proof.} Apply matrix-block scaling and Kronecker commutativity to place the matrix factor on the right.
\end{proof}

\renewcommand{\thetheorem}{1.2}
\begin{equationresult}[Introduction: value upper bound]
\label{eq:schonhages-ineq}
\manuscriptproof
\lean{MatrixMultiplication.Tensor.fractionalValue_le_asymptoticRank}
\leanok
For $2\le3\tau\le\omega$, $V_\tau(T)\le\widetilde R(T)$.
\end{equationresult}

\noindent\textit{Paper equation 1.2. Source:} \texttt{material/main.tex:196}.
\textit{Status:} \texttt{formalized}. Additional manuscript coverage.
\textit{Proof supplied in manuscript.}
\textit{Checked Lean declarations:} \leandecl{MatrixMultiplication.\allowbreak{}Tensor.\allowbreak{}fractionalValue\_\allowbreak{}le\_\allowbreak{}asymptoticRank}.
\textit{Formalization note:} This introduction restatement is covered by the checked rank-value corollary.

\begin{proof}
\proves{eq:schonhages-ineq}
\leanok
\uses{lem:rank-value}
\textit{Checked proof.} This is the introduction's restatement of the rank-value corollary.
\end{proof}

\renewcommand{\thetheorem}{1.3}
\begin{equationresult}[Introduction: partial-sum inequality]
\label{eq:partial-direct-sum-ineq}
\manuscriptproof
\lean{MatrixMultiplication.Tensor.fractionalValue_partialSum_superadditive}
\leanok
For admissible $\tau$, $V_\tau(\mathop{\diamondplus}_hT_h)^{1/\tau}\ge\sum_hV_\tau(T_h)^{1/\tau}$.
\end{equationresult}

\noindent\textit{Paper equation 1.3. Source:} \texttt{material/main.tex:209}.
\textit{Status:} \texttt{formalized}. Additional manuscript coverage.
\textit{Proof supplied in manuscript.}
\textit{Checked Lean declarations:} \leandecl{MatrixMultiplication.\allowbreak{}Tensor.\allowbreak{}fractionalValue\_\allowbreak{}partialSum\_\allowbreak{}superadditive}.
\textit{Formalization note:} Checked by the same unconditional finite-family theorem as Lemma 4.12.

\begin{proof}
\proves{eq:partial-direct-sum-ineq}
\leanok
\uses{lem:value-diamondsum}
\textit{Checked proof.} This is the introduction's restatement of partial-sum superadditivity.
\end{proof}

\renewcommand{\thetheorem}{1.4}
\begin{equationresult}[Introduction: asymptotic rank of polynomial multiplication]
\label{eq:a-rk-pm}
\manuscriptproof
\lean{MatrixMultiplication.Tensor.asymptoticRank_polynomialMul}
\leanok
For positive $a,b$, $\widetilde R(\mathrm{PM}_{a,b})=a+b-1$.
\end{equationresult}

\noindent\textit{Paper equation 1.4. Source:} \texttt{material/main.tex:237}.
\textit{Status:} \texttt{formalized}. Additional manuscript coverage.
\textit{Proof supplied in manuscript.}
\textit{Checked Lean declarations:} \leandecl{MatrixMultiplication.\allowbreak{}Tensor.\allowbreak{}asymptoticRank\_\allowbreak{}polynomialMul}.

\begin{proof}
\proves{eq:a-rk-pm}
\leanok
\uses{lem:poly-mult-from-sum}
\textit{Checked proof.} Use the exact polynomial-multiplication asymptotic-rank theorem.
\end{proof}
